% Options for packages loaded elsewhere
% Options for packages loaded elsewhere
\PassOptionsToPackage{unicode}{hyperref}
\PassOptionsToPackage{hyphens}{url}
%
\documentclass[
  ignorenonframetext,
  aspectratio=169,
  spanish,
]{beamer}
\newif\ifbibliography
\usepackage{pgfpages}
\setbeamertemplate{caption}[numbered]
\setbeamertemplate{caption label separator}{: }
\setbeamercolor{caption name}{fg=normal text.fg}
\beamertemplatenavigationsymbolsempty
% remove section numbering
\setbeamertemplate{part page}{
  \centering
  \begin{beamercolorbox}[sep=16pt,center]{part title}
    \usebeamerfont{part title}\insertpart\par
  \end{beamercolorbox}
}
\setbeamertemplate{section page}{
  \centering
  \begin{beamercolorbox}[sep=12pt,center]{section title}
    \usebeamerfont{section title}\insertsection\par
  \end{beamercolorbox}
}
\setbeamertemplate{subsection page}{
  \centering
  \begin{beamercolorbox}[sep=8pt,center]{subsection title}
    \usebeamerfont{subsection title}\insertsubsection\par
  \end{beamercolorbox}
}
% Prevent slide breaks in the middle of a paragraph
\widowpenalties 1 10000
\raggedbottom
\AtBeginPart{
  \frame{\partpage}
}
\AtBeginSection{
  \ifbibliography
  \else
    \frame{\sectionpage}
  \fi
}
\AtBeginSubsection{
  \frame{\subsectionpage}
}
\usepackage{iftex}
\ifPDFTeX
  \usepackage[T1]{fontenc}
  \usepackage[utf8]{inputenc}
  \usepackage{textcomp} % provide euro and other symbols
\else % if luatex or xetex
  \usepackage{unicode-math} % this also loads fontspec
  \defaultfontfeatures{Scale=MatchLowercase}
  \defaultfontfeatures[\rmfamily]{Ligatures=TeX,Scale=1}
\fi
\usepackage{lmodern}

\usetheme[]{Madrid}
\usecolortheme[]{dolphin}
\usefonttheme[]{professionalfonts}
\ifPDFTeX\else
  % xetex/luatex font selection
\fi
% Use upquote if available, for straight quotes in verbatim environments
\IfFileExists{upquote.sty}{\usepackage{upquote}}{}
\IfFileExists{microtype.sty}{% use microtype if available
  \usepackage[]{microtype}
  \UseMicrotypeSet[protrusion]{basicmath} % disable protrusion for tt fonts
}{}
\makeatletter
\@ifundefined{KOMAClassName}{% if non-KOMA class
  \IfFileExists{parskip.sty}{%
    \usepackage{parskip}
  }{% else
    \setlength{\parindent}{0pt}
    \setlength{\parskip}{6pt plus 2pt minus 1pt}}
}{% if KOMA class
  \KOMAoptions{parskip=half}}
\makeatother

\usepackage{color}
\usepackage{fancyvrb}
\newcommand{\VerbBar}{|}
\newcommand{\VERB}{\Verb[commandchars=\\\{\}]}
\DefineVerbatimEnvironment{Highlighting}{Verbatim}{commandchars=\\\{\}}
% Add ',fontsize=\small' for more characters per line
\usepackage{framed}
\definecolor{shadecolor}{RGB}{241,243,245}
\newenvironment{Shaded}{\begin{snugshade}}{\end{snugshade}}
\newcommand{\AlertTok}[1]{\textcolor[rgb]{0.68,0.00,0.00}{#1}}
\newcommand{\AnnotationTok}[1]{\textcolor[rgb]{0.37,0.37,0.37}{#1}}
\newcommand{\AttributeTok}[1]{\textcolor[rgb]{0.40,0.45,0.13}{#1}}
\newcommand{\BaseNTok}[1]{\textcolor[rgb]{0.68,0.00,0.00}{#1}}
\newcommand{\BuiltInTok}[1]{\textcolor[rgb]{0.00,0.23,0.31}{#1}}
\newcommand{\CharTok}[1]{\textcolor[rgb]{0.13,0.47,0.30}{#1}}
\newcommand{\CommentTok}[1]{\textcolor[rgb]{0.37,0.37,0.37}{#1}}
\newcommand{\CommentVarTok}[1]{\textcolor[rgb]{0.37,0.37,0.37}{\textit{#1}}}
\newcommand{\ConstantTok}[1]{\textcolor[rgb]{0.56,0.35,0.01}{#1}}
\newcommand{\ControlFlowTok}[1]{\textcolor[rgb]{0.00,0.23,0.31}{\textbf{#1}}}
\newcommand{\DataTypeTok}[1]{\textcolor[rgb]{0.68,0.00,0.00}{#1}}
\newcommand{\DecValTok}[1]{\textcolor[rgb]{0.68,0.00,0.00}{#1}}
\newcommand{\DocumentationTok}[1]{\textcolor[rgb]{0.37,0.37,0.37}{\textit{#1}}}
\newcommand{\ErrorTok}[1]{\textcolor[rgb]{0.68,0.00,0.00}{#1}}
\newcommand{\ExtensionTok}[1]{\textcolor[rgb]{0.00,0.23,0.31}{#1}}
\newcommand{\FloatTok}[1]{\textcolor[rgb]{0.68,0.00,0.00}{#1}}
\newcommand{\FunctionTok}[1]{\textcolor[rgb]{0.28,0.35,0.67}{#1}}
\newcommand{\ImportTok}[1]{\textcolor[rgb]{0.00,0.46,0.62}{#1}}
\newcommand{\InformationTok}[1]{\textcolor[rgb]{0.37,0.37,0.37}{#1}}
\newcommand{\KeywordTok}[1]{\textcolor[rgb]{0.00,0.23,0.31}{\textbf{#1}}}
\newcommand{\NormalTok}[1]{\textcolor[rgb]{0.00,0.23,0.31}{#1}}
\newcommand{\OperatorTok}[1]{\textcolor[rgb]{0.37,0.37,0.37}{#1}}
\newcommand{\OtherTok}[1]{\textcolor[rgb]{0.00,0.23,0.31}{#1}}
\newcommand{\PreprocessorTok}[1]{\textcolor[rgb]{0.68,0.00,0.00}{#1}}
\newcommand{\RegionMarkerTok}[1]{\textcolor[rgb]{0.00,0.23,0.31}{#1}}
\newcommand{\SpecialCharTok}[1]{\textcolor[rgb]{0.37,0.37,0.37}{#1}}
\newcommand{\SpecialStringTok}[1]{\textcolor[rgb]{0.13,0.47,0.30}{#1}}
\newcommand{\StringTok}[1]{\textcolor[rgb]{0.13,0.47,0.30}{#1}}
\newcommand{\VariableTok}[1]{\textcolor[rgb]{0.07,0.07,0.07}{#1}}
\newcommand{\VerbatimStringTok}[1]{\textcolor[rgb]{0.13,0.47,0.30}{#1}}
\newcommand{\WarningTok}[1]{\textcolor[rgb]{0.37,0.37,0.37}{\textit{#1}}}

\providecommand{\faExclamation}{\faIcon{exclamation}}
% O si prefieres que renderice el triángulo de advertencia:
% \providecommand{\faExclamation}{\faExclamationTriangle}

\usepackage{longtable,booktabs,array}
\newcounter{none} % for unnumbered tables
\usepackage{calc} % for calculating minipage widths
\usepackage{caption}
% Make caption package work with longtable
\makeatletter
\def\fnum@table{\tablename~\thetable}
\makeatother
\usepackage{graphicx}
\makeatletter
\newsavebox\pandoc@box
\newcommand*\pandocbounded[1]{% scales image to fit in text height/width
  \sbox\pandoc@box{#1}%
  \Gscale@div\@tempa{\textheight}{\dimexpr\ht\pandoc@box+\dp\pandoc@box\relax}%
  \Gscale@div\@tempb{\linewidth}{\wd\pandoc@box}%
  \ifdim\@tempb\p@<\@tempa\p@\let\@tempa\@tempb\fi% select the smaller of both
  \ifdim\@tempa\p@<\p@\scalebox{\@tempa}{\usebox\pandoc@box}%
  \else\usebox{\pandoc@box}%
  \fi%
}
% Set default figure placement to htbp
\def\fps@figure{htbp}
\makeatother


% definitions for citeproc citations
\NewDocumentCommand\citeproctext{}{}
\NewDocumentCommand\citeproc{mm}{%
  \begingroup\def\citeproctext{#2}\cite{#1}\endgroup}
\makeatletter
 % allow citations to break across lines
 \let\@cite@ofmt\@firstofone
 % avoid brackets around text for \cite:
 \def\@biblabel#1{}
 \def\@cite#1#2{{#1\if@tempswa , #2\fi}}
\makeatother
\newlength{\cslhangindent}
\setlength{\cslhangindent}{1.5em}
\newlength{\csllabelwidth}
\setlength{\csllabelwidth}{3em}
\newenvironment{CSLReferences}[2] % #1 hanging-indent, #2 entry-spacing
 {\begin{list}{}{%
  \setlength{\itemindent}{0pt}
  \setlength{\leftmargin}{0pt}
  \setlength{\parsep}{0pt}
  % turn on hanging indent if param 1 is 1
  \ifodd #1
   \setlength{\leftmargin}{\cslhangindent}
   \setlength{\itemindent}{-1\cslhangindent}
  \fi
  % set entry spacing
  \setlength{\itemsep}{#2\baselineskip}}}
 {\end{list}}
\usepackage{calc}
\newcommand{\CSLBlock}[1]{\hfill\break\parbox[t]{\linewidth}{\strut\ignorespaces#1\strut}}
\newcommand{\CSLLeftMargin}[1]{\parbox[t]{\csllabelwidth}{\strut#1\strut}}
\newcommand{\CSLRightInline}[1]{\parbox[t]{\linewidth - \csllabelwidth}{\strut#1\strut}}
\newcommand{\CSLIndent}[1]{\hspace{\cslhangindent}#1}

\ifLuaTeX
\usepackage[bidi=basic,shorthands=off]{babel}
\else
\usepackage[bidi=default,shorthands=off]{babel}
\fi
\ifLuaTeX
  \usepackage{selnolig} % disable illegal ligatures
\fi


\setlength{\emergencystretch}{3em} % prevent overfull lines

\providecommand{\tightlist}{%
  \setlength{\itemsep}{0pt}\setlength{\parskip}{0pt}}



 


\usepackage[T1]{fontenc}
\usepackage{helvet}
\renewcommand{\familydefault}{\sfdefault}
\usepackage{booktabs}
\usepackage{amsmath,amssymb}
\usepackage{microtype}
\usepackage{ragged2e}
\definecolor{BioNavy}{HTML}{0B3954}
\definecolor{BioTeal}{HTML}{006D77}
\definecolor{BioGold}{HTML}{E09F3E}
\definecolor{BioPaper}{HTML}{F7F9FC}
\setbeamercolor{structure}{fg=BioNavy}
\setbeamercolor{frametitle}{fg=white,bg=BioNavy}
\setbeamercolor{title}{fg=BioNavy}
\setbeamercolor{block title}{fg=white,bg=BioTeal}
\setbeamercolor{block body}{fg=black,bg=BioTeal!8}
\setbeamercolor{alerted text}{fg=BioGold!80!black}
\setbeamertemplate{navigation symbols}{}
\setbeamertemplate{itemize item}{\color{BioGold}\small$\blacktriangleright$}
\setbeamertemplate{itemize subitem}{\color{BioTeal}\scriptsize$\bullet$}
\setbeamersize{text margin left=7mm,text margin right=7mm}
\setbeamerfont{frametitle}{size=\Large,series=\bfseries}
\setbeamerfont{normal text}{size=\small}
\setbeamertemplate{footline}{%
  \leavevmode\hbox{%
  \begin{beamercolorbox}[wd=.84\paperwidth,ht=2.6ex,dp=1ex,leftskip=7mm]{author in head/foot}%
  \color{BioNavy}\scriptsize SYSB · Semanas 1--5
  \end{beamercolorbox}%
  \begin{beamercolorbox}[wd=.16\paperwidth,ht=2.6ex,dp=1ex,center]{date in head/foot}%
  \color{BioNavy}\scriptsize\insertframenumber/\inserttotalframenumber
  \end{beamercolorbox}}}
\makeatletter
\@ifpackageloaded{tcolorbox}{}{\usepackage[skins,breakable]{tcolorbox}}
\@ifpackageloaded{fontawesome5}{}{\usepackage{fontawesome5}}
\definecolor{quarto-callout-color}{HTML}{909090}
\definecolor{quarto-callout-note-color}{HTML}{0758E5}
\definecolor{quarto-callout-important-color}{HTML}{CC1914}
\definecolor{quarto-callout-warning-color}{HTML}{EB9113}
\definecolor{quarto-callout-tip-color}{HTML}{00A047}
\definecolor{quarto-callout-caution-color}{HTML}{FC5300}
\definecolor{quarto-callout-color-frame}{HTML}{acacac}
\definecolor{quarto-callout-note-color-frame}{HTML}{4582ec}
\definecolor{quarto-callout-important-color-frame}{HTML}{d9534f}
\definecolor{quarto-callout-warning-color-frame}{HTML}{f0ad4e}
\definecolor{quarto-callout-tip-color-frame}{HTML}{02b875}
\definecolor{quarto-callout-caution-color-frame}{HTML}{fd7e14}
\makeatother
\makeatletter
\@ifpackageloaded{caption}{}{\usepackage{caption}}
\AtBeginDocument{%
\ifdefined\contentsname
  \renewcommand*\contentsname{Tabla de contenidos}
\else
  \newcommand\contentsname{Tabla de contenidos}
\fi
\ifdefined\listfigurename
  \renewcommand*\listfigurename{Listado de Figuras}
\else
  \newcommand\listfigurename{Listado de Figuras}
\fi
\ifdefined\listtablename
  \renewcommand*\listtablename{Listado de Tablas}
\else
  \newcommand\listtablename{Listado de Tablas}
\fi
\ifdefined\figurename
  \renewcommand*\figurename{Figura}
\else
  \newcommand\figurename{Figura}
\fi
\ifdefined\tablename
  \renewcommand*\tablename{Tabla}
\else
  \newcommand\tablename{Tabla}
\fi
}
\@ifpackageloaded{float}{}{\usepackage{float}}
\floatstyle{ruled}
\@ifundefined{c@chapter}{\newfloat{codelisting}{h}{lop}}{\newfloat{codelisting}{h}{lop}[chapter]}
\floatname{codelisting}{Listado}
\newcommand*\listoflistings{\listof{codelisting}{Listado de Listados}}
\makeatother
\makeatletter
\makeatother
\makeatletter
\@ifpackageloaded{caption}{}{\usepackage{caption}}
\@ifpackageloaded{subcaption}{}{\usepackage{subcaption}}
\makeatother

\usepackage{bookmark}
\IfFileExists{xurl.sty}{\usepackage{xurl}}{} % add URL line breaks if available
\urlstyle{same}
\hypersetup{
  pdftitle={Sistemas y Señales Biomédicos},
  pdfauthor={Ph. D. Pablo Eduardo Caicedo-Rodríguez},
  pdflang={es},
  hidelinks,
  pdfcreator={LaTeX via pandoc}}


\title{\texorpdfstring{Sistemas y Señales
Biomédicos}{Sistemas y Señales Biomédicos}}
\subtitle{\texorpdfstring{Del Espectro Finito hasta los Sistemas
LTI}{Del Espectro Finito hasta los Sistemas LTI}}
\author{Ph. D. Pablo Eduardo Caicedo-Rodríguez}
\date{}
\logo{\includegraphics{../../../recursos/imagenes/generales/Escuela\_Rosario\_logo.png}}

\begin{document}
\frame{\titlepage}


\section{Cómo estudiar las semanas
  6--10}\label{cuxf3mo-estudiar-las-semanas-610}

\begin{frame}{Del espectro ideal a decisiones reproducibles}
  \protect\phantomsection\label{del-espectro-ideal-a-decisiones-reproducibles}
  Al terminar la semana 10, el estudiante debe poder:

  \begin{enumerate}
    \tightlist
    \item
          predecir el efecto de operaciones temporales sobre la transformada de
          Fourier;
    \item
          calcular e interpretar DFT y FFT con escalamiento explícito;
    \item
          explicar resolución frecuencial, fuga y efecto de las ventanas;
    \item
          relacionar adquisición, muestreo y almacenamiento con el espectro
          observado;
    \item
          estimar PSD mediante periodograma y Welch;
    \item
          analizar señales no estacionarias con STFT;
    \item
          verificar linealidad, invariancia temporal y respuesta al impulso.
  \end{enumerate}

  \begin{tcolorbox}[enhanced jigsaw, arc=.35mm, bottomrule=.15mm, bottomtitle=1mm, breakable, colback=white, colbacktitle=quarto-callout-important-color!10!white, colframe=quarto-callout-important-color-frame, coltitle=black, left=2mm, leftrule=.75mm, opacityback=0, opacitybacktitle=0.6, rightrule=.15mm, title=\textcolor{quarto-callout-important-color}{\faExclamation}\hspace{0.5em}{Idea conductora}, titlerule=0mm, toprule=.15mm, toptitle=1mm]

    El espectro calculado depende tanto de la señal fisiológica como de la
    adquisición, la duración observada y el estimador seleccionado.

  \end{tcolorbox}
\end{frame}

\begin{frame}{Mapa de las diez sesiones}
  \protect\phantomsection\label{mapa-de-las-diez-sesiones}
  {\def\LTcaptype{none} % do not increment counter
    \begin{longtable}[]{@{}
      >{\raggedright\arraybackslash}p{(\linewidth - 4\tabcolsep) * \real{0.3333}}
      >{\raggedright\arraybackslash}p{(\linewidth - 4\tabcolsep) * \real{0.3333}}
      >{\raggedright\arraybackslash}p{(\linewidth - 4\tabcolsep) * \real{0.3333}}@{}}
      \toprule\noalign{}
      \begin{minipage}[b]{\linewidth}\raggedright
        Semana
      \end{minipage} & \begin{minipage}[b]{\linewidth}\raggedright
                         Sesión 1
                       \end{minipage} & \begin{minipage}[b]{\linewidth}\raggedright
                                          Sesión 2
                                        \end{minipage}                                                    \\
      \midrule\noalign{}
      \endhead
      6                                           & Propiedades de Fourier                      & DFT y FFT                           \\
      7                                           & Resolución, fuga y ventanas                 & Cadena de adquisición               \\
      8                                           & Muestreo, Nyquist y aliasing                & Potencia, PSD y periodograma        \\
      9                                           & Promediado y Welch                          & No estacionariedad y STFT           \\
      10                                          & Ventanas tiempo-frecuencia                  & Sistemas LTI y respuesta al impulso \\
      \bottomrule\noalign{}
    \end{longtable}
  }

  Cada sesión integra fundamento matemático, interpretación biomédica,
  ejemplo computacional y verificación conceptual.
\end{frame}

\begin{frame}{Convenciones del bloque}
  \protect\phantomsection\label{convenciones-del-bloque}
  Transformada continua en frecuencia ordinaria:

  \[X(f)=\int_{-\infty}^{\infty}x(t)e^{-j2\pi ft}\,dt\]

  DFT de \(N\) muestras:

  \[X[k]=\sum_{n=0}^{N-1}x[n]e^{-j2\pi kn/N}\]

  \begin{tcolorbox}[enhanced jigsaw, arc=.35mm, bottomrule=.15mm, bottomtitle=1mm, breakable, colback=white, colbacktitle=quarto-callout-warning-color!10!white, colframe=quarto-callout-warning-color-frame, coltitle=black, left=2mm, leftrule=.75mm, opacityback=0, opacitybacktitle=0.6, rightrule=.15mm, title=\textcolor{quarto-callout-warning-color}{\faExclamationTriangle}\hspace{0.5em}{Advertencia}, titlerule=0mm, toprule=.15mm, toptitle=1mm]

    Las constantes de normalización cambian entre textos y bibliotecas. Toda
    interpretación cuantitativa debe declarar la convención utilizada.

  \end{tcolorbox}
\end{frame}

\section{Semana 6 · Propiedades de Fourier y
  DFT}\label{semana-6-propiedades-de-fourier-y-dft}

\begin{frame}{Sesión 11 · Las propiedades permiten razonar sin
    reintegrar}
  \protect\phantomsection\label{sesiuxf3n-11-las-propiedades-permiten-razonar-sin-reintegrar}
  Si conocemos el par

  \[x(t)\longleftrightarrow X(f),\]

  podemos anticipar la transformada de una señal modificada mediante
  propiedades algebraicas.

  \begin{tcolorbox}[enhanced jigsaw, arc=.35mm, bottomrule=.15mm, bottomtitle=1mm, breakable, colback=white, colbacktitle=quarto-callout-important-color!10!white, colframe=quarto-callout-important-color-frame, coltitle=black, left=2mm, leftrule=.75mm, opacityback=0, opacitybacktitle=0.6, rightrule=.15mm, title=\textcolor{quarto-callout-important-color}{\faExclamation}\hspace{0.5em}{Propósito}, titlerule=0mm, toprule=.15mm, toptitle=1mm]

    Las propiedades no reemplazan la interpretación. Permiten conectar una
    operación en tiempo con su consecuencia en magnitud y fase.

  \end{tcolorbox}
\end{frame}

\begin{frame}{Linealidad conserva la combinación}
  \protect\phantomsection\label{linealidad-conserva-la-combinaciuxf3n}
  Si

  \[x_1(t)\longleftrightarrow X_1(f),\qquad x_2(t)\longleftrightarrow X_2(f),\]

  entonces

  \[a x_1(t)+b x_2(t)\longleftrightarrow aX_1(f)+bX_2(f).\]

  \begin{tcolorbox}[enhanced jigsaw, arc=.35mm, bottomrule=.15mm, bottomtitle=1mm, breakable, colback=white, colbacktitle=quarto-callout-note-color!10!white, colframe=quarto-callout-note-color-frame, coltitle=black, left=2mm, leftrule=.75mm, opacityback=0, opacitybacktitle=0.6, rightrule=.15mm, title=\textcolor{quarto-callout-note-color}{\faInfo}\hspace{0.5em}{Nota}, titlerule=0mm, toprule=.15mm, toptitle=1mm]

    En un ECG con deriva e interferencia de red, el espectro observado
    combina las transformadas de cada componente solo si el modelo aditivo
    resulta razonable.

  \end{tcolorbox}
\end{frame}

\begin{frame}{Linealidad no separa fuentes automáticamente}
  \protect\phantomsection\label{linealidad-no-separa-fuentes-automuxe1ticamente}
  Modelo aditivo:

  \[x_{obs}(t)=s(t)+i(t)+a(t).\]

  \begin{itemize}
    \tightlist
    \item
          \(s(t)\): señal de interés.
    \item
          \(i(t)\): interferencia estructurada.
    \item
          \(a(t)\): artefacto.
  \end{itemize}

  Entonces \(X_{obs}(f)=S(f)+I(f)+A(f)\).

  \begin{tcolorbox}[enhanced jigsaw, arc=.35mm, bottomrule=.15mm, bottomtitle=1mm, breakable, colback=white, colbacktitle=quarto-callout-warning-color!10!white, colframe=quarto-callout-warning-color-frame, coltitle=black, left=2mm, leftrule=.75mm, opacityback=0, opacitybacktitle=0.6, rightrule=.15mm, title=\textcolor{quarto-callout-warning-color}{\faExclamationTriangle}\hspace{0.5em}{Advertencia}, titlerule=0mm, toprule=.15mm, toptitle=1mm]

    La suma espectral puede producir cancelación por fase. La magnitud de la
    suma no equivale a sumar magnitudes.

  \end{tcolorbox}
\end{frame}

\begin{frame}{Un retraso añade fase lineal}
  \protect\phantomsection\label{un-retraso-auxf1ade-fase-lineal}
  \[x(t-t_0)\longleftrightarrow X(f)e^{-j2\pi ft_0}\]

  Por tanto:

  \[|X_{ret}(f)|=|X(f)|,\]

  \[\angle X_{ret}(f)=\angle X(f)-2\pi ft_0.\]

  \begin{tcolorbox}[enhanced jigsaw, arc=.35mm, bottomrule=.15mm, bottomtitle=1mm, breakable, colback=white, colbacktitle=quarto-callout-important-color!10!white, colframe=quarto-callout-important-color-frame, coltitle=black, left=2mm, leftrule=.75mm, opacityback=0, opacitybacktitle=0.6, rightrule=.15mm, title=\textcolor{quarto-callout-important-color}{\faExclamation}\hspace{0.5em}{Importante}, titlerule=0mm, toprule=.15mm, toptitle=1mm]

    Un retraso puro conserva la magnitud. La información temporal aparece en
    la pendiente de fase.

  \end{tcolorbox}
\end{frame}

\begin{frame}{La fase lineal representa un retardo constante}
  \protect\phantomsection\label{la-fase-lineal-representa-un-retardo-constante}
  El retardo de grupo se define como

  \[\tau_g(f)=-\frac{1}{2\pi}\frac{d\phi(f)}{df}.\]

  Si \(\phi(f)=-2\pi ft_0\), entonces \(\tau_g(f)=t_0\).

  \begin{tcolorbox}[enhanced jigsaw, arc=.35mm, bottomrule=.15mm, bottomtitle=1mm, breakable, colback=white, colbacktitle=quarto-callout-note-color!10!white, colframe=quarto-callout-note-color-frame, coltitle=black, left=2mm, leftrule=.75mm, opacityback=0, opacitybacktitle=0.6, rightrule=.15mm, title=\textcolor{quarto-callout-note-color}{\faInfo}\hspace{0.5em}{Nota}, titlerule=0mm, toprule=.15mm, toptitle=1mm]

    Cuando todas las frecuencias sufren el mismo retardo, la forma temporal
    puede conservarse. Una fase no lineal puede deformar morfologías como el
    complejo QRS.

  \end{tcolorbox}
\end{frame}

\begin{frame}{Desplazar en frecuencia equivale a modular}
  \protect\phantomsection\label{desplazar-en-frecuencia-equivale-a-modular}
  \[x(t)e^{j2\pi f_0t}\longleftrightarrow X(f-f_0).\]

  Para modulación real:

  \[x(t)\cos(2\pi f_0t)
    \longleftrightarrow \frac{1}{2}X(f-f_0)+\frac{1}{2}X(f+f_0).\]

  \begin{tcolorbox}[enhanced jigsaw, arc=.35mm, bottomrule=.15mm, bottomtitle=1mm, breakable, colback=white, colbacktitle=quarto-callout-tip-color!10!white, colframe=quarto-callout-tip-color-frame, coltitle=black, left=2mm, leftrule=.75mm, opacityback=0, opacitybacktitle=0.6, rightrule=.15mm, title=\textcolor{quarto-callout-tip-color}{\faLightbulb}\hspace{0.5em}{Tip}, titlerule=0mm, toprule=.15mm, toptitle=1mm]

    La multiplicación por una portadora crea copias desplazadas del
    espectro. Este principio aparece en instrumentación y comunicaciones
    biomédicas.

  \end{tcolorbox}
\end{frame}

\begin{frame}{El escalamiento temporal modifica ancho y amplitud}
  \protect\phantomsection\label{el-escalamiento-temporal-modifica-ancho-y-amplitud}
  \[x(at)\longleftrightarrow \frac{1}{|a|}X\left(\frac{f}{a}\right).\]

  \begin{itemize}
    \tightlist
    \item
          \(|a|>1\): la señal se comprime en tiempo y su espectro se expande.
    \item
          \(0<|a|<1\): la señal se expande en tiempo y su espectro se comprime.
    \item
          \(a<0\): también ocurre inversión temporal.
  \end{itemize}

  \begin{tcolorbox}[enhanced jigsaw, arc=.35mm, bottomrule=.15mm, bottomtitle=1mm, breakable, colback=white, colbacktitle=quarto-callout-warning-color!10!white, colframe=quarto-callout-warning-color-frame, coltitle=black, left=2mm, leftrule=.75mm, opacityback=0, opacitybacktitle=0.6, rightrule=.15mm, title=\textcolor{quarto-callout-warning-color}{\faExclamationTriangle}\hspace{0.5em}{Advertencia}, titlerule=0mm, toprule=.15mm, toptitle=1mm]

    El factor \(1/|a|\) forma parte de la propiedad. Omitirlo cambia la
    escala de la transformada.

  \end{tcolorbox}
\end{frame}

\begin{frame}{Derivar resalta componentes rápidas}
  \protect\phantomsection\label{derivar-resalta-componentes-ruxe1pidas}
  \[\frac{dx(t)}{dt}\longleftrightarrow j2\pi fX(f).\]

  Consecuencias:

  \begin{itemize}
    \tightlist
    \item
          frecuencia cero queda anulada;
    \item
          las componentes altas reciben mayor ponderación;
    \item
          el ruido de alta frecuencia puede amplificarse.
  \end{itemize}

  \begin{tcolorbox}[enhanced jigsaw, arc=.35mm, bottomrule=.15mm, bottomtitle=1mm, breakable, colback=white, colbacktitle=quarto-callout-warning-color!10!white, colframe=quarto-callout-warning-color-frame, coltitle=black, left=2mm, leftrule=.75mm, opacityback=0, opacitybacktitle=0.6, rightrule=.15mm, title=\textcolor{quarto-callout-warning-color}{\faExclamationTriangle}\hspace{0.5em}{Advertencia}, titlerule=0mm, toprule=.15mm, toptitle=1mm]

    Un diferenciador ayuda a localizar cambios rápidos, pero puede
    deteriorar una bioseñal si no se controla el ruido.

  \end{tcolorbox}
\end{frame}

\begin{frame}{Integrar favorece variaciones lentas}
  \protect\phantomsection\label{integrar-favorece-variaciones-lentas}
  De forma conceptual, para condiciones adecuadas:

  \[\int_{-\infty}^{t}x(\tau)d\tau
    \longleftrightarrow \frac{X(f)}{j2\pi f}+\text{término en }\delta(f).\]

  \begin{tcolorbox}[enhanced jigsaw, arc=.35mm, bottomrule=.15mm, bottomtitle=1mm, breakable, colback=white, colbacktitle=quarto-callout-note-color!10!white, colframe=quarto-callout-note-color-frame, coltitle=black, left=2mm, leftrule=.75mm, opacityback=0, opacitybacktitle=0.6, rightrule=.15mm, title=\textcolor{quarto-callout-note-color}{\faInfo}\hspace{0.5em}{Nota}, titlerule=0mm, toprule=.15mm, toptitle=1mm]

    La singularidad en \(f=0\) obliga a considerar media, condiciones
    iniciales y convergencia. La integración no debe aplicarse como una
    división numérica directa por cero.

  \end{tcolorbox}
\end{frame}

\begin{frame}{La convolución se convierte en multiplicación}
  \protect\phantomsection\label{la-convoluciuxf3n-se-convierte-en-multiplicaciuxf3n}
  \[y(t)=x(t)*h(t)\]

  \[Y(f)=X(f)H(f).\]

  Interpretación:

  \begin{itemize}
    \tightlist
    \item
          \(X(f)\) describe el contenido de la entrada;
    \item
          \(H(f)\) describe la acción frecuencial del sistema;
    \item
          \(Y(f)\) conserva, atenúa o cambia la fase de cada componente.
  \end{itemize}

  \begin{tcolorbox}[enhanced jigsaw, arc=.35mm, bottomrule=.15mm, bottomtitle=1mm, breakable, colback=white, colbacktitle=quarto-callout-important-color!10!white, colframe=quarto-callout-important-color-frame, coltitle=black, left=2mm, leftrule=.75mm, opacityback=0, opacitybacktitle=0.6, rightrule=.15mm, title=\textcolor{quarto-callout-important-color}{\faExclamation}\hspace{0.5em}{Importante}, titlerule=0mm, toprule=.15mm, toptitle=1mm]

    El teorema de convolución conecta el análisis de señales con el análisis
    de sistemas LTI.

  \end{tcolorbox}
\end{frame}

\begin{frame}{Multiplicar en tiempo produce convolución espectral}
  \protect\phantomsection\label{multiplicar-en-tiempo-produce-convoluciuxf3n-espectral}
  \[y(t)=x(t)w(t)\]

  \[Y(f)=X(f)*W(f).\]

  \begin{tcolorbox}[enhanced jigsaw, arc=.35mm, bottomrule=.15mm, bottomtitle=1mm, breakable, colback=white, colbacktitle=quarto-callout-warning-color!10!white, colframe=quarto-callout-warning-color-frame, coltitle=black, left=2mm, leftrule=.75mm, opacityback=0, opacitybacktitle=0.6, rightrule=.15mm, title=\textcolor{quarto-callout-warning-color}{\faExclamationTriangle}\hspace{0.5em}{Advertencia}, titlerule=0mm, toprule=.15mm, toptitle=1mm]

    Seleccionar un segmento mediante una ventana mezcla el espectro de la
    señal con el de la ventana. Esta es la base matemática de la fuga
    espectral.

  \end{tcolorbox}
\end{frame}

\begin{frame}{Dualidad: tiempo y frecuencia intercambian estructuras}
  \protect\phantomsection\label{dualidad-tiempo-y-frecuencia-intercambian-estructuras}
  Bajo la convención empleada:

  \[x(t)\longleftrightarrow X(f)\]

  implica

  \[X(t)\longleftrightarrow x(-f).\]

  \begin{tcolorbox}[enhanced jigsaw, arc=.35mm, bottomrule=.15mm, bottomtitle=1mm, breakable, colback=white, colbacktitle=quarto-callout-note-color!10!white, colframe=quarto-callout-note-color-frame, coltitle=black, left=2mm, leftrule=.75mm, opacityback=0, opacitybacktitle=0.6, rightrule=.15mm, title=\textcolor{quarto-callout-note-color}{\faInfo}\hspace{0.5em}{Nota}, titlerule=0mm, toprule=.15mm, toptitle=1mm]

    La dualidad ayuda a generar pares de transformada, pero requiere
    atención a signos y factores según la convención.

  \end{tcolorbox}
\end{frame}

\begin{frame}{Tabla de predicción rápida}
  \protect\phantomsection\label{tabla-de-predicciuxf3n-ruxe1pida}
  {\def\LTcaptype{none} % do not increment counter
    \begin{longtable}[]{@{}ll@{}}
      \toprule\noalign{}
      Operación temporal  & Efecto principal en frecuencia \\
      \midrule\noalign{}
      \endhead
      suma ponderada      & suma compleja de espectros     \\
      retraso             & fase lineal                    \\
      compresión temporal & expansión espectral            \\
      derivada            & ponderación por \(j2\pi f\)    \\
      convolución         & producto de espectros          \\
      ventaneo            & convolución con \(W(f)\)       \\
      \bottomrule\noalign{}
    \end{longtable}
  }

  \begin{tcolorbox}[enhanced jigsaw, arc=.35mm, bottomrule=.15mm, bottomtitle=1mm, breakable, colback=white, colbacktitle=quarto-callout-tip-color!10!white, colframe=quarto-callout-tip-color-frame, coltitle=black, left=2mm, leftrule=.75mm, opacityback=0, opacitybacktitle=0.6, rightrule=.15mm, title=\textcolor{quarto-callout-tip-color}{\faLightbulb}\hspace{0.5em}{Tip}, titlerule=0mm, toprule=.15mm, toptitle=1mm]

    Antes de calcular, prediga qué aspectos deberían cambiar y cuáles
    deberían conservarse.

  \end{tcolorbox}
\end{frame}

\begin{frame}{Caso biomédico: tres modificaciones del ECG}
  \protect\phantomsection\label{caso-biomuxe9dico-tres-modificaciones-del-ecg}
  Considere un latido ideal \(x(t)\).

  \begin{enumerate}
    \tightlist
    \item
          \(x(t-0.2)\) conserva magnitud y añade fase lineal.
    \item
          \(x(2t)\) dura la mitad y extiende el contenido frecuencial.
    \item
          \(dx/dt\) enfatiza las pendientes rápidas del QRS.
  \end{enumerate}

  \begin{tcolorbox}[enhanced jigsaw, arc=.35mm, bottomrule=.15mm, bottomtitle=1mm, breakable, colback=white, colbacktitle=quarto-callout-warning-color!10!white, colframe=quarto-callout-warning-color-frame, coltitle=black, left=2mm, leftrule=.75mm, opacityback=0, opacitybacktitle=0.6, rightrule=.15mm, title=\textcolor{quarto-callout-warning-color}{\faExclamationTriangle}\hspace{0.5em}{Advertencia}, titlerule=0mm, toprule=.15mm, toptitle=1mm]

    Ninguna operación garantiza mejorar detección o diagnóstico. La utilidad
    depende del objetivo y de la contaminación presente.

  \end{tcolorbox}
\end{frame}

\begin{frame}{Verificación de la sesión 11}
  \protect\phantomsection\label{verificaciuxf3n-de-la-sesiuxf3n-11}
  Para \(y(t)=3x(2t-4)\):

  \begin{enumerate}
    \tightlist
    \item
          reescriba el argumento para identificar escalamiento y retraso;
    \item
          obtenga \(Y(f)\) a partir de \(X(f)\);
    \item
          indique qué ocurre con el ancho espectral;
    \item
          determine si la fase cambia por el desplazamiento.
  \end{enumerate}

  \begin{tcolorbox}[enhanced jigsaw, arc=.35mm, bottomrule=.15mm, bottomtitle=1mm, breakable, colback=white, colbacktitle=quarto-callout-tip-color!10!white, colframe=quarto-callout-tip-color-frame, coltitle=black, left=2mm, leftrule=.75mm, opacityback=0, opacitybacktitle=0.6, rightrule=.15mm, title=\textcolor{quarto-callout-tip-color}{\faLightbulb}\hspace{0.5em}{Criterio de logro}, titlerule=0mm, toprule=.15mm, toptitle=1mm]

    La solución debe incluir el factor de amplitud, el escalamiento
    frecuencial y la fase asociada al retardo efectivo.

  \end{tcolorbox}
\end{frame}

\begin{frame}{Sesión 12 · La DFT analiza un bloque finito}
  \protect\phantomsection\label{sesiuxf3n-12-la-dft-analiza-un-bloque-finito}
  La DFT opera sobre \(N\) muestras:

  \[X[k]=\sum_{n=0}^{N-1}x[n]e^{-j2\pi kn/N},\qquad k=0,\ldots,N-1.\]

  \begin{tcolorbox}[enhanced jigsaw, arc=.35mm, bottomrule=.15mm, bottomtitle=1mm, breakable, colback=white, colbacktitle=quarto-callout-important-color!10!white, colframe=quarto-callout-important-color-frame, coltitle=black, left=2mm, leftrule=.75mm, opacityback=0, opacitybacktitle=0.6, rightrule=.15mm, title=\textcolor{quarto-callout-important-color}{\faExclamation}\hspace{0.5em}{Importante}, titlerule=0mm, toprule=.15mm, toptitle=1mm]

    La DFT produce coeficientes de una representación discreta y periódica.
    No entrega un continuo de frecuencias.

  \end{tcolorbox}
\end{frame}

\begin{frame}{La transformada inversa reconstruye el bloque}
  \protect\phantomsection\label{la-transformada-inversa-reconstruye-el-bloque}
  \[x[n]=\frac{1}{N}\sum_{k=0}^{N-1}X[k]e^{j2\pi kn/N}.\]

  La DFT y la IDFT forman un par exacto salvo error numérico.

  \begin{tcolorbox}[enhanced jigsaw, arc=.35mm, bottomrule=.15mm, bottomtitle=1mm, breakable, colback=white, colbacktitle=quarto-callout-warning-color!10!white, colframe=quarto-callout-warning-color-frame, coltitle=black, left=2mm, leftrule=.75mm, opacityback=0, opacitybacktitle=0.6, rightrule=.15mm, title=\textcolor{quarto-callout-warning-color}{\faExclamationTriangle}\hspace{0.5em}{Advertencia}, titlerule=0mm, toprule=.15mm, toptitle=1mm]

    Esta convención asigna el factor \(1/N\) a la inversa. Otras
    convenciones reparten la normalización entre ambas transformadas.

  \end{tcolorbox}
\end{frame}

\begin{frame}{Cada índice corresponde a una frecuencia}
  \protect\phantomsection\label{cada-uxedndice-corresponde-a-una-frecuencia}
  Para muestreo uniforme a \(f_s\):

  \[f_k=\frac{k}{N}f_s,\qquad \Delta f=\frac{f_s}{N}=\frac{1}{T}.\]

  donde \(T=N/f_s\) es la duración.

  \begin{tcolorbox}[enhanced jigsaw, arc=.35mm, bottomrule=.15mm, bottomtitle=1mm, breakable, colback=white, colbacktitle=quarto-callout-important-color!10!white, colframe=quarto-callout-important-color-frame, coltitle=black, left=2mm, leftrule=.75mm, opacityback=0, opacitybacktitle=0.6, rightrule=.15mm, title=\textcolor{quarto-callout-important-color}{\faExclamation}\hspace{0.5em}{Importante}, titlerule=0mm, toprule=.15mm, toptitle=1mm]

    La separación entre bins depende de la duración observada. Aumentar
    \(f_s\) sin aumentar \(N\) no mejora la resolución frecuencial.

  \end{tcolorbox}
\end{frame}

\begin{frame}[fragile]{Los índices superiores representan frecuencias
    negativas}
  \protect\phantomsection\label{los-uxedndices-superiores-representan-frecuencias-negativas}
  Para una DFT bilateral:

  \[f_k=\begin{cases}
      kf_s/N,     & 0\leq k\leq N/2, \\
      (k-N)f_s/N, & N/2<k<N.
    \end{cases}\]

  \begin{tcolorbox}[enhanced jigsaw, arc=.35mm, bottomrule=.15mm, bottomtitle=1mm, breakable, colback=white, colbacktitle=quarto-callout-note-color!10!white, colframe=quarto-callout-note-color-frame, coltitle=black, left=2mm, leftrule=.75mm, opacityback=0, opacitybacktitle=0.6, rightrule=.15mm, title=\textcolor{quarto-callout-note-color}{\faInfo}\hspace{0.5em}{Nota}, titlerule=0mm, toprule=.15mm, toptitle=1mm]

    \texttt{np.fft.fftfreq} construye este eje y evita asignaciones manuales
    incorrectas.

  \end{tcolorbox}
\end{frame}

\begin{frame}[fragile]{Señales reales poseen simetría conjugada}
  \protect\phantomsection\label{seuxf1ales-reales-poseen-simetruxeda-conjugada}
  Si \(x[n]\in\mathbb{R}\):

  \[X[N-k]=X^*[k].\]

  Por tanto, la mitad no negativa contiene información suficiente, con
  cuidado especial en DC y Nyquist.

  \begin{tcolorbox}[enhanced jigsaw, arc=.35mm, bottomrule=.15mm, bottomtitle=1mm, breakable, colback=white, colbacktitle=quarto-callout-tip-color!10!white, colframe=quarto-callout-tip-color-frame, coltitle=black, left=2mm, leftrule=.75mm, opacityback=0, opacitybacktitle=0.6, rightrule=.15mm, title=\textcolor{quarto-callout-tip-color}{\faLightbulb}\hspace{0.5em}{Tip}, titlerule=0mm, toprule=.15mm, toptitle=1mm]

    \texttt{rfft} calcula solo las frecuencias no negativas de una entrada
    real y reduce costo y almacenamiento.

  \end{tcolorbox}
\end{frame}

\begin{frame}{La FFT es un algoritmo}
  \protect\phantomsection\label{la-fft-es-un-algoritmo}
  \begin{itemize}
    \tightlist
    \item
          DFT: transformación matemática.
    \item
          FFT: familia de algoritmos que calcula la DFT eficientemente.
    \item
          La FFT no cambia bins, resolución, fuga ni significado físico.
  \end{itemize}

  Complejidad aproximada:

  \[\mathcal{O}(N^2)\quad\text{frente a}\quad\mathcal{O}(N\log N).\]

  \begin{tcolorbox}[enhanced jigsaw, arc=.35mm, bottomrule=.15mm, bottomtitle=1mm, breakable, colback=white, colbacktitle=quarto-callout-warning-color!10!white, colframe=quarto-callout-warning-color-frame, coltitle=black, left=2mm, leftrule=.75mm, opacityback=0, opacitybacktitle=0.6, rightrule=.15mm, title=\textcolor{quarto-callout-warning-color}{\faExclamationTriangle}\hspace{0.5em}{Advertencia}, titlerule=0mm, toprule=.15mm, toptitle=1mm]

    Una FFT más rápida no corrige un segmento mal elegido ni una frecuencia
    de muestreo inadecuada.

  \end{tcolorbox}
\end{frame}

\begin{frame}{La DFT supone extensión periódica del bloque}
  \protect\phantomsection\label{la-dft-supone-extensiuxf3n-periuxf3dica-del-bloque}
  El bloque \(x[0],\ldots,x[N-1]\) se interpreta como un periodo de una
  secuencia periódica.

  Si los extremos no enlazan suavemente, la repetición introduce
  discontinuidades.

  \begin{tcolorbox}[enhanced jigsaw, arc=.35mm, bottomrule=.15mm, bottomtitle=1mm, breakable, colback=white, colbacktitle=quarto-callout-important-color!10!white, colframe=quarto-callout-important-color-frame, coltitle=black, left=2mm, leftrule=.75mm, opacityback=0, opacitybacktitle=0.6, rightrule=.15mm, title=\textcolor{quarto-callout-important-color}{\faExclamation}\hspace{0.5em}{Importante}, titlerule=0mm, toprule=.15mm, toptitle=1mm]

    La fuga espectral surge porque el bloque finito y su extensión periódica
    rara vez contienen un número entero de ciclos.

  \end{tcolorbox}
\end{frame}

\begin{frame}{Magnitud bilateral y unilateral}
  \protect\phantomsection\label{magnitud-bilateral-y-unilateral}
  Con \(X[k]=\mathrm{DFT}\{x[n]\}\):

  \[A_{2}[k]=\frac{|X[k]|}{N}.\]

  Para espectro unilateral de una señal real, se duplican los bins
  interiores:

  \[A_{1}[k]=2A_{2}[k],\]

  excepto DC y, si existe, Nyquist.

  \begin{tcolorbox}[enhanced jigsaw, arc=.35mm, bottomrule=.15mm, bottomtitle=1mm, breakable, colback=white, colbacktitle=quarto-callout-warning-color!10!white, colframe=quarto-callout-warning-color-frame, coltitle=black, left=2mm, leftrule=.75mm, opacityback=0, opacitybacktitle=0.6, rightrule=.15mm, title=\textcolor{quarto-callout-warning-color}{\faExclamationTriangle}\hspace{0.5em}{Advertencia}, titlerule=0mm, toprule=.15mm, toptitle=1mm]

    Duplicar todos los bins produce una amplitud incorrecta en los extremos
    del espectro.

  \end{tcolorbox}
\end{frame}

\begin{frame}{La fase necesita una referencia de amplitud}
  \protect\phantomsection\label{la-fase-necesita-una-referencia-de-amplitud}
  \[\phi[k]=\operatorname{angle}(X[k]).\]

  En bins donde \(|X[k]|\) es casi cero, la fase resulta numéricamente
  inestable.

  \begin{tcolorbox}[enhanced jigsaw, arc=.35mm, bottomrule=.15mm, bottomtitle=1mm, breakable, colback=white, colbacktitle=quarto-callout-note-color!10!white, colframe=quarto-callout-note-color-frame, coltitle=black, left=2mm, leftrule=.75mm, opacityback=0, opacitybacktitle=0.6, rightrule=.15mm, title=\textcolor{quarto-callout-note-color}{\faInfo}\hspace{0.5em}{Nota}, titlerule=0mm, toprule=.15mm, toptitle=1mm]

    La fase debe interpretarse solo donde la magnitud o coherencia aporten
    evidencia suficiente.

  \end{tcolorbox}
\end{frame}

\begin{frame}{Convolución circular y periodicidad}
  \protect\phantomsection\label{convoluciuxf3n-circular-y-periodicidad}
  La multiplicación de DFT corresponde a convolución circular:

  \[\mathrm{IDFT}\{X[k]H[k]\}=x[n]\circledast h[n].\]

  Para obtener convolución lineal sin solapamiento circular, se usa
  relleno con ceros hasta al menos \(N_x+N_h-1\).

  \begin{tcolorbox}[enhanced jigsaw, arc=.35mm, bottomrule=.15mm, bottomtitle=1mm, breakable, colback=white, colbacktitle=quarto-callout-important-color!10!white, colframe=quarto-callout-important-color-frame, coltitle=black, left=2mm, leftrule=.75mm, opacityback=0, opacitybacktitle=0.6, rightrule=.15mm, title=\textcolor{quarto-callout-important-color}{\faExclamation}\hspace{0.5em}{Importante}, titlerule=0mm, toprule=.15mm, toptitle=1mm]

    El relleno evita plegamiento circular en convolución, pero no añade
    información espectral al registro.

  \end{tcolorbox}
\end{frame}

\begin{frame}[fragile]{Código: eje y espectro unilateral}
  \protect\phantomsection\label{cuxf3digo-eje-y-espectro-unilateral}
  \begin{Shaded}
    \begin{Highlighting}[]
      \ImportTok{import}\NormalTok{ numpy }\ImportTok{as}\NormalTok{ np}

      \NormalTok{x0 }\OperatorTok{=}\NormalTok{ x }\OperatorTok{{-}}\NormalTok{ np.mean(x)}
      \NormalTok{N }\OperatorTok{=} \BuiltInTok{len}\NormalTok{(x0)}
      \NormalTok{X }\OperatorTok{=}\NormalTok{ np.fft.rfft(x0)}
      \NormalTok{f }\OperatorTok{=}\NormalTok{ np.fft.rfftfreq(N, d}\OperatorTok{=}\DecValTok{1}\OperatorTok{/}\NormalTok{fs)}
      \NormalTok{A }\OperatorTok{=}\NormalTok{ np.}\BuiltInTok{abs}\NormalTok{(X) }\OperatorTok{/}\NormalTok{ N}
      \ControlFlowTok{if}\NormalTok{ N }\OperatorTok{\textgreater{}} \DecValTok{2}\NormalTok{:}
      \NormalTok{    A[}\DecValTok{1}\NormalTok{:}\OperatorTok{{-}}\DecValTok{1}\NormalTok{] }\OperatorTok{*=} \DecValTok{2}
      \NormalTok{fase }\OperatorTok{=}\NormalTok{ np.angle(X)}
    \end{Highlighting}
  \end{Shaded}

  \begin{tcolorbox}[enhanced jigsaw, arc=.35mm, bottomrule=.15mm, bottomtitle=1mm, breakable, colback=white, colbacktitle=quarto-callout-warning-color!10!white, colframe=quarto-callout-warning-color-frame, coltitle=black, left=2mm, leftrule=.75mm, opacityback=0, opacitybacktitle=0.6, rightrule=.15mm, title=\textcolor{quarto-callout-warning-color}{\faExclamationTriangle}\hspace{0.5em}{Advertencia}, titlerule=0mm, toprule=.15mm, toptitle=1mm]

    Para \(N\) impar, el último bin no corresponde exactamente a Nyquist y
    también requiere duplicación. El código debe tratar paridad de forma
    explícita.

  \end{tcolorbox}
\end{frame}

\begin{frame}[fragile]{Código robusto para \(N\) par o impar}
  \protect\phantomsection\label{cuxf3digo-robusto-para-n-par-o-impar}
  \begin{Shaded}
    \begin{Highlighting}[]
      \NormalTok{X }\OperatorTok{=}\NormalTok{ np.fft.rfft(x0)}
      \NormalTok{A }\OperatorTok{=}\NormalTok{ np.}\BuiltInTok{abs}\NormalTok{(X) }\OperatorTok{/}\NormalTok{ N}

      \ControlFlowTok{if}\NormalTok{ N }\OperatorTok{\%} \DecValTok{2} \OperatorTok{==} \DecValTok{0}\NormalTok{:      }\CommentTok{\# existe bin de Nyquist}
      \NormalTok{    A[}\DecValTok{1}\NormalTok{:}\OperatorTok{{-}}\DecValTok{1}\NormalTok{] }\OperatorTok{*=} \DecValTok{2}
      \ControlFlowTok{else}\NormalTok{:}
      \NormalTok{    A[}\DecValTok{1}\NormalTok{:] }\OperatorTok{*=} \DecValTok{2}
    \end{Highlighting}
  \end{Shaded}

  \begin{tcolorbox}[enhanced jigsaw, arc=.35mm, bottomrule=.15mm, bottomtitle=1mm, breakable, colback=white, colbacktitle=quarto-callout-tip-color!10!white, colframe=quarto-callout-tip-color-frame, coltitle=black, left=2mm, leftrule=.75mm, opacityback=0, opacitybacktitle=0.6, rightrule=.15mm, title=\textcolor{quarto-callout-tip-color}{\faLightbulb}\hspace{0.5em}{Tip}, titlerule=0mm, toprule=.15mm, toptitle=1mm]

    Una prueba útil consiste en aplicar el código a una cosenoide de
    amplitud conocida alineada con un bin.

  \end{tcolorbox}
\end{frame}

\begin{frame}[fragile]{Verificación de la sesión 12}
  \protect\phantomsection\label{verificaciuxf3n-de-la-sesiuxf3n-12}
  Un registro tiene \(f_s=500\) Hz y \(N=2000\).

  \begin{enumerate}
    \tightlist
    \item
          calcule su duración;
    \item
          calcule \(\Delta f\);
    \item
          identifique el rango de \texttt{rfft};
    \item
          explique qué mejora y qué no mejora al usar 4096 puntos mediante zero
          padding;
    \item
          indique por qué la fase de bins pequeños puede ser inestable.
  \end{enumerate}

  \begin{tcolorbox}[enhanced jigsaw, arc=.35mm, bottomrule=.15mm, bottomtitle=1mm, breakable, colback=white, colbacktitle=quarto-callout-tip-color!10!white, colframe=quarto-callout-tip-color-frame, coltitle=black, left=2mm, leftrule=.75mm, opacityback=0, opacitybacktitle=0.6, rightrule=.15mm, title=\textcolor{quarto-callout-tip-color}{\faLightbulb}\hspace{0.5em}{Criterio de logro}, titlerule=0mm, toprule=.15mm, toptitle=1mm]

    Debe distinguir duración, resolución, densidad gráfica y costo
    computacional.

  \end{tcolorbox}
\end{frame}

\begin{frame}{Cierre de la semana 6}
  \protect\phantomsection\label{cierre-de-la-semana-6}
  \begin{itemize}
    \tightlist
    \item
          Las propiedades conectan operaciones temporales con cambios
          espectrales.
    \item
          El retraso actúa sobre fase; el escalamiento temporal modifica el
          ancho.
    \item
          La convolución temporal equivale a multiplicación frecuencial.
    \item
          La DFT representa un bloque finito sobre bins discretos.
    \item
          La FFT acelera el cálculo sin alterar la interpretación.
  \end{itemize}

  \begin{tcolorbox}[enhanced jigsaw, arc=.35mm, bottomrule=.15mm, bottomtitle=1mm, breakable, colback=white, colbacktitle=quarto-callout-note-color!10!white, colframe=quarto-callout-note-color-frame, coltitle=black, left=2mm, leftrule=.75mm, opacityback=0, opacitybacktitle=0.6, rightrule=.15mm, title=\textcolor{quarto-callout-note-color}{\faInfo}\hspace{0.5em}{Pregunta de salida}, titlerule=0mm, toprule=.15mm, toptitle=1mm]

    ¿Por qué dos registros con la misma frecuencia de muestreo pueden tener
    resoluciones frecuenciales distintas?

  \end{tcolorbox}
\end{frame}

\section{Semana 7 · Resolución, fuga y
  adquisición}\label{semana-7-resoluciuxf3n-fuga-y-adquisiciuxf3n}

\begin{frame}{Sesión 13 · La resolución proviene de la duración}
  \protect\phantomsection\label{sesiuxf3n-13-la-resoluciuxf3n-proviene-de-la-duraciuxf3n}
  \[\Delta f=\frac{f_s}{N}=\frac{1}{T}.\]

  Ejemplo: \(T=10\) s produce bins separados 0.1 Hz, con independencia de
  la forma usada para calcular \(N=f_sT\).

  \begin{tcolorbox}[enhanced jigsaw, arc=.35mm, bottomrule=.15mm, bottomtitle=1mm, breakable, colback=white, colbacktitle=quarto-callout-important-color!10!white, colframe=quarto-callout-important-color-frame, coltitle=black, left=2mm, leftrule=.75mm, opacityback=0, opacitybacktitle=0.6, rightrule=.15mm, title=\textcolor{quarto-callout-important-color}{\faExclamation}\hspace{0.5em}{Importante}, titlerule=0mm, toprule=.15mm, toptitle=1mm]

    Observar durante más tiempo permite separar componentes más próximas,
    siempre que la señal mantenga propiedades suficientemente estables.

  \end{tcolorbox}
\end{frame}

\begin{frame}{Resolución nominal no garantiza separación efectiva}
  \protect\phantomsection\label{resoluciuxf3n-nominal-no-garantiza-separaciuxf3n-efectiva}
  Dos tonos separados por \(\Delta f\) pueden seguir sin distinguirse si:

  \begin{itemize}
    \tightlist
    \item
          la ventana ensancha el lóbulo principal;
    \item
          una componente es mucho mayor que la otra;
    \item
          el ruido oculta la componente débil;
    \item
          la frecuencia cambia durante el registro.
  \end{itemize}

  \begin{tcolorbox}[enhanced jigsaw, arc=.35mm, bottomrule=.15mm, bottomtitle=1mm, breakable, colback=white, colbacktitle=quarto-callout-warning-color!10!white, colframe=quarto-callout-warning-color-frame, coltitle=black, left=2mm, leftrule=.75mm, opacityback=0, opacitybacktitle=0.6, rightrule=.15mm, title=\textcolor{quarto-callout-warning-color}{\faExclamationTriangle}\hspace{0.5em}{Advertencia}, titlerule=0mm, toprule=.15mm, toptitle=1mm]

    El ancho del lóbulo principal y el nivel de lóbulos laterales importan
    tanto como el espaciado entre bins.

  \end{tcolorbox}
\end{frame}

\begin{frame}{Una frecuencia alineada con un bin produce un caso
    especial}
  \protect\phantomsection\label{una-frecuencia-alineada-con-un-bin-produce-un-caso-especial}
  Si una cosenoide completa un número entero de ciclos en \(T\):

  \[f_0=m\Delta f,\]

  su energía coincide con bins de la DFT bajo una ventana rectangular.

  \begin{tcolorbox}[enhanced jigsaw, arc=.35mm, bottomrule=.15mm, bottomtitle=1mm, breakable, colback=white, colbacktitle=quarto-callout-note-color!10!white, colframe=quarto-callout-note-color-frame, coltitle=black, left=2mm, leftrule=.75mm, opacityback=0, opacitybacktitle=0.6, rightrule=.15mm, title=\textcolor{quarto-callout-note-color}{\faInfo}\hspace{0.5em}{Nota}, titlerule=0mm, toprule=.15mm, toptitle=1mm]

    Este caso facilita pruebas de escalamiento, pero rara vez describe
    exactamente una bioseñal real.

  \end{tcolorbox}
\end{frame}

\begin{frame}{La fuga distribuye energía entre bins}
  \protect\phantomsection\label{la-fuga-distribuye-energuxeda-entre-bins}
  Una componente no alineada con los bins produce valores en múltiples
  frecuencias.

  La causa matemática es el producto temporal:

  \[x_w[n]=x[n]w[n],\]

  que corresponde a convolución espectral.

  \begin{tcolorbox}[enhanced jigsaw, arc=.35mm, bottomrule=.15mm, bottomtitle=1mm, breakable, colback=white, colbacktitle=quarto-callout-important-color!10!white, colframe=quarto-callout-important-color-frame, coltitle=black, left=2mm, leftrule=.75mm, opacityback=0, opacitybacktitle=0.6, rightrule=.15mm, title=\textcolor{quarto-callout-important-color}{\faExclamation}\hspace{0.5em}{Importante}, titlerule=0mm, toprule=.15mm, toptitle=1mm]

    La fuga no implica que la señal contenga todas esas frecuencias como
    oscilaciones independientes.

  \end{tcolorbox}
\end{frame}

\begin{frame}{La ventana rectangular conserva amplitud y genera lóbulos
    altos}
  \protect\phantomsection\label{la-ventana-rectangular-conserva-amplitud-y-genera-luxf3bulos-altos}
  La ventana rectangular equivale a recortar el registro sin suavizar
  extremos.

  \begin{itemize}
    \tightlist
    \item
          lóbulo principal estrecho;
    \item
          lóbulos laterales relativamente altos;
    \item
          buena separación para tonos comparables y alineados;
    \item
          mayor contaminación alrededor de componentes intensas.
  \end{itemize}

  \begin{tcolorbox}[enhanced jigsaw, arc=.35mm, bottomrule=.15mm, bottomtitle=1mm, breakable, colback=white, colbacktitle=quarto-callout-warning-color!10!white, colframe=quarto-callout-warning-color-frame, coltitle=black, left=2mm, leftrule=.75mm, opacityback=0, opacitybacktitle=0.6, rightrule=.15mm, title=\textcolor{quarto-callout-warning-color}{\faExclamationTriangle}\hspace{0.5em}{Advertencia}, titlerule=0mm, toprule=.15mm, toptitle=1mm]

    Usar ``sin ventana'' significa usar una ventana rectangular.

  \end{tcolorbox}
\end{frame}

\begin{frame}{Hann reduce lóbulos laterales y ensancha el principal}
  \protect\phantomsection\label{hann-reduce-luxf3bulos-laterales-y-ensancha-el-principal}
  \[w[n]=\frac{1}{2}\left(1-\cos\frac{2\pi n}{N-1}\right).\]

  La ventana lleva los extremos hacia cero y reduce discontinuidades en la
  extensión periódica.

  \begin{tcolorbox}[enhanced jigsaw, arc=.35mm, bottomrule=.15mm, bottomtitle=1mm, breakable, colback=white, colbacktitle=quarto-callout-important-color!10!white, colframe=quarto-callout-important-color-frame, coltitle=black, left=2mm, leftrule=.75mm, opacityback=0, opacitybacktitle=0.6, rightrule=.15mm, title=\textcolor{quarto-callout-important-color}{\faExclamation}\hspace{0.5em}{Importante}, titlerule=0mm, toprule=.15mm, toptitle=1mm]

    La reducción de fuga se paga con menor capacidad para separar
    componentes cercanas.

  \end{tcolorbox}
\end{frame}

\begin{frame}{Hamming y Blackman ofrecen otros compromisos}
  \protect\phantomsection\label{hamming-y-blackman-ofrecen-otros-compromisos}
  {\def\LTcaptype{none} % do not increment counter
    \begin{longtable}[]{@{}
      >{\raggedright\arraybackslash}p{(\linewidth - 4\tabcolsep) * \real{0.3333}}
      >{\raggedright\arraybackslash}p{(\linewidth - 4\tabcolsep) * \real{0.3333}}
      >{\raggedright\arraybackslash}p{(\linewidth - 4\tabcolsep) * \real{0.3333}}@{}}
      \toprule\noalign{}
      \begin{minipage}[b]{\linewidth}\raggedright
        Ventana
      \end{minipage} & \begin{minipage}[b]{\linewidth}\raggedright
                         Uso orientativo
                       \end{minipage} & \begin{minipage}[b]{\linewidth}\raggedright
                                          Compromiso
                                        \end{minipage}                                        \\
      \midrule\noalign{}
      \endhead
      Rectangular                                 & tonos alineados y cercanos                  & lóbulos laterales altos \\
      Hann                                        & análisis general y Welch                    & balance habitual        \\
      Hamming                                     & atenuación del primer lateral               & extremos no nulos       \\
      Blackman                                    & componente débil junto a fuerte             & lóbulo principal más
      ancho                                                                                                               \\
      \bottomrule\noalign{}
    \end{longtable}
  }

  \begin{tcolorbox}[enhanced jigsaw, arc=.35mm, bottomrule=.15mm, bottomtitle=1mm, breakable, colback=white, colbacktitle=quarto-callout-warning-color!10!white, colframe=quarto-callout-warning-color-frame, coltitle=black, left=2mm, leftrule=.75mm, opacityback=0, opacitybacktitle=0.6, rightrule=.15mm, title=\textcolor{quarto-callout-warning-color}{\faExclamationTriangle}\hspace{0.5em}{Advertencia}, titlerule=0mm, toprule=.15mm, toptitle=1mm]

    La selección no debe basarse en una lista memorizada. Debe responder a
    separación, rango dinámico y objetivo de amplitud.

  \end{tcolorbox}
\end{frame}

\begin{frame}{La ganancia coherente corrige amplitud}
  \protect\phantomsection\label{la-ganancia-coherente-corrige-amplitud}
  Una ventana reduce la amplitud de un tono. Su ganancia coherente es

  \[G_c=\frac{1}{N}\sum_{n=0}^{N-1}w[n].\]

  Para estimar amplitud se corrige por \(G_c\) cuando el modelo de tono y
  alineación lo justifican.

  \begin{tcolorbox}[enhanced jigsaw, arc=.35mm, bottomrule=.15mm, bottomtitle=1mm, breakable, colback=white, colbacktitle=quarto-callout-note-color!10!white, colframe=quarto-callout-note-color-frame, coltitle=black, left=2mm, leftrule=.75mm, opacityback=0, opacitybacktitle=0.6, rightrule=.15mm, title=\textcolor{quarto-callout-note-color}{\faInfo}\hspace{0.5em}{Nota}, titlerule=0mm, toprule=.15mm, toptitle=1mm]

    La corrección de amplitud no elimina fuga ni convierte una señal no
    estacionaria en estacionaria.

  \end{tcolorbox}
\end{frame}

\begin{frame}{El ancho de banda equivalente afecta el ruido}
  \protect\phantomsection\label{el-ancho-de-banda-equivalente-afecta-el-ruido}
  Una forma discreta normalizada en bins es

  \[\mathrm{ENBW}=N\frac{\sum_n w^2[n]}{\left(\sum_n w[n]\right)^2}.\]

  Una ENBW mayor integra más potencia de ruido por bin espectral.

  \begin{tcolorbox}[enhanced jigsaw, arc=.35mm, bottomrule=.15mm, bottomtitle=1mm, breakable, colback=white, colbacktitle=quarto-callout-important-color!10!white, colframe=quarto-callout-important-color-frame, coltitle=black, left=2mm, leftrule=.75mm, opacityback=0, opacitybacktitle=0.6, rightrule=.15mm, title=\textcolor{quarto-callout-important-color}{\faExclamation}\hspace{0.5em}{Importante}, titlerule=0mm, toprule=.15mm, toptitle=1mm]

    La ventana afecta tanto la forma de una componente tonal como el piso de
    ruido estimado.

  \end{tcolorbox}
\end{frame}

\begin{frame}{Zero padding densifica el eje}
  \protect\phantomsection\label{zero-padding-densifica-el-eje}
  Agregar ceros hasta \(N_{FFT}>N\) produce separación gráfica

  \[\Delta f_{grid}=\frac{f_s}{N_{FFT}},\]

  pero la resolución física sigue determinada principalmente por
  \(T=N/f_s\) y la ventana.

  \begin{tcolorbox}[enhanced jigsaw, arc=.35mm, bottomrule=.15mm, bottomtitle=1mm, breakable, colback=white, colbacktitle=quarto-callout-warning-color!10!white, colframe=quarto-callout-warning-color-frame, coltitle=black, left=2mm, leftrule=.75mm, opacityback=0, opacitybacktitle=0.6, rightrule=.15mm, title=\textcolor{quarto-callout-warning-color}{\faExclamationTriangle}\hspace{0.5em}{Advertencia}, titlerule=0mm, toprule=.15mm, toptitle=1mm]

    Una curva más suave no contiene más observación ni permite separar
    componentes que el registro original no resolvía.

  \end{tcolorbox}
\end{frame}

\begin{frame}{Duraciones largas también tienen costos}
  \protect\phantomsection\label{duraciones-largas-tambiuxe9n-tienen-costos}
  \begin{itemize}
    \tightlist
    \item
          pueden mezclar estados fisiológicos;
    \item
          reducen localización temporal;
    \item
          aumentan probabilidad de artefactos;
    \item
          exigen mayor estabilidad del reloj y del protocolo.
  \end{itemize}

  \begin{tcolorbox}[enhanced jigsaw, arc=.35mm, bottomrule=.15mm, bottomtitle=1mm, breakable, colback=white, colbacktitle=quarto-callout-important-color!10!white, colframe=quarto-callout-important-color-frame, coltitle=black, left=2mm, leftrule=.75mm, opacityback=0, opacitybacktitle=0.6, rightrule=.15mm, title=\textcolor{quarto-callout-important-color}{\faExclamation}\hspace{0.5em}{Importante}, titlerule=0mm, toprule=.15mm, toptitle=1mm]

    La duración óptima surge del equilibrio entre resolución frecuencial y
    estacionariedad local.

  \end{tcolorbox}
\end{frame}

\begin{frame}{Ejemplo biomédico: ritmo respiratorio}
  \protect\phantomsection\label{ejemplo-biomuxe9dico-ritmo-respiratorio}
  Para distinguir 0.20 Hz de 0.25 Hz se necesita una observación
  suficientemente larga.

  \begin{itemize}
    \tightlist
    \item
          con \(T=10\) s, \(\Delta f=0.1\) Hz;
    \item
          con \(T=40\) s, \(\Delta f=0.025\) Hz.
  \end{itemize}

  \begin{tcolorbox}[enhanced jigsaw, arc=.35mm, bottomrule=.15mm, bottomtitle=1mm, breakable, colback=white, colbacktitle=quarto-callout-warning-color!10!white, colframe=quarto-callout-warning-color-frame, coltitle=black, left=2mm, leftrule=.75mm, opacityback=0, opacitybacktitle=0.6, rightrule=.15mm, title=\textcolor{quarto-callout-warning-color}{\faExclamationTriangle}\hspace{0.5em}{Advertencia}, titlerule=0mm, toprule=.15mm, toptitle=1mm]

    La mejor separación nominal de 40 s solo es útil si el ritmo no cambia
    de manera importante durante esos 40 s.

  \end{tcolorbox}
\end{frame}

\begin{frame}[fragile]{Código: comparar ventanas}
  \protect\phantomsection\label{cuxf3digo-comparar-ventanas}
  \begin{Shaded}
    \begin{Highlighting}[]
      \ImportTok{import}\NormalTok{ numpy }\ImportTok{as}\NormalTok{ np}
      \ImportTok{from}\NormalTok{ scipy.signal }\ImportTok{import}\NormalTok{ get\_window}

      \NormalTok{N }\OperatorTok{=} \BuiltInTok{len}\NormalTok{(x)}
      \ControlFlowTok{for}\NormalTok{ nombre }\KeywordTok{in}\NormalTok{ [}\StringTok{"boxcar"}\NormalTok{, }\StringTok{"hann"}\NormalTok{, }\StringTok{"hamming"}\NormalTok{, }\StringTok{"blackman"}\NormalTok{]:}
      \NormalTok{    w }\OperatorTok{=}\NormalTok{ get\_window(nombre, N, fftbins}\OperatorTok{=}\VariableTok{False}\NormalTok{)}
      \NormalTok{    Xw }\OperatorTok{=}\NormalTok{ np.fft.rfft((x }\OperatorTok{{-}}\NormalTok{ np.mean(x)) }\OperatorTok{*}\NormalTok{ w)}
      \NormalTok{    Gc }\OperatorTok{=}\NormalTok{ np.mean(w)}
      \NormalTok{    A }\OperatorTok{=}\NormalTok{ np.}\BuiltInTok{abs}\NormalTok{(Xw) }\OperatorTok{/}\NormalTok{ (N }\OperatorTok{*}\NormalTok{ Gc)}
    \end{Highlighting}
  \end{Shaded}

  \begin{tcolorbox}[enhanced jigsaw, arc=.35mm, bottomrule=.15mm, bottomtitle=1mm, breakable, colback=white, colbacktitle=quarto-callout-note-color!10!white, colframe=quarto-callout-note-color-frame, coltitle=black, left=2mm, leftrule=.75mm, opacityback=0, opacitybacktitle=0.6, rightrule=.15mm, title=\textcolor{quarto-callout-note-color}{\faInfo}\hspace{0.5em}{Nota}, titlerule=0mm, toprule=.15mm, toptitle=1mm]

    La comparación debe mantener el mismo segmento y declarar la corrección
    aplicada.

  \end{tcolorbox}
\end{frame}

\begin{frame}{Verificación de la sesión 13}
  \protect\phantomsection\label{verificaciuxf3n-de-la-sesiuxf3n-13}
  Un EEG se analiza con \(f_s=256\) Hz y ventanas de 2 s.

  \begin{enumerate}
    \tightlist
    \item
          determine \(N\) y \(\Delta f\);
    \item
          explique si zero padding a 2048 puntos cambia la resolución real;
    \item
          compare el efecto esperado de rectangular y Hann;
    \item
          proponga una duración para separar 9.8 Hz de 10.2 Hz;
    \item
          discuta el supuesto fisiológico de esa duración.
  \end{enumerate}

  \begin{tcolorbox}[enhanced jigsaw, arc=.35mm, bottomrule=.15mm, bottomtitle=1mm, breakable, colback=white, colbacktitle=quarto-callout-tip-color!10!white, colframe=quarto-callout-tip-color-frame, coltitle=black, left=2mm, leftrule=.75mm, opacityback=0, opacitybacktitle=0.6, rightrule=.15mm, title=\textcolor{quarto-callout-tip-color}{\faLightbulb}\hspace{0.5em}{Criterio de logro}, titlerule=0mm, toprule=.15mm, toptitle=1mm]

    La respuesta debe separar espaciado de bins, ancho del lóbulo y
    estabilidad temporal.

  \end{tcolorbox}
\end{frame}

\begin{frame}{Sesión 14 · La cadena de adquisición condiciona el dato}
  \protect\phantomsection\label{sesiuxf3n-14-la-cadena-de-adquisiciuxf3n-condiciona-el-dato}
  \begin{enumerate}
    \tightlist
    \item
          El fenómeno fisiológico origina una magnitud observable.
    \item
          El sensor y el acoplamiento convierten esa magnitud.
    \item
          El acondicionamiento limita banda y adapta nivel.
    \item
          El ADC discretiza tiempo y amplitud.
    \item
          El registro conserva valores, unidades y metadatos.
  \end{enumerate}

  \begin{tcolorbox}[enhanced jigsaw, arc=.35mm, bottomrule=.15mm, bottomtitle=1mm, breakable, colback=white, colbacktitle=quarto-callout-important-color!10!white, colframe=quarto-callout-important-color-frame, coltitle=black, left=2mm, leftrule=.75mm, opacityback=0, opacitybacktitle=0.6, rightrule=.15mm, title=\textcolor{quarto-callout-important-color}{\faExclamation}\hspace{0.5em}{Importante}, titlerule=0mm, toprule=.15mm, toptitle=1mm]

    Cada etapa limita la información disponible para el procesamiento
    posterior.

  \end{tcolorbox}
\end{frame}

\begin{frame}{El sensor define sensibilidad y ancho de banda}
  \protect\phantomsection\label{el-sensor-define-sensibilidad-y-ancho-de-banda}
  Modelo local:

  \[v_s(t)=Sx(t)+b+n_s(t).\]

  \begin{itemize}
    \tightlist
    \item
          \(S\): sensibilidad con unidades físicas.
    \item
          \(b\): offset.
    \item
          \(n_s(t)\): perturbación propia del sensor.
  \end{itemize}

  \begin{tcolorbox}[enhanced jigsaw, arc=.35mm, bottomrule=.15mm, bottomtitle=1mm, breakable, colback=white, colbacktitle=quarto-callout-warning-color!10!white, colframe=quarto-callout-warning-color-frame, coltitle=black, left=2mm, leftrule=.75mm, opacityback=0, opacitybacktitle=0.6, rightrule=.15mm, title=\textcolor{quarto-callout-warning-color}{\faExclamationTriangle}\hspace{0.5em}{Advertencia}, titlerule=0mm, toprule=.15mm, toptitle=1mm]

    La relación lineal suele ser válida solo dentro de un rango. Saturación,
    histéresis y acoplamiento alteran el modelo.

  \end{tcolorbox}
\end{frame}

\begin{frame}{El acoplamiento forma parte de la medición}
  \protect\phantomsection\label{el-acoplamiento-forma-parte-de-la-mediciuxf3n}
  Ejemplos:

  \begin{itemize}
    \tightlist
    \item
          impedancia electrodo-piel en ECG y EEG;
    \item
          presión de contacto en PPG;
    \item
          orientación y fijación en acelerometría;
    \item
          posición de banda en respiración.
  \end{itemize}

  \begin{tcolorbox}[enhanced jigsaw, arc=.35mm, bottomrule=.15mm, bottomtitle=1mm, breakable, colback=white, colbacktitle=quarto-callout-important-color!10!white, colframe=quarto-callout-important-color-frame, coltitle=black, left=2mm, leftrule=.75mm, opacityback=0, opacitybacktitle=0.6, rightrule=.15mm, title=\textcolor{quarto-callout-important-color}{\faExclamation}\hspace{0.5em}{Importante}, titlerule=0mm, toprule=.15mm, toptitle=1mm]

    Un sensor correcto con acoplamiento deficiente puede producir un
    registro inválido.

  \end{tcolorbox}
\end{frame}

\begin{frame}{El acondicionamiento protege y adapta la señal}
  \protect\phantomsection\label{el-acondicionamiento-protege-y-adapta-la-seuxf1al}
  Funciones frecuentes:

  \begin{itemize}
    \tightlist
    \item
          amplificación;
    \item
          rechazo de modo común;
    \item
          limitación de banda;
    \item
          ajuste de nivel y protección;
    \item
          aislamiento cuando corresponde.
  \end{itemize}

  \begin{tcolorbox}[enhanced jigsaw, arc=.35mm, bottomrule=.15mm, bottomtitle=1mm, breakable, colback=white, colbacktitle=quarto-callout-note-color!10!white, colframe=quarto-callout-note-color-frame, coltitle=black, left=2mm, leftrule=.75mm, opacityback=0, opacitybacktitle=0.6, rightrule=.15mm, title=\textcolor{quarto-callout-note-color}{\faInfo}\hspace{0.5em}{Nota}, titlerule=0mm, toprule=.15mm, toptitle=1mm]

    La ganancia debe usar el rango del convertidor sin recortar transitorios
    fisiológicos ni artefactos plausibles.

  \end{tcolorbox}
\end{frame}

\begin{frame}{El filtro antialias actúa antes del ADC}
  \protect\phantomsection\label{el-filtro-antialias-actuxfaa-antes-del-adc}
  La etapa analógica atenúa componentes por encima de la banda que el
  sistema digital puede representar.

  \begin{tcolorbox}[enhanced jigsaw, arc=.35mm, bottomrule=.15mm, bottomtitle=1mm, breakable, colback=white, colbacktitle=quarto-callout-warning-color!10!white, colframe=quarto-callout-warning-color-frame, coltitle=black, left=2mm, leftrule=.75mm, opacityback=0, opacitybacktitle=0.6, rightrule=.15mm, title=\textcolor{quarto-callout-warning-color}{\faExclamationTriangle}\hspace{0.5em}{Advertencia}, titlerule=0mm, toprule=.15mm, toptitle=1mm]

    Un filtro digital posterior no puede distinguir una frecuencia real de
    su alias cuando ambas ya producen las mismas muestras.

  \end{tcolorbox}
\end{frame}

\begin{frame}{El ADC discretiza tiempo y amplitud}
  \protect\phantomsection\label{el-adc-discretiza-tiempo-y-amplitud}
  Dos decisiones distintas:

  \begin{itemize}
    \tightlist
    \item
          muestreo: intervalos temporales \(T_s=1/f_s\);
    \item
          cuantización: niveles de amplitud.
  \end{itemize}

  Para un rango \(V_{FS}\) y \(B\) bits:

  \[\Delta_q\approx\frac{V_{FS}}{2^B}.\]

  \begin{tcolorbox}[enhanced jigsaw, arc=.35mm, bottomrule=.15mm, bottomtitle=1mm, breakable, colback=white, colbacktitle=quarto-callout-important-color!10!white, colframe=quarto-callout-important-color-frame, coltitle=black, left=2mm, leftrule=.75mm, opacityback=0, opacitybacktitle=0.6, rightrule=.15mm, title=\textcolor{quarto-callout-important-color}{\faExclamation}\hspace{0.5em}{Importante}, titlerule=0mm, toprule=.15mm, toptitle=1mm]

    Más bits reducen el escalón ideal de cuantización, pero no evitan ruido
    analógico, mala calibración ni clipping.

  \end{tcolorbox}
\end{frame}

\begin{frame}{El rango dinámico exige margen}
  \protect\phantomsection\label{el-rango-dinuxe1mico-exige-margen}
  Si la señal supera el rango del ADC, aparece saturación. Si ocupa una
  fracción muy pequeña, se desaprovecha resolución.

  \begin{tcolorbox}[enhanced jigsaw, arc=.35mm, bottomrule=.15mm, bottomtitle=1mm, breakable, colback=white, colbacktitle=quarto-callout-warning-color!10!white, colframe=quarto-callout-warning-color-frame, coltitle=black, left=2mm, leftrule=.75mm, opacityback=0, opacitybacktitle=0.6, rightrule=.15mm, title=\textcolor{quarto-callout-warning-color}{\faExclamationTriangle}\hspace{0.5em}{Advertencia}, titlerule=0mm, toprule=.15mm, toptitle=1mm]

    El clipping es una distorsión no lineal. El filtrado no reconstruye de
    manera confiable las amplitudes recortadas.

  \end{tcolorbox}
\end{frame}

\begin{frame}{El reloj determina el eje temporal}
  \protect\phantomsection\label{el-reloj-determina-el-eje-temporal}
  Un registro requiere:

  \begin{itemize}
    \tightlist
    \item
          frecuencia nominal de muestreo;
    \item
          sellos temporales o instante inicial;
    \item
          información sobre muestras perdidas;
    \item
          tolerancia o deriva del reloj;
    \item
          sincronización con otros equipos.
  \end{itemize}

  \begin{tcolorbox}[enhanced jigsaw, arc=.35mm, bottomrule=.15mm, bottomtitle=1mm, breakable, colback=white, colbacktitle=quarto-callout-note-color!10!white, colframe=quarto-callout-note-color-frame, coltitle=black, left=2mm, leftrule=.75mm, opacityback=0, opacitybacktitle=0.6, rightrule=.15mm, title=\textcolor{quarto-callout-note-color}{\faInfo}\hspace{0.5em}{Nota}, titlerule=0mm, toprule=.15mm, toptitle=1mm]

    Pequeños errores de reloj pueden alterar fase, coherencia y alineación
    de eventos en registros prolongados.

  \end{tcolorbox}
\end{frame}

\begin{frame}{La sincronización sostiene análisis multicanal}
  \protect\phantomsection\label{la-sincronizaciuxf3n-sostiene-anuxe1lisis-multicanal}
  Para comparar ECG, presión y respiración se necesita una referencia
  temporal común.

  Fuentes de desalineación:

  \begin{itemize}
    \tightlist
    \item
          retardos distintos de sensores y filtros;
    \item
          buffers de software;
    \item
          marcas manuales;
    \item
          relojes independientes.
  \end{itemize}

  \begin{tcolorbox}[enhanced jigsaw, arc=.35mm, bottomrule=.15mm, bottomtitle=1mm, breakable, colback=white, colbacktitle=quarto-callout-important-color!10!white, colframe=quarto-callout-important-color-frame, coltitle=black, left=2mm, leftrule=.75mm, opacityback=0, opacitybacktitle=0.6, rightrule=.15mm, title=\textcolor{quarto-callout-important-color}{\faExclamation}\hspace{0.5em}{Importante}, titlerule=0mm, toprule=.15mm, toptitle=1mm]

    Un desfase instrumental puede confundirse con un retardo fisiológico.

  \end{tcolorbox}
\end{frame}

\begin{frame}{Almacenar valores exige conservar contexto}
  \protect\phantomsection\label{almacenar-valores-exige-conservar-contexto}
  Metadatos mínimos:

  \begin{itemize}
    \tightlist
    \item
          unidades y escala;
    \item
          frecuencia de muestreo;
    \item
          nombres y ubicación de canales;
    \item
          configuración de ganancia y filtros;
    \item
          eventos del protocolo;
    \item
          identificador seudonimizado del participante.
  \end{itemize}

  \begin{tcolorbox}[enhanced jigsaw, arc=.35mm, bottomrule=.15mm, bottomtitle=1mm, breakable, colback=white, colbacktitle=quarto-callout-warning-color!10!white, colframe=quarto-callout-warning-color-frame, coltitle=black, left=2mm, leftrule=.75mm, opacityback=0, opacitybacktitle=0.6, rightrule=.15mm, title=\textcolor{quarto-callout-warning-color}{\faExclamationTriangle}\hspace{0.5em}{Advertencia}, titlerule=0mm, toprule=.15mm, toptitle=1mm]

    Un archivo numérico sin metadatos puede ser ilegible desde el punto de
    vista científico aunque abra correctamente.

  \end{tcolorbox}
\end{frame}

\begin{frame}{Formato tabular no siempre basta}
  \protect\phantomsection\label{formato-tabular-no-siempre-basta}
  {\def\LTcaptype{none} % do not increment counter
    \begin{longtable}[]{@{}ll@{}}
      \toprule\noalign{}
      Necesidad          & Riesgo de un CSV aislado     \\
      \midrule\noalign{}
      \endhead
      tiempo irregular   & se pierde si solo hay índice \\
      múltiples unidades & columnas ambiguas            \\
      eventos            & quedan en archivos separados \\
      calibración        & no se documenta              \\
      datos faltantes    & pueden confundirse con cero  \\
      \bottomrule\noalign{}
    \end{longtable}
  }

  \begin{tcolorbox}[enhanced jigsaw, arc=.35mm, bottomrule=.15mm, bottomtitle=1mm, breakable, colback=white, colbacktitle=quarto-callout-tip-color!10!white, colframe=quarto-callout-tip-color-frame, coltitle=black, left=2mm, leftrule=.75mm, opacityback=0, opacitybacktitle=0.6, rightrule=.15mm, title=\textcolor{quarto-callout-tip-color}{\faLightbulb}\hspace{0.5em}{Tip}, titlerule=0mm, toprule=.15mm, toptitle=1mm]

    El contrato de datos debe describir estructura, tipos, unidades, valores
    faltantes y relación entre archivos.

  \end{tcolorbox}
\end{frame}

\begin{frame}{Control de calidad antes del análisis}
  \protect\phantomsection\label{control-de-calidad-antes-del-anuxe1lisis}
  \begin{enumerate}
    \tightlist
    \item
          verificar longitud y tasa efectiva;
    \item
          comprobar monotonicidad del tiempo;
    \item
          identificar faltantes y saturación;
    \item
          revisar unidades y calibración;
    \item
          contrastar eventos con el protocolo;
    \item
          inspeccionar canales auxiliares.
  \end{enumerate}

  \begin{tcolorbox}[enhanced jigsaw, arc=.35mm, bottomrule=.15mm, bottomtitle=1mm, breakable, colback=white, colbacktitle=quarto-callout-important-color!10!white, colframe=quarto-callout-important-color-frame, coltitle=black, left=2mm, leftrule=.75mm, opacityback=0, opacitybacktitle=0.6, rightrule=.15mm, title=\textcolor{quarto-callout-important-color}{\faExclamation}\hspace{0.5em}{Importante}, titlerule=0mm, toprule=.15mm, toptitle=1mm]

    La inspección debe preceder al filtrado para evitar que el procesamiento
    oculte problemas de adquisición.

  \end{tcolorbox}
\end{frame}

\begin{frame}{Verificación de la sesión 14}
  \protect\phantomsection\label{verificaciuxf3n-de-la-sesiuxf3n-14}
  Un ECG de 12 bits usa un rango de \(\pm 5\) mV y \(f_s=500\) Hz.

  \begin{enumerate}
    \tightlist
    \item
          estime el escalón ideal de cuantización;
    \item
          identifique el máximo tiempo entre muestras;
    \item
          explique por qué necesita filtro antialias analógico;
    \item
          enumere metadatos para interpretar derivaciones;
    \item
          proponga una prueba de saturación y pérdida de muestras.
  \end{enumerate}

  \begin{tcolorbox}[enhanced jigsaw, arc=.35mm, bottomrule=.15mm, bottomtitle=1mm, breakable, colback=white, colbacktitle=quarto-callout-tip-color!10!white, colframe=quarto-callout-tip-color-frame, coltitle=black, left=2mm, leftrule=.75mm, opacityback=0, opacitybacktitle=0.6, rightrule=.15mm, title=\textcolor{quarto-callout-tip-color}{\faLightbulb}\hspace{0.5em}{Criterio de logro}, titlerule=0mm, toprule=.15mm, toptitle=1mm]

    La solución debe conectar especificaciones del sensor,
    acondicionamiento, ADC y archivo final.

  \end{tcolorbox}
\end{frame}

\begin{frame}{Cierre de la semana 7}
  \protect\phantomsection\label{cierre-de-la-semana-7}
  \begin{itemize}
    \tightlist
    \item
          La duración determina el espaciado básico de frecuencias.
    \item
          La ventana modifica resolución, fuga y potencia de ruido.
    \item
          Zero padding interpola el eje sin mejorar la información.
    \item
          La adquisición define qué contenido llega al computador.
    \item
          Metadatos y sincronización forman parte de la validez.
  \end{itemize}

  \begin{tcolorbox}[enhanced jigsaw, arc=.35mm, bottomrule=.15mm, bottomtitle=1mm, breakable, colback=white, colbacktitle=quarto-callout-note-color!10!white, colframe=quarto-callout-note-color-frame, coltitle=black, left=2mm, leftrule=.75mm, opacityback=0, opacitybacktitle=0.6, rightrule=.15mm, title=\textcolor{quarto-callout-note-color}{\faInfo}\hspace{0.5em}{Pregunta de salida}, titlerule=0mm, toprule=.15mm, toptitle=1mm]

    ¿Qué problema espectral puede parecer un fenómeno fisiológico cuando dos
    equipos no comparten el mismo reloj?

  \end{tcolorbox}
\end{frame}

\section{Semana 8 · Muestreo y estimación de
  potencia}\label{semana-8-muestreo-y-estimaciuxf3n-de-potencia}

\begin{frame}{Sesión 15 · Muestrear observa la señal en instantes
    discretos}
  \protect\phantomsection\label{sesiuxf3n-15-muestrear-observa-la-seuxf1al-en-instantes-discretos}
  Un modelo ideal usa un tren de impulsos:

  \[p(t)=\sum_{n=-\infty}^{\infty}\delta(t-nT_s).\]

  La señal muestreada es

  \[x_s(t)=x(t)p(t)=\sum_n x(nT_s)\delta(t-nT_s).\]

  \begin{tcolorbox}[enhanced jigsaw, arc=.35mm, bottomrule=.15mm, bottomtitle=1mm, breakable, colback=white, colbacktitle=quarto-callout-important-color!10!white, colframe=quarto-callout-important-color-frame, coltitle=black, left=2mm, leftrule=.75mm, opacityback=0, opacitybacktitle=0.6, rightrule=.15mm, title=\textcolor{quarto-callout-important-color}{\faExclamation}\hspace{0.5em}{Importante}, titlerule=0mm, toprule=.15mm, toptitle=1mm]

    El muestreo conserva valores en instantes separados por \(T_s\). Entre
    muestras no existe observación directa.

  \end{tcolorbox}
\end{frame}

\begin{frame}{Muestrear replica el espectro}
  \protect\phantomsection\label{muestrear-replica-el-espectro}
  El espectro del tren de impulsos también es periódico:

  \[P(f)=f_s\sum_{k=-\infty}^{\infty}\delta(f-kf_s).\]

  Por tanto:

  \[X_s(f)=f_s\sum_{k=-\infty}^{\infty}X(f-kf_s).\]

  \begin{tcolorbox}[enhanced jigsaw, arc=.35mm, bottomrule=.15mm, bottomtitle=1mm, breakable, colback=white, colbacktitle=quarto-callout-important-color!10!white, colframe=quarto-callout-important-color-frame, coltitle=black, left=2mm, leftrule=.75mm, opacityback=0, opacitybacktitle=0.6, rightrule=.15mm, title=\textcolor{quarto-callout-important-color}{\faExclamation}\hspace{0.5em}{Importante}, titlerule=0mm, toprule=.15mm, toptitle=1mm]

    El muestreo crea réplicas del espectro separadas por \(f_s\). El
    aliasing aparece cuando esas réplicas se superponen.

  \end{tcolorbox}
\end{frame}

\begin{frame}{Nyquist establece una condición ideal}
  \protect\phantomsection\label{nyquist-establece-una-condiciuxf3n-ideal}
  Si \(X(f)=0\) para \(|f|>B\), la reconstrucción ideal requiere

  \[f_s>2B.\]

  La frecuencia \(f_s/2\) se denomina frecuencia de Nyquist.

  \begin{tcolorbox}[enhanced jigsaw, arc=.35mm, bottomrule=.15mm, bottomtitle=1mm, breakable, colback=white, colbacktitle=quarto-callout-warning-color!10!white, colframe=quarto-callout-warning-color-frame, coltitle=black, left=2mm, leftrule=.75mm, opacityback=0, opacitybacktitle=0.6, rightrule=.15mm, title=\textcolor{quarto-callout-warning-color}{\faExclamationTriangle}\hspace{0.5em}{Advertencia}, titlerule=0mm, toprule=.15mm, toptitle=1mm]

    La condición supone banda estrictamente limitada y un filtro ideal. Las
    señales y filtros físicos requieren margen de transición.

  \end{tcolorbox}
\end{frame}

\begin{frame}{La frecuencia de Nyquist no es la tasa de Nyquist}
  \protect\phantomsection\label{la-frecuencia-de-nyquist-no-es-la-tasa-de-nyquist}
  \begin{itemize}
    \tightlist
    \item
          frecuencia de Nyquist: \(f_s/2\) para una tasa dada;
    \item
          tasa de Nyquist: al menos \(2B\) para una señal limitada a \(B\).
  \end{itemize}

  \begin{tcolorbox}[enhanced jigsaw, arc=.35mm, bottomrule=.15mm, bottomtitle=1mm, breakable, colback=white, colbacktitle=quarto-callout-important-color!10!white, colframe=quarto-callout-important-color-frame, coltitle=black, left=2mm, leftrule=.75mm, opacityback=0, opacitybacktitle=0.6, rightrule=.15mm, title=\textcolor{quarto-callout-important-color}{\faExclamation}\hspace{0.5em}{Importante}, titlerule=0mm, toprule=.15mm, toptitle=1mm]

    Intercambiar estos términos produce especificaciones ambiguas. Debe
    indicarse si se habla de la señal o del sistema de adquisición.

  \end{tcolorbox}
\end{frame}

\begin{frame}{El alias de una frecuencia depende de \(f_s\)}
  \protect\phantomsection\label{el-alias-de-una-frecuencia-depende-de-f_s}
  Una frecuencia \(f_0\) se pliega dentro de \([0,f_s/2]\). Una forma
  operativa es

  \[f_a=\left|\left((f_0+f_s/2)\bmod f_s\right)-f_s/2\right|.\]

  Ejemplo: con \(f_s=100\) Hz, una componente de 70 Hz aparece a 30 Hz.

  \begin{tcolorbox}[enhanced jigsaw, arc=.35mm, bottomrule=.15mm, bottomtitle=1mm, breakable, colback=white, colbacktitle=quarto-callout-warning-color!10!white, colframe=quarto-callout-warning-color-frame, coltitle=black, left=2mm, leftrule=.75mm, opacityback=0, opacitybacktitle=0.6, rightrule=.15mm, title=\textcolor{quarto-callout-warning-color}{\faExclamationTriangle}\hspace{0.5em}{Advertencia}, titlerule=0mm, toprule=.15mm, toptitle=1mm]

    El alias conserva una apariencia sinusoidal válida. El registro aislado
    no permite saber si provino de 30 Hz o de 70 Hz.

  \end{tcolorbox}
\end{frame}

\begin{frame}{El filtro antialias necesita una banda de transición}
  \protect\phantomsection\label{el-filtro-antialias-necesita-una-banda-de-transiciuxf3n}
  Especificación conceptual:

  \begin{itemize}
    \tightlist
    \item
          banda útil hasta \(f_p\);
    \item
          banda de rechazo desde \(f_a\);
    \item
          borde digital máximo \(f_s/2\);
    \item
          atenuación suficiente antes de Nyquist.
  \end{itemize}

  \begin{tcolorbox}[enhanced jigsaw, arc=.35mm, bottomrule=.15mm, bottomtitle=1mm, breakable, colback=white, colbacktitle=quarto-callout-important-color!10!white, colframe=quarto-callout-important-color-frame, coltitle=black, left=2mm, leftrule=.75mm, opacityback=0, opacitybacktitle=0.6, rightrule=.15mm, title=\textcolor{quarto-callout-important-color}{\faExclamation}\hspace{0.5em}{Importante}, titlerule=0mm, toprule=.15mm, toptitle=1mm]

    Elegir \(f_s\) implica reservar espacio para que el filtro analógico
    pase de la banda útil a la banda de rechazo.

  \end{tcolorbox}
\end{frame}

\begin{frame}{Sobremuestrear simplifica parte del diseño}
  \protect\phantomsection\label{sobremuestrear-simplifica-parte-del-diseuxf1o}
  Una tasa mayor puede:

  \begin{itemize}
    \tightlist
    \item
          ampliar la transición del filtro antialias;
    \item
          mejorar la localización temporal;
    \item
          reducir sensibilidad a ciertos errores analógicos.
  \end{itemize}

  También aumenta datos, consumo y costo de transmisión.

  \begin{tcolorbox}[enhanced jigsaw, arc=.35mm, bottomrule=.15mm, bottomtitle=1mm, breakable, colback=white, colbacktitle=quarto-callout-note-color!10!white, colframe=quarto-callout-note-color-frame, coltitle=black, left=2mm, leftrule=.75mm, opacityback=0, opacitybacktitle=0.6, rightrule=.15mm, title=\textcolor{quarto-callout-note-color}{\faInfo}\hspace{0.5em}{Nota}, titlerule=0mm, toprule=.15mm, toptitle=1mm]

    Sobremuestrear no aumenta por sí mismo la duración ni la resolución
    frecuencial de un registro fijo en segundos.

  \end{tcolorbox}
\end{frame}

\begin{frame}{La resolución temporal básica es \(T_s\)}
  \protect\phantomsection\label{la-resoluciuxf3n-temporal-buxe1sica-es-t_s}
  \[T_s=\frac{1}{f_s}.\]

  Si \(f_s=500\) Hz, la separación nominal entre muestras es 2 ms.

  \begin{tcolorbox}[enhanced jigsaw, arc=.35mm, bottomrule=.15mm, bottomtitle=1mm, breakable, colback=white, colbacktitle=quarto-callout-warning-color!10!white, colframe=quarto-callout-warning-color-frame, coltitle=black, left=2mm, leftrule=.75mm, opacityback=0, opacitybacktitle=0.6, rightrule=.15mm, title=\textcolor{quarto-callout-warning-color}{\faExclamationTriangle}\hspace{0.5em}{Advertencia}, titlerule=0mm, toprule=.15mm, toptitle=1mm]

    La precisión de un evento no queda determinada solo por \(T_s\).
    Interpolación, ruido, ancho de banda y algoritmo también influyen.

  \end{tcolorbox}
\end{frame}

\begin{frame}{Duración y muestreo controlan aspectos distintos}
  \protect\phantomsection\label{duraciuxf3n-y-muestreo-controlan-aspectos-distintos}
  {\def\LTcaptype{none} % do not increment counter
    \begin{longtable}[]{@{}ll@{}}
      \toprule\noalign{}
      Parámetro  & Relación principal                                 \\
      \midrule\noalign{}
      \endhead
      \(f_s\)    & rango frecuencial representable y paso temporal    \\
      \(T\)      & resolución frecuencial básica y contexto observado \\
      \(N=f_sT\) & tamaño del registro y costo computacional          \\
      \bottomrule\noalign{}
    \end{longtable}
  }

  \begin{tcolorbox}[enhanced jigsaw, arc=.35mm, bottomrule=.15mm, bottomtitle=1mm, breakable, colback=white, colbacktitle=quarto-callout-important-color!10!white, colframe=quarto-callout-important-color-frame, coltitle=black, left=2mm, leftrule=.75mm, opacityback=0, opacitybacktitle=0.6, rightrule=.15mm, title=\textcolor{quarto-callout-important-color}{\faExclamation}\hspace{0.5em}{Importante}, titlerule=0mm, toprule=.15mm, toptitle=1mm]

    Una mayor cantidad de muestras puede provenir de mayor duración o mayor
    tasa. Esas decisiones no son equivalentes.

  \end{tcolorbox}
\end{frame}

\begin{frame}{Muestreo irregular exige otro modelo}
  \protect\phantomsection\label{muestreo-irregular-exige-otro-modelo}
  Si los instantes son \(t_n\) y no cumplen \(t_n=nT_s\), la DFT estándar
  no representa exactamente el eje temporal real.

  Opciones:

  \begin{itemize}
    \tightlist
    \item
          corregir o interpolar con justificación;
    \item
          usar métodos para muestreo no uniforme;
    \item
          modelar explícitamente los sellos temporales.
  \end{itemize}

  \begin{tcolorbox}[enhanced jigsaw, arc=.35mm, bottomrule=.15mm, bottomtitle=1mm, breakable, colback=white, colbacktitle=quarto-callout-warning-color!10!white, colframe=quarto-callout-warning-color-frame, coltitle=black, left=2mm, leftrule=.75mm, opacityback=0, opacitybacktitle=0.6, rightrule=.15mm, title=\textcolor{quarto-callout-warning-color}{\faExclamationTriangle}\hspace{0.5em}{Advertencia}, titlerule=0mm, toprule=.15mm, toptitle=1mm]

    Asignar una frecuencia nominal a datos irregulares puede introducir
    errores de fase y frecuencia.

  \end{tcolorbox}
\end{frame}

\begin{frame}[fragile]{Código: observar aliasing}
  \protect\phantomsection\label{cuxf3digo-observar-aliasing}
  \begin{Shaded}
    \begin{Highlighting}[]
      \ImportTok{import}\NormalTok{ numpy }\ImportTok{as}\NormalTok{ np}

      \NormalTok{fs }\OperatorTok{=} \FloatTok{100.0}
      \NormalTok{t }\OperatorTok{=}\NormalTok{ np.arange(}\DecValTok{0}\NormalTok{, }\DecValTok{1}\NormalTok{, }\DecValTok{1}\OperatorTok{/}\NormalTok{fs)}
      \NormalTok{x30 }\OperatorTok{=}\NormalTok{ np.cos(}\DecValTok{2}\OperatorTok{*}\NormalTok{np.pi}\OperatorTok{*}\DecValTok{30}\OperatorTok{*}\NormalTok{t)}
      \NormalTok{x70 }\OperatorTok{=}\NormalTok{ np.cos(}\DecValTok{2}\OperatorTok{*}\NormalTok{np.pi}\OperatorTok{*}\DecValTok{70}\OperatorTok{*}\NormalTok{t)}

      \NormalTok{error }\OperatorTok{=}\NormalTok{ np.}\BuiltInTok{max}\NormalTok{(np.}\BuiltInTok{abs}\NormalTok{(x30 }\OperatorTok{{-}}\NormalTok{ x70))}
    \end{Highlighting}
  \end{Shaded}

  \begin{tcolorbox}[enhanced jigsaw, arc=.35mm, bottomrule=.15mm, bottomtitle=1mm, breakable, colback=white, colbacktitle=quarto-callout-note-color!10!white, colframe=quarto-callout-note-color-frame, coltitle=black, left=2mm, leftrule=.75mm, opacityback=0, opacitybacktitle=0.6, rightrule=.15mm, title=\textcolor{quarto-callout-note-color}{\faInfo}\hspace{0.5em}{Nota}, titlerule=0mm, toprule=.15mm, toptitle=1mm]

    Con esos instantes, ambas secuencias coinciden salvo error numérico. El
    ejemplo demuestra indistinguibilidad discreta, no igualdad de las
    señales continuas.

  \end{tcolorbox}
\end{frame}

\begin{frame}{Selección de \(f_s\) para una bioseñal}
  \protect\phantomsection\label{selecciuxf3n-de-f_s-para-una-bioseuxf1al}
  Procedimiento:

  \begin{enumerate}
    \tightlist
    \item
          definir la banda fisiológica necesaria para la tarea;
    \item
          incluir transitorios que deban conservarse;
    \item
          establecer margen para el filtro antialias;
    \item
          revisar precisión temporal y sincronización;
    \item
          evaluar almacenamiento y transmisión;
    \item
          documentar la tasa efectiva.
  \end{enumerate}

  \begin{tcolorbox}[enhanced jigsaw, arc=.35mm, bottomrule=.15mm, bottomtitle=1mm, breakable, colback=white, colbacktitle=quarto-callout-tip-color!10!white, colframe=quarto-callout-tip-color-frame, coltitle=black, left=2mm, leftrule=.75mm, opacityback=0, opacitybacktitle=0.6, rightrule=.15mm, title=\textcolor{quarto-callout-tip-color}{\faLightbulb}\hspace{0.5em}{Tip}, titlerule=0mm, toprule=.15mm, toptitle=1mm]

    La tasa se justifica desde la pregunta y la cadena completa, no desde
    una regla fija por modalidad.

  \end{tcolorbox}
\end{frame}

\begin{frame}{Verificación de la sesión 15}
  \protect\phantomsection\label{verificaciuxf3n-de-la-sesiuxf3n-15}
  Una señal contiene información hasta 120 Hz y se muestrea a 250 Hz.

  \begin{enumerate}
    \tightlist
    \item
          calcule la frecuencia de Nyquist;
    \item
          explique por qué el margen de 5 Hz resulta insuficiente para un filtro
          real;
    \item
          determine el alias de una interferencia de 180 Hz;
    \item
          calcule \(T_s\);
    \item
          proponga una tasa razonada, sin asumir un valor universal.
  \end{enumerate}

  \begin{tcolorbox}[enhanced jigsaw, arc=.35mm, bottomrule=.15mm, bottomtitle=1mm, breakable, colback=white, colbacktitle=quarto-callout-tip-color!10!white, colframe=quarto-callout-tip-color-frame, coltitle=black, left=2mm, leftrule=.75mm, opacityback=0, opacitybacktitle=0.6, rightrule=.15mm, title=\textcolor{quarto-callout-tip-color}{\faLightbulb}\hspace{0.5em}{Criterio de logro}, titlerule=0mm, toprule=.15mm, toptitle=1mm]

    La propuesta debe incluir banda de transición, atenuación y precisión
    temporal.

  \end{tcolorbox}
\end{frame}

\begin{frame}{Sesión 16 · Potencia espectral describe distribución
    cuadrática}
  \protect\phantomsection\label{sesiuxf3n-16-potencia-espectral-describe-distribuciuxf3n-cuadruxe1tica}
  Para una señal de potencia, la PSD \(S_{xx}(f)\) distribuye potencia por
  unidad de frecuencia.

  \[P=\int_{-\infty}^{\infty}S_{xx}(f)\,df.\]

  \begin{tcolorbox}[enhanced jigsaw, arc=.35mm, bottomrule=.15mm, bottomtitle=1mm, breakable, colback=white, colbacktitle=quarto-callout-important-color!10!white, colframe=quarto-callout-important-color-frame, coltitle=black, left=2mm, leftrule=.75mm, opacityback=0, opacitybacktitle=0.6, rightrule=.15mm, title=\textcolor{quarto-callout-important-color}{\faExclamation}\hspace{0.5em}{Importante}, titlerule=0mm, toprule=.15mm, toptitle=1mm]

    La PSD permite integrar potencia en bandas, comparar ritmos y
    cuantificar cómo se distribuye la variabilidad.

  \end{tcolorbox}
\end{frame}

\begin{frame}{Energía y potencia no son intercambiables}
  \protect\phantomsection\label{energuxeda-y-potencia-no-son-intercambiables}
  Energía discreta:

  \[E_x=\sum_{n=-\infty}^{\infty}|x[n]|^2.\]

  Potencia media:

  \[P_x=\lim_{N\to\infty}\frac{1}{2N+1}\sum_{n=-N}^{N}|x[n]|^2.\]

  \begin{tcolorbox}[enhanced jigsaw, arc=.35mm, bottomrule=.15mm, bottomtitle=1mm, breakable, colback=white, colbacktitle=quarto-callout-warning-color!10!white, colframe=quarto-callout-warning-color-frame, coltitle=black, left=2mm, leftrule=.75mm, opacityback=0, opacitybacktitle=0.6, rightrule=.15mm, title=\textcolor{quarto-callout-warning-color}{\faExclamationTriangle}\hspace{0.5em}{Advertencia}, titlerule=0mm, toprule=.15mm, toptitle=1mm]

    Un registro finito siempre tiene suma cuadrática finita, pero su modelo
    fisiológico puede tratarse como realización de un proceso de potencia.

  \end{tcolorbox}
\end{frame}

\begin{frame}{Espectro de energía y PSD responden a modelos distintos}
  \protect\phantomsection\label{espectro-de-energuxeda-y-psd-responden-a-modelos-distintos}
  \begin{itemize}
    \tightlist
    \item
          \(|X(f)|^2\) describe densidad de energía para señales de energía.
    \item
          \(S_{xx}(f)\) describe densidad de potencia para procesos o señales
          persistentes.
    \item
          Un registro finito permite estimar estas cantidades, no observarlas
          sin incertidumbre.
  \end{itemize}

  \begin{tcolorbox}[enhanced jigsaw, arc=.35mm, bottomrule=.15mm, bottomtitle=1mm, breakable, colback=white, colbacktitle=quarto-callout-important-color!10!white, colframe=quarto-callout-important-color-frame, coltitle=black, left=2mm, leftrule=.75mm, opacityback=0, opacitybacktitle=0.6, rightrule=.15mm, title=\textcolor{quarto-callout-important-color}{\faExclamation}\hspace{0.5em}{Importante}, titlerule=0mm, toprule=.15mm, toptitle=1mm]

    El periodograma es un estimador de PSD. No es la PSD verdadera del
    proceso.

  \end{tcolorbox}
\end{frame}

\begin{frame}{El periodograma parte de la DFT}
  \protect\phantomsection\label{el-periodograma-parte-de-la-dft}
  Una forma bilateral para \(N\) muestras y tasa \(f_s\) es

  \[\widehat S_{xx}[k]=\frac{1}{Nf_s}|X[k]|^2.\]

  La constante cambia si se usa ventana, espectro unilateral o una
  convención diferente.

  \begin{tcolorbox}[enhanced jigsaw, arc=.35mm, bottomrule=.15mm, bottomtitle=1mm, breakable, colback=white, colbacktitle=quarto-callout-warning-color!10!white, colframe=quarto-callout-warning-color-frame, coltitle=black, left=2mm, leftrule=.75mm, opacityback=0, opacitybacktitle=0.6, rightrule=.15mm, title=\textcolor{quarto-callout-warning-color}{\faExclamationTriangle}\hspace{0.5em}{Advertencia}, titlerule=0mm, toprule=.15mm, toptitle=1mm]

    No se puede comparar PSD entre programas si no se verifican
    escalamiento, ventana, detrend y lados del espectro.

  \end{tcolorbox}
\end{frame}

\begin{frame}{Las unidades distinguen espectro y densidad}
  \protect\phantomsection\label{las-unidades-distinguen-espectro-y-densidad}
  Si \(x\) se mide en mV:

  \begin{itemize}
    \tightlist
    \item
          amplitud espectral: mV;
    \item
          potencia por bin: mV\(^2\);
    \item
          PSD: mV\(^2\)/Hz.
  \end{itemize}

  \begin{tcolorbox}[enhanced jigsaw, arc=.35mm, bottomrule=.15mm, bottomtitle=1mm, breakable, colback=white, colbacktitle=quarto-callout-important-color!10!white, colframe=quarto-callout-important-color-frame, coltitle=black, left=2mm, leftrule=.75mm, opacityback=0, opacitybacktitle=0.6, rightrule=.15mm, title=\textcolor{quarto-callout-important-color}{\faExclamation}\hspace{0.5em}{Importante}, titlerule=0mm, toprule=.15mm, toptitle=1mm]

    La integración de PSD sobre Hz debe recuperar una potencia en unidades
    cuadráticas de la señal.

  \end{tcolorbox}
\end{frame}

\begin{frame}{La PSD unilateral conserva potencia}
  \protect\phantomsection\label{la-psd-unilateral-conserva-potencia}
  Para señales reales se combinan contribuciones positivas y negativas.

  \begin{itemize}
    \tightlist
    \item
          DC no se duplica.
    \item
          Nyquist no se duplica cuando existe.
    \item
          Los bins interiores se duplican.
  \end{itemize}

  \begin{tcolorbox}[enhanced jigsaw, arc=.35mm, bottomrule=.15mm, bottomtitle=1mm, breakable, colback=white, colbacktitle=quarto-callout-warning-color!10!white, colframe=quarto-callout-warning-color-frame, coltitle=black, left=2mm, leftrule=.75mm, opacityback=0, opacitybacktitle=0.6, rightrule=.15mm, title=\textcolor{quarto-callout-warning-color}{\faExclamationTriangle}\hspace{0.5em}{Advertencia}, titlerule=0mm, toprule=.15mm, toptitle=1mm]

    La regla de duplicación depende de si se trabaja con amplitud, potencia
    o densidad. Debe aplicarse en la escala correspondiente.

  \end{tcolorbox}
\end{frame}

\begin{frame}{Parseval ofrece una prueba de consistencia}
  \protect\phantomsection\label{parseval-ofrece-una-prueba-de-consistencia}
  Para la DFT adoptada:

  \[\sum_{n=0}^{N-1}|x[n]|^2=\frac{1}{N}\sum_{k=0}^{N-1}|X[k]|^2.\]

  \begin{tcolorbox}[enhanced jigsaw, arc=.35mm, bottomrule=.15mm, bottomtitle=1mm, breakable, colback=white, colbacktitle=quarto-callout-tip-color!10!white, colframe=quarto-callout-tip-color-frame, coltitle=black, left=2mm, leftrule=.75mm, opacityback=0, opacitybacktitle=0.6, rightrule=.15mm, title=\textcolor{quarto-callout-tip-color}{\faLightbulb}\hspace{0.5em}{Tip}, titlerule=0mm, toprule=.15mm, toptitle=1mm]

    Comparar potencia temporal con la integral numérica de PSD permite
    detectar errores de escalamiento.

  \end{tcolorbox}
\end{frame}

\begin{frame}{Retirar la media cambia la PSD cerca de 0 Hz}
  \protect\phantomsection\label{retirar-la-media-cambia-la-psd-cerca-de-0-hz}
  La componente DC representa el promedio del segmento.

  \[x_0[n]=x[n]-\bar{x}.\]

  \begin{tcolorbox}[enhanced jigsaw, arc=.35mm, bottomrule=.15mm, bottomtitle=1mm, breakable, colback=white, colbacktitle=quarto-callout-warning-color!10!white, colframe=quarto-callout-warning-color-frame, coltitle=black, left=2mm, leftrule=.75mm, opacityback=0, opacitybacktitle=0.6, rightrule=.15mm, title=\textcolor{quarto-callout-warning-color}{\faExclamationTriangle}\hspace{0.5em}{Advertencia}, titlerule=0mm, toprule=.15mm, toptitle=1mm]

    Retirar media o tendencia debe declararse. Puede ser conveniente para
    estudiar oscilaciones, pero elimina información de nivel.

  \end{tcolorbox}
\end{frame}

\begin{frame}{El periodograma tiene alta variabilidad}
  \protect\phantomsection\label{el-periodograma-tiene-alta-variabilidad}
  Aumentar \(N\) mejora el espaciado frecuencial, pero el periodograma
  simple no se vuelve necesariamente estable como estimador puntual.

  \begin{tcolorbox}[enhanced jigsaw, arc=.35mm, bottomrule=.15mm, bottomtitle=1mm, breakable, colback=white, colbacktitle=quarto-callout-important-color!10!white, colframe=quarto-callout-important-color-frame, coltitle=black, left=2mm, leftrule=.75mm, opacityback=0, opacitybacktitle=0.6, rightrule=.15mm, title=\textcolor{quarto-callout-important-color}{\faExclamation}\hspace{0.5em}{Importante}, titlerule=0mm, toprule=.15mm, toptitle=1mm]

    La varianza del estimador motiva el promediado de segmentos y el método
    de Welch.

  \end{tcolorbox}
\end{frame}

\begin{frame}{Potencia en una banda}
  \protect\phantomsection\label{potencia-en-una-banda}
  Para una banda \([f_1,f_2]\):

  \[P_{banda}=\int_{f_1}^{f_2}S_{xx}(f)\,df.\]

  En un eje discreto uniforme se aproxima mediante suma o integración
  trapezoidal.

  \begin{tcolorbox}[enhanced jigsaw, arc=.35mm, bottomrule=.15mm, bottomtitle=1mm, breakable, colback=white, colbacktitle=quarto-callout-warning-color!10!white, colframe=quarto-callout-warning-color-frame, coltitle=black, left=2mm, leftrule=.75mm, opacityback=0, opacitybacktitle=0.6, rightrule=.15mm, title=\textcolor{quarto-callout-warning-color}{\faExclamationTriangle}\hspace{0.5em}{Advertencia}, titlerule=0mm, toprule=.15mm, toptitle=1mm]

    Las bandas deben justificarse para la modalidad, población, estado y
    método de adquisición.

  \end{tcolorbox}
\end{frame}

\begin{frame}{Potencia absoluta y relativa}
  \protect\phantomsection\label{potencia-absoluta-y-relativa}
  \[P_{rel}=\frac{P_{banda}}{P_{referencia}}.\]

  La potencia relativa reduce algunas diferencias de escala, pero acopla
  el numerador al denominador.

  \begin{tcolorbox}[enhanced jigsaw, arc=.35mm, bottomrule=.15mm, bottomtitle=1mm, breakable, colback=white, colbacktitle=quarto-callout-note-color!10!white, colframe=quarto-callout-note-color-frame, coltitle=black, left=2mm, leftrule=.75mm, opacityback=0, opacitybacktitle=0.6, rightrule=.15mm, title=\textcolor{quarto-callout-note-color}{\faInfo}\hspace{0.5em}{Nota}, titlerule=0mm, toprule=.15mm, toptitle=1mm]

    Un aumento de potencia relativa puede deberse a aumento de la banda o
    disminución de la potencia de referencia.

  \end{tcolorbox}
\end{frame}

\begin{frame}[fragile]{Código: periodograma con unidades explícitas}
  \protect\phantomsection\label{cuxf3digo-periodograma-con-unidades-expluxedcitas}
  \begin{Shaded}
    \begin{Highlighting}[]
      \ImportTok{from}\NormalTok{ scipy.signal }\ImportTok{import}\NormalTok{ periodogram}

      \NormalTok{f, Pxx }\OperatorTok{=}\NormalTok{ periodogram(}
      \NormalTok{    x,}
      \NormalTok{    fs}\OperatorTok{=}\NormalTok{fs,}
      \NormalTok{    window}\OperatorTok{=}\StringTok{"hann"}\NormalTok{,}
      \NormalTok{    detrend}\OperatorTok{=}\StringTok{"constant"}\NormalTok{,}
      \NormalTok{    scaling}\OperatorTok{=}\StringTok{"density"}\NormalTok{,}
      \NormalTok{    return\_onesided}\OperatorTok{=}\VariableTok{True}\NormalTok{,}
      \NormalTok{)}
      \NormalTok{potencia }\OperatorTok{=}\NormalTok{ np.trapz(Pxx, f)}
    \end{Highlighting}
  \end{Shaded}

  \begin{tcolorbox}[enhanced jigsaw, arc=.35mm, bottomrule=.15mm, bottomtitle=1mm, breakable, colback=white, colbacktitle=quarto-callout-note-color!10!white, colframe=quarto-callout-note-color-frame, coltitle=black, left=2mm, leftrule=.75mm, opacityback=0, opacitybacktitle=0.6, rightrule=.15mm, title=\textcolor{quarto-callout-note-color}{\faInfo}\hspace{0.5em}{Nota}, titlerule=0mm, toprule=.15mm, toptitle=1mm]

    Con \texttt{scaling="density"}, \texttt{Pxx} tiene unidades de la señal
    al cuadrado por Hz.

  \end{tcolorbox}
\end{frame}

\begin{frame}{Lectura biomédica prudente}
  \protect\phantomsection\label{lectura-biomuxe9dica-prudente}
  Antes de interpretar un pico o una banda:

  \begin{itemize}
    \tightlist
    \item
          verificar estado fisiológico y protocolo;
    \item
          revisar artefactos y canales auxiliares;
    \item
          confirmar duración y resolución;
    \item
          comprobar preprocesamiento;
    \item
          estimar incertidumbre o variabilidad.
  \end{itemize}

  \begin{tcolorbox}[enhanced jigsaw, arc=.35mm, bottomrule=.15mm, bottomtitle=1mm, breakable, colback=white, colbacktitle=quarto-callout-important-color!10!white, colframe=quarto-callout-important-color-frame, coltitle=black, left=2mm, leftrule=.75mm, opacityback=0, opacitybacktitle=0.6, rightrule=.15mm, title=\textcolor{quarto-callout-important-color}{\faExclamation}\hspace{0.5em}{Importante}, titlerule=0mm, toprule=.15mm, toptitle=1mm]

    Una asociación entre banda y fenómeno requiere evidencia adicional. La
    frecuencia por sí sola no identifica la fuente.

  \end{tcolorbox}
\end{frame}

\begin{frame}[fragile]{Verificación de la sesión 16}
  \protect\phantomsection\label{verificaciuxf3n-de-la-sesiuxf3n-16}
  Una PSD de aceleración se expresa en \((m/s^2)^2/Hz\).

  \begin{enumerate}
    \tightlist
    \item
          indique las unidades de su integral entre 0 y 20 Hz;
    \item
          explique qué hace \texttt{detrend="constant"};
    \item
          compare periodograma unilateral y bilateral;
    \item
          proponga una prueba de Parseval;
    \item
          explique por qué un pico no identifica automáticamente un mecanismo.
  \end{enumerate}

  \begin{tcolorbox}[enhanced jigsaw, arc=.35mm, bottomrule=.15mm, bottomtitle=1mm, breakable, colback=white, colbacktitle=quarto-callout-tip-color!10!white, colframe=quarto-callout-tip-color-frame, coltitle=black, left=2mm, leftrule=.75mm, opacityback=0, opacitybacktitle=0.6, rightrule=.15mm, title=\textcolor{quarto-callout-tip-color}{\faLightbulb}\hspace{0.5em}{Criterio de logro}, titlerule=0mm, toprule=.15mm, toptitle=1mm]

    La respuesta debe conservar unidades y diferenciar estimador, parámetro
    y explicación fisiológica.

  \end{tcolorbox}
\end{frame}

\begin{frame}{Cierre de la semana 8}
  \protect\phantomsection\label{cierre-de-la-semana-8}
  \begin{itemize}
    \tightlist
    \item
          El muestreo replica el espectro y puede producir aliasing.
    \item
          Nyquist ofrece una condición ideal, no un margen de diseño suficiente.
    \item
          \(f_s\) y \(T\) controlan propiedades distintas.
    \item
          El periodograma estima densidad de potencia a partir de un bloque
          finito.
    \item
          Las unidades y Parseval permiten verificar el escalamiento.
  \end{itemize}

  \begin{tcolorbox}[enhanced jigsaw, arc=.35mm, bottomrule=.15mm, bottomtitle=1mm, breakable, colback=white, colbacktitle=quarto-callout-note-color!10!white, colframe=quarto-callout-note-color-frame, coltitle=black, left=2mm, leftrule=.75mm, opacityback=0, opacitybacktitle=0.6, rightrule=.15mm, title=\textcolor{quarto-callout-note-color}{\faInfo}\hspace{0.5em}{Pregunta de salida}, titlerule=0mm, toprule=.15mm, toptitle=1mm]

    ¿Por qué duplicar la frecuencia de muestreo sin cambiar la duración no
    divide a la mitad la incertidumbre del periodograma?

  \end{tcolorbox}
\end{frame}

\section{Semana 9 · Welch y análisis
  tiempo-frecuencia}\label{semana-9-welch-y-anuxe1lisis-tiempo-frecuencia}

\begin{frame}{Sesión 17 · Welch reduce variabilidad mediante promediado}
  \protect\phantomsection\label{sesiuxf3n-17-welch-reduce-variabilidad-mediante-promediado}
  El método divide el registro en segmentos, aplica una ventana y promedia
  periodogramas modificados.

  \[\widehat S_{Welch}(f)=\frac{1}{K}\sum_{r=1}^{K}\widehat S_r(f).\]

  \begin{tcolorbox}[enhanced jigsaw, arc=.35mm, bottomrule=.15mm, bottomtitle=1mm, breakable, colback=white, colbacktitle=quarto-callout-important-color!10!white, colframe=quarto-callout-important-color-frame, coltitle=black, left=2mm, leftrule=.75mm, opacityback=0, opacitybacktitle=0.6, rightrule=.15mm, title=\textcolor{quarto-callout-important-color}{\faExclamation}\hspace{0.5em}{Importante}, titlerule=0mm, toprule=.15mm, toptitle=1mm]

    El promediado reduce variabilidad del estimador a cambio de resolución
    frecuencial y dependencia entre segmentos solapados.

  \end{tcolorbox}
\end{frame}

\begin{frame}{El procedimiento contiene cuatro decisiones}
  \protect\phantomsection\label{el-procedimiento-contiene-cuatro-decisiones}
  \begin{enumerate}
    \tightlist
    \item
          longitud de segmento \(L\);
    \item
          solapamiento \(O\);
    \item
          ventana \(w[n]\);
    \item
          número de FFT \(N_{FFT}\).
  \end{enumerate}

  \begin{tcolorbox}[enhanced jigsaw, arc=.35mm, bottomrule=.15mm, bottomtitle=1mm, breakable, colback=white, colbacktitle=quarto-callout-warning-color!10!white, colframe=quarto-callout-warning-color-frame, coltitle=black, left=2mm, leftrule=.75mm, opacityback=0, opacitybacktitle=0.6, rightrule=.15mm, title=\textcolor{quarto-callout-warning-color}{\faExclamationTriangle}\hspace{0.5em}{Advertencia}, titlerule=0mm, toprule=.15mm, toptitle=1mm]

    Dos PSD calculadas con parámetros distintos pueden diferir aunque
    provengan del mismo registro.

  \end{tcolorbox}
\end{frame}

\begin{frame}{Segmentos cortos producen más promedios}
  \protect\phantomsection\label{segmentos-cortos-producen-muxe1s-promedios}
  Al reducir \(L\):

  \begin{itemize}
    \tightlist
    \item
          aumenta el número de segmentos;
    \item
          suele disminuir la variabilidad;
    \item
          empeora la resolución frecuencial;
    \item
          mejora la aproximación de estacionariedad local.
  \end{itemize}

  \begin{tcolorbox}[enhanced jigsaw, arc=.35mm, bottomrule=.15mm, bottomtitle=1mm, breakable, colback=white, colbacktitle=quarto-callout-important-color!10!white, colframe=quarto-callout-important-color-frame, coltitle=black, left=2mm, leftrule=.75mm, opacityback=0, opacitybacktitle=0.6, rightrule=.15mm, title=\textcolor{quarto-callout-important-color}{\faExclamation}\hspace{0.5em}{Importante}, titlerule=0mm, toprule=.15mm, toptitle=1mm]

    La longitud debe relacionarse con la separación frecuencial necesaria y
    la escala temporal del fenómeno.

  \end{tcolorbox}
\end{frame}

\begin{frame}{El solapamiento reutiliza muestras}
  \protect\phantomsection\label{el-solapamiento-reutiliza-muestras}
  Con avance \(H=L-O\):

  \[\text{solapamiento relativo}=\frac{O}{L}=1-\frac{H}{L}.\]

  \begin{tcolorbox}[enhanced jigsaw, arc=.35mm, bottomrule=.15mm, bottomtitle=1mm, breakable, colback=white, colbacktitle=quarto-callout-warning-color!10!white, colframe=quarto-callout-warning-color-frame, coltitle=black, left=2mm, leftrule=.75mm, opacityback=0, opacitybacktitle=0.6, rightrule=.15mm, title=\textcolor{quarto-callout-warning-color}{\faExclamationTriangle}\hspace{0.5em}{Advertencia}, titlerule=0mm, toprule=.15mm, toptitle=1mm]

    Más solapamiento no crea observaciones independientes. Su beneficio
    depende de la ventana y del costo computacional.

  \end{tcolorbox}
\end{frame}

\begin{frame}{Hann y 50 \% forman una elección frecuente}
  \protect\phantomsection\label{hann-y-50-forman-una-elecciuxf3n-frecuente}
  La combinación suele ofrecer cobertura uniforme razonable y eficiencia
  computacional.

  \begin{tcolorbox}[enhanced jigsaw, arc=.35mm, bottomrule=.15mm, bottomtitle=1mm, breakable, colback=white, colbacktitle=quarto-callout-note-color!10!white, colframe=quarto-callout-note-color-frame, coltitle=black, left=2mm, leftrule=.75mm, opacityback=0, opacitybacktitle=0.6, rightrule=.15mm, title=\textcolor{quarto-callout-note-color}{\faInfo}\hspace{0.5em}{Nota}, titlerule=0mm, toprule=.15mm, toptitle=1mm]

    ``Frecuente'' no significa universal. Eventos transitorios, registros
    breves o necesidades de amplitud pueden exigir otra selección.

  \end{tcolorbox}
\end{frame}

\begin{frame}{El escalamiento incluye la energía de la ventana}
  \protect\phantomsection\label{el-escalamiento-incluye-la-energuxeda-de-la-ventana}
  El periodograma modificado compensa, según la convención, la suma
  cuadrática de \(w[n]\):

  \[U=\frac{1}{L}\sum_{n=0}^{L-1}w^2[n].\]

  \begin{tcolorbox}[enhanced jigsaw, arc=.35mm, bottomrule=.15mm, bottomtitle=1mm, breakable, colback=white, colbacktitle=quarto-callout-important-color!10!white, colframe=quarto-callout-important-color-frame, coltitle=black, left=2mm, leftrule=.75mm, opacityback=0, opacitybacktitle=0.6, rightrule=.15mm, title=\textcolor{quarto-callout-important-color}{\faExclamation}\hspace{0.5em}{Importante}, titlerule=0mm, toprule=.15mm, toptitle=1mm]

    La normalización correcta permite que la integral de la PSD represente
    potencia del segmento.

  \end{tcolorbox}
\end{frame}

\begin{frame}{\(N_{FFT}\) cambia la malla, no la evidencia}
  \protect\phantomsection\label{n_fft-cambia-la-malla-no-la-evidencia}
  Si \(N_{FFT}>L\), se aplica zero padding dentro de cada segmento.

  \begin{tcolorbox}[enhanced jigsaw, arc=.35mm, bottomrule=.15mm, bottomtitle=1mm, breakable, colback=white, colbacktitle=quarto-callout-warning-color!10!white, colframe=quarto-callout-warning-color-frame, coltitle=black, left=2mm, leftrule=.75mm, opacityback=0, opacitybacktitle=0.6, rightrule=.15mm, title=\textcolor{quarto-callout-warning-color}{\faExclamationTriangle}\hspace{0.5em}{Advertencia}, titlerule=0mm, toprule=.15mm, toptitle=1mm]

    La densidad de puntos puede facilitar lectura o integración, pero la
    resolución física sigue limitada por \(L/f_s\) y la ventana.

  \end{tcolorbox}
\end{frame}

\begin{frame}{Promediar requiere segmentos comparables}
  \protect\phantomsection\label{promediar-requiere-segmentos-comparables}
  Welch supone que los segmentos describen propiedades espectrales
  suficientemente semejantes.

  Problemas:

  \begin{itemize}
    \tightlist
    \item
          cambio de tarea;
    \item
          transición fisiológica;
    \item
          artefactos concentrados;
    \item
          deriva prolongada.
  \end{itemize}

  \begin{tcolorbox}[enhanced jigsaw, arc=.35mm, bottomrule=.15mm, bottomtitle=1mm, breakable, colback=white, colbacktitle=quarto-callout-warning-color!10!white, colframe=quarto-callout-warning-color-frame, coltitle=black, left=2mm, leftrule=.75mm, opacityback=0, opacitybacktitle=0.6, rightrule=.15mm, title=\textcolor{quarto-callout-warning-color}{\faExclamationTriangle}\hspace{0.5em}{Advertencia}, titlerule=0mm, toprule=.15mm, toptitle=1mm]

    Promediar estados diferentes produce una PSD precisa para una mezcla que
    puede carecer de significado fisiológico.

  \end{tcolorbox}
\end{frame}

\begin{frame}{Ancho de banda puede definirse de varias formas}
  \protect\phantomsection\label{ancho-de-banda-puede-definirse-de-varias-formas}
  Definiciones habituales:

  \begin{itemize}
    \tightlist
    \item
          intervalo entre bordes fijados por especificación;
    \item
          banda que contiene un porcentaje de potencia;
    \item
          frecuencia media o mediana espectral;
    \item
          soporte efectivo por encima de un umbral.
  \end{itemize}

  \begin{tcolorbox}[enhanced jigsaw, arc=.35mm, bottomrule=.15mm, bottomtitle=1mm, breakable, colback=white, colbacktitle=quarto-callout-important-color!10!white, colframe=quarto-callout-important-color-frame, coltitle=black, left=2mm, leftrule=.75mm, opacityback=0, opacitybacktitle=0.6, rightrule=.15mm, title=\textcolor{quarto-callout-important-color}{\faExclamation}\hspace{0.5em}{Importante}, titlerule=0mm, toprule=.15mm, toptitle=1mm]

    ``Ancho de banda'' necesita una definición operacional. Sin ella, la
    comparación no es reproducible.

  \end{tcolorbox}
\end{frame}

\begin{frame}{Frecuencia mediana espectral}
  \protect\phantomsection\label{frecuencia-mediana-espectral}
  La frecuencia \(f_m\) satisface

  \[\int_{0}^{f_m}S_{xx}(f)df
    =\frac{1}{2}\int_{0}^{f_{max}}S_{xx}(f)df.\]

  \begin{tcolorbox}[enhanced jigsaw, arc=.35mm, bottomrule=.15mm, bottomtitle=1mm, breakable, colback=white, colbacktitle=quarto-callout-note-color!10!white, colframe=quarto-callout-note-color-frame, coltitle=black, left=2mm, leftrule=.75mm, opacityback=0, opacitybacktitle=0.6, rightrule=.15mm, title=\textcolor{quarto-callout-note-color}{\faInfo}\hspace{0.5em}{Nota}, titlerule=0mm, toprule=.15mm, toptitle=1mm]

    En EMG puede describir redistribución espectral, pero depende de fatiga,
    fuerza, geometría del registro y preprocesamiento.

  \end{tcolorbox}
\end{frame}

\begin{frame}{Frecuencia media espectral}
  \protect\phantomsection\label{frecuencia-media-espectral}
  \[f_{mean}=\frac{\int fS_{xx}(f)df}{\int S_{xx}(f)df}.\]

  \begin{tcolorbox}[enhanced jigsaw, arc=.35mm, bottomrule=.15mm, bottomtitle=1mm, breakable, colback=white, colbacktitle=quarto-callout-warning-color!10!white, colframe=quarto-callout-warning-color-frame, coltitle=black, left=2mm, leftrule=.75mm, opacityback=0, opacitybacktitle=0.6, rightrule=.15mm, title=\textcolor{quarto-callout-warning-color}{\faExclamationTriangle}\hspace{0.5em}{Advertencia}, titlerule=0mm, toprule=.15mm, toptitle=1mm]

    La frecuencia media es sensible a contaminación de alta frecuencia y a
    los límites escogidos para la integración.

  \end{tcolorbox}
\end{frame}

\begin{frame}[fragile]{Potencia de banda mediante integración numérica}
  \protect\phantomsection\label{potencia-de-banda-mediante-integraciuxf3n-numuxe9rica}
  \begin{Shaded}
    \begin{Highlighting}[]
      \ImportTok{from}\NormalTok{ scipy.signal }\ImportTok{import}\NormalTok{ welch}
      \ImportTok{import}\NormalTok{ numpy }\ImportTok{as}\NormalTok{ np}

      \NormalTok{f, Pxx }\OperatorTok{=}\NormalTok{ welch(x, fs}\OperatorTok{=}\NormalTok{fs, window}\OperatorTok{=}\StringTok{"hann"}\NormalTok{,}
      \NormalTok{               nperseg}\OperatorTok{=}\NormalTok{L, noverlap}\OperatorTok{=}\NormalTok{L}\OperatorTok{//}\DecValTok{2}\NormalTok{,}
      \NormalTok{               detrend}\OperatorTok{=}\StringTok{"constant"}\NormalTok{, scaling}\OperatorTok{=}\StringTok{"density"}\NormalTok{)}
      \NormalTok{mask }\OperatorTok{=}\NormalTok{ (f }\OperatorTok{\textgreater{}=}\NormalTok{ f1) }\OperatorTok{\&}\NormalTok{ (f }\OperatorTok{\textless{}=}\NormalTok{ f2)}
      \NormalTok{P\_band }\OperatorTok{=}\NormalTok{ np.trapz(Pxx[mask], f[mask])}
    \end{Highlighting}
  \end{Shaded}

  \begin{tcolorbox}[enhanced jigsaw, arc=.35mm, bottomrule=.15mm, bottomtitle=1mm, breakable, colback=white, colbacktitle=quarto-callout-note-color!10!white, colframe=quarto-callout-note-color-frame, coltitle=black, left=2mm, leftrule=.75mm, opacityback=0, opacitybacktitle=0.6, rightrule=.15mm, title=\textcolor{quarto-callout-note-color}{\faInfo}\hspace{0.5em}{Nota}, titlerule=0mm, toprule=.15mm, toptitle=1mm]

    La integración debe incluir suficientes puntos y límites coherentes con
    la resolución disponible.

  \end{tcolorbox}
\end{frame}

\begin{frame}{Selección razonada de parámetros}
  \protect\phantomsection\label{selecciuxf3n-razonada-de-paruxe1metros}
  \begin{enumerate}
    \tightlist
    \item
          fijar la mínima separación frecuencial de interés;
    \item
          elegir una duración compatible con esa separación;
    \item
          verificar estacionariedad local;
    \item
          seleccionar ventana según fuga y resolución;
    \item
          usar solapamiento justificado;
    \item
          validar potencia y estabilidad.
  \end{enumerate}

  \begin{tcolorbox}[enhanced jigsaw, arc=.35mm, bottomrule=.15mm, bottomtitle=1mm, breakable, colback=white, colbacktitle=quarto-callout-tip-color!10!white, colframe=quarto-callout-tip-color-frame, coltitle=black, left=2mm, leftrule=.75mm, opacityback=0, opacitybacktitle=0.6, rightrule=.15mm, title=\textcolor{quarto-callout-tip-color}{\faLightbulb}\hspace{0.5em}{Tip}, titlerule=0mm, toprule=.15mm, toptitle=1mm]

    Los parámetros forman parte del método y deben quedar registrados junto
    con el resultado.

  \end{tcolorbox}
\end{frame}

\begin{frame}{Comparación entre condiciones}
  \protect\phantomsection\label{comparaciuxf3n-entre-condiciones}
  Para comparar reposo y tarea:

  \begin{itemize}
    \tightlist
    \item
          igualar unidades y preprocesamiento;
    \item
          usar parámetros espectrales iguales;
    \item
          controlar duración y calidad;
    \item
          comparar por participante cuando corresponda;
    \item
          reportar variabilidad, no solo promedios.
  \end{itemize}

  \begin{tcolorbox}[enhanced jigsaw, arc=.35mm, bottomrule=.15mm, bottomtitle=1mm, breakable, colback=white, colbacktitle=quarto-callout-warning-color!10!white, colframe=quarto-callout-warning-color-frame, coltitle=black, left=2mm, leftrule=.75mm, opacityback=0, opacitybacktitle=0.6, rightrule=.15mm, title=\textcolor{quarto-callout-warning-color}{\faExclamationTriangle}\hspace{0.5em}{Advertencia}, titlerule=0mm, toprule=.15mm, toptitle=1mm]

    Normalizar cada condición por un denominador distinto puede cambiar la
    interpretación de las diferencias.

  \end{tcolorbox}
\end{frame}

\begin{frame}{Verificación de la sesión 17}
  \protect\phantomsection\label{verificaciuxf3n-de-la-sesiuxf3n-17}
  Un registro de 60 s a 1000 Hz usa Welch con segmentos de 2 s y 50 \% de
  solapamiento.

  \begin{enumerate}
    \tightlist
    \item
          determine \(L\) y el avance \(H\);
    \item
          calcule la separación nominal de bins;
    \item
          explique qué cambia con segmentos de 4 s;
    \item
          identifique el supuesto al promediar;
    \item
          proponga una prueba de conservación de potencia.
  \end{enumerate}

  \begin{tcolorbox}[enhanced jigsaw, arc=.35mm, bottomrule=.15mm, bottomtitle=1mm, breakable, colback=white, colbacktitle=quarto-callout-tip-color!10!white, colframe=quarto-callout-tip-color-frame, coltitle=black, left=2mm, leftrule=.75mm, opacityback=0, opacitybacktitle=0.6, rightrule=.15mm, title=\textcolor{quarto-callout-tip-color}{\faLightbulb}\hspace{0.5em}{Criterio de logro}, titlerule=0mm, toprule=.15mm, toptitle=1mm]

    La respuesta debe discutir resolución, varianza, dependencia y
    estacionariedad.

  \end{tcolorbox}
\end{frame}

\begin{frame}{Sesión 18 · Una señal no estacionaria cambia durante el
    registro}
  \protect\phantomsection\label{sesiuxf3n-18-una-seuxf1al-no-estacionaria-cambia-durante-el-registro}
  Ejemplos:

  \begin{itemize}
    \tightlist
    \item
          transición de reposo a ejercicio;
    \item
          ráfagas EEG;
    \item
          activación muscular durante marcha;
    \item
          cambios de ritmo cardíaco.
  \end{itemize}

  \begin{tcolorbox}[enhanced jigsaw, arc=.35mm, bottomrule=.15mm, bottomtitle=1mm, breakable, colback=white, colbacktitle=quarto-callout-important-color!10!white, colframe=quarto-callout-important-color-frame, coltitle=black, left=2mm, leftrule=.75mm, opacityback=0, opacitybacktitle=0.6, rightrule=.15mm, title=\textcolor{quarto-callout-important-color}{\faExclamation}\hspace{0.5em}{Importante}, titlerule=0mm, toprule=.15mm, toptitle=1mm]

    Un espectro global mezcla tiempos. El análisis tiempo-frecuencia busca
    localizar cómo cambia el contenido espectral.

  \end{tcolorbox}
\end{frame}

\begin{frame}{La estacionariedad es un supuesto de modelamiento}
  \protect\phantomsection\label{la-estacionariedad-es-un-supuesto-de-modelamiento}
  En sentido amplio:

  \[\mathbb{E}\{x(t)\}=\mu,\]

  \[R_x(t_1,t_2)=R_x(t_1-t_2).\]

  \begin{tcolorbox}[enhanced jigsaw, arc=.35mm, bottomrule=.15mm, bottomtitle=1mm, breakable, colback=white, colbacktitle=quarto-callout-warning-color!10!white, colframe=quarto-callout-warning-color-frame, coltitle=black, left=2mm, leftrule=.75mm, opacityback=0, opacitybacktitle=0.6, rightrule=.15mm, title=\textcolor{quarto-callout-warning-color}{\faExclamationTriangle}\hspace{0.5em}{Advertencia}, titlerule=0mm, toprule=.15mm, toptitle=1mm]

    Las bioseñales rara vez cumplen estacionariedad global. Se trabaja con
    segmentos donde las propiedades cambian poco para el objetivo.

  \end{tcolorbox}
\end{frame}

\begin{frame}{La STFT aplica Fourier localmente}
  \protect\phantomsection\label{la-stft-aplica-fourier-localmente}
  \[X(\tau,f)=\int_{-\infty}^{\infty}x(t)w(t-\tau)e^{-j2\pi ft}dt.\]

  \begin{itemize}
    \tightlist
    \item
          \(\tau\): posición temporal de la ventana.
    \item
          \(f\): frecuencia analizada.
    \item
          \(w\): ventana de análisis.
  \end{itemize}

  \begin{tcolorbox}[enhanced jigsaw, arc=.35mm, bottomrule=.15mm, bottomtitle=1mm, breakable, colback=white, colbacktitle=quarto-callout-important-color!10!white, colframe=quarto-callout-important-color-frame, coltitle=black, left=2mm, leftrule=.75mm, opacityback=0, opacitybacktitle=0.6, rightrule=.15mm, title=\textcolor{quarto-callout-important-color}{\faExclamation}\hspace{0.5em}{Importante}, titlerule=0mm, toprule=.15mm, toptitle=1mm]

    La STFT produce una colección de espectros locales con una ventana de
    forma y longitud fijas.

  \end{tcolorbox}
\end{frame}

\begin{frame}{El espectrograma muestra magnitud o potencia}
  \protect\phantomsection\label{el-espectrograma-muestra-magnitud-o-potencia}
  Una definición habitual:

  \[\mathcal{S}(\tau,f)=|X(\tau,f)|^2.\]

  \begin{tcolorbox}[enhanced jigsaw, arc=.35mm, bottomrule=.15mm, bottomtitle=1mm, breakable, colback=white, colbacktitle=quarto-callout-warning-color!10!white, colframe=quarto-callout-warning-color-frame, coltitle=black, left=2mm, leftrule=.75mm, opacityback=0, opacitybacktitle=0.6, rightrule=.15mm, title=\textcolor{quarto-callout-warning-color}{\faExclamationTriangle}\hspace{0.5em}{Advertencia}, titlerule=0mm, toprule=.15mm, toptitle=1mm]

    Color, escala logarítmica y normalización cambian la apariencia. Una
    imagen atractiva no garantiza unidades correctas.

  \end{tcolorbox}
\end{frame}

\begin{frame}{La ventana fija impone un compromiso}
  \protect\phantomsection\label{la-ventana-fija-impone-un-compromiso}
  Ventana corta:

  \begin{itemize}
    \tightlist
    \item
          mejor localización temporal;
    \item
          peor separación frecuencial.
  \end{itemize}

  Ventana larga:

  \begin{itemize}
    \tightlist
    \item
          mejor separación frecuencial;
    \item
          mayor mezcla temporal.
  \end{itemize}

  \begin{tcolorbox}[enhanced jigsaw, arc=.35mm, bottomrule=.15mm, bottomtitle=1mm, breakable, colback=white, colbacktitle=quarto-callout-important-color!10!white, colframe=quarto-callout-important-color-frame, coltitle=black, left=2mm, leftrule=.75mm, opacityback=0, opacitybacktitle=0.6, rightrule=.15mm, title=\textcolor{quarto-callout-important-color}{\faExclamation}\hspace{0.5em}{Importante}, titlerule=0mm, toprule=.15mm, toptitle=1mm]

    La STFT no ofrece máxima resolución en ambos dominios. La pregunta
    define el compromiso aceptable.

  \end{tcolorbox}
\end{frame}

\begin{frame}{La incertidumbre expresa ese límite}
  \protect\phantomsection\label{la-incertidumbre-expresa-ese-luxedmite}
  De forma conceptual:

  \[\Delta t\,\Delta f\geq C,\]

  donde \(C\) depende de la definición y la ventana.

  \begin{tcolorbox}[enhanced jigsaw, arc=.35mm, bottomrule=.15mm, bottomtitle=1mm, breakable, colback=white, colbacktitle=quarto-callout-note-color!10!white, colframe=quarto-callout-note-color-frame, coltitle=black, left=2mm, leftrule=.75mm, opacityback=0, opacitybacktitle=0.6, rightrule=.15mm, title=\textcolor{quarto-callout-note-color}{\faInfo}\hspace{0.5em}{Nota}, titlerule=0mm, toprule=.15mm, toptitle=1mm]

    Reducir indefinidamente la ventana mejora tiempo, pero ensancha la
    respuesta frecuencial de cada componente.

  \end{tcolorbox}
\end{frame}

\begin{frame}{El avance controla densidad temporal}
  \protect\phantomsection\label{el-avance-controla-densidad-temporal}
  Con ventana \(L\) y avance \(H\):

  \[t_r=\frac{rH+\text{referencia de ventana}}{f_s}.\]

  \begin{tcolorbox}[enhanced jigsaw, arc=.35mm, bottomrule=.15mm, bottomtitle=1mm, breakable, colback=white, colbacktitle=quarto-callout-warning-color!10!white, colframe=quarto-callout-warning-color-frame, coltitle=black, left=2mm, leftrule=.75mm, opacityback=0, opacitybacktitle=0.6, rightrule=.15mm, title=\textcolor{quarto-callout-warning-color}{\faExclamationTriangle}\hspace{0.5em}{Advertencia}, titlerule=0mm, toprule=.15mm, toptitle=1mm]

    Mayor solapamiento produce más columnas en el espectrograma, pero no
    aumenta la resolución temporal física en la misma proporción.

  \end{tcolorbox}
\end{frame}

\begin{frame}{Las condiciones de borde necesitan atención}
  \protect\phantomsection\label{las-condiciones-de-borde-necesitan-atenciuxf3n}
  Cerca del inicio y final, la ventana puede:

  \begin{itemize}
    \tightlist
    \item
          contener menos datos reales;
    \item
          usar relleno;
    \item
          introducir transitorios de borde;
    \item
          cambiar la interpretación temporal.
  \end{itemize}

  \begin{tcolorbox}[enhanced jigsaw, arc=.35mm, bottomrule=.15mm, bottomtitle=1mm, breakable, colback=white, colbacktitle=quarto-callout-important-color!10!white, colframe=quarto-callout-important-color-frame, coltitle=black, left=2mm, leftrule=.75mm, opacityback=0, opacitybacktitle=0.6, rightrule=.15mm, title=\textcolor{quarto-callout-important-color}{\faExclamation}\hspace{0.5em}{Importante}, titlerule=0mm, toprule=.15mm, toptitle=1mm]

    Debe documentarse si se usa padding, centrado de ventanas y extensión
    del registro.

  \end{tcolorbox}
\end{frame}

\begin{frame}{La escala logarítmica revela componentes débiles}
  \protect\phantomsection\label{la-escala-logaruxedtmica-revela-componentes-duxe9biles}
  Una representación frecuente es

  \[P_{dB}(\tau,f)=10\log_{10}\left(\frac{P(\tau,f)}{P_{ref}}\right).\]

  \begin{tcolorbox}[enhanced jigsaw, arc=.35mm, bottomrule=.15mm, bottomtitle=1mm, breakable, colback=white, colbacktitle=quarto-callout-warning-color!10!white, colframe=quarto-callout-warning-color-frame, coltitle=black, left=2mm, leftrule=.75mm, opacityback=0, opacitybacktitle=0.6, rightrule=.15mm, title=\textcolor{quarto-callout-warning-color}{\faExclamationTriangle}\hspace{0.5em}{Advertencia}, titlerule=0mm, toprule=.15mm, toptitle=1mm]

    Sin \(P_{ref}\) y unidades, ``dB'' queda incompleto. Ceros o valores muy
    pequeños requieren un piso numérico declarado.

  \end{tcolorbox}
\end{frame}

\begin{frame}{STFT de una activación EMG}
  \protect\phantomsection\label{stft-de-una-activaciuxf3n-emg}
  Una ventana corta puede localizar inicio y final de activación. Una
  ventana más larga estabiliza la distribución espectral dentro de la
  contracción.

  \begin{tcolorbox}[enhanced jigsaw, arc=.35mm, bottomrule=.15mm, bottomtitle=1mm, breakable, colback=white, colbacktitle=quarto-callout-note-color!10!white, colframe=quarto-callout-note-color-frame, coltitle=black, left=2mm, leftrule=.75mm, opacityback=0, opacitybacktitle=0.6, rightrule=.15mm, title=\textcolor{quarto-callout-note-color}{\faInfo}\hspace{0.5em}{Nota}, titlerule=0mm, toprule=.15mm, toptitle=1mm]

    La envolvente responde a variación de amplitud lenta. El espectrograma
    describe además cómo se distribuye energía dentro de la banda.

  \end{tcolorbox}
\end{frame}

\begin{frame}{STFT de una ráfaga EEG}
  \protect\phantomsection\label{stft-de-una-ruxe1faga-eeg}
  Preguntas relevantes:

  \begin{itemize}
    \tightlist
    \item
          ¿cuándo comienza la ráfaga?
    \item
          ¿qué banda domina?
    \item
          ¿cambia su frecuencia central?
    \item
          ¿aparece simultáneamente en varios canales?
  \end{itemize}

  \begin{tcolorbox}[enhanced jigsaw, arc=.35mm, bottomrule=.15mm, bottomtitle=1mm, breakable, colback=white, colbacktitle=quarto-callout-warning-color!10!white, colframe=quarto-callout-warning-color-frame, coltitle=black, left=2mm, leftrule=.75mm, opacityback=0, opacitybacktitle=0.6, rightrule=.15mm, title=\textcolor{quarto-callout-warning-color}{\faExclamationTriangle}\hspace{0.5em}{Advertencia}, titlerule=0mm, toprule=.15mm, toptitle=1mm]

    Una banda visible no identifica por sí sola una fuente neural.
    Artefactos o referencia común pueden generar patrones similares.

  \end{tcolorbox}
\end{frame}

\begin{frame}[fragile]{Código: espectrograma con parámetros explícitos}
  \protect\phantomsection\label{cuxf3digo-espectrograma-con-paruxe1metros-expluxedcitos}
  \begin{Shaded}
    \begin{Highlighting}[]
      \ImportTok{from}\NormalTok{ scipy.signal }\ImportTok{import}\NormalTok{ stft}

      \NormalTok{f, t, Z }\OperatorTok{=}\NormalTok{ stft(}
      \NormalTok{    x, fs}\OperatorTok{=}\NormalTok{fs, window}\OperatorTok{=}\StringTok{"hann"}\NormalTok{,}
      \NormalTok{    nperseg}\OperatorTok{=}\NormalTok{L, noverlap}\OperatorTok{=}\NormalTok{L}\OperatorTok{{-}}\NormalTok{H,}
      \NormalTok{    nfft}\OperatorTok{=}\NormalTok{Nfft, detrend}\OperatorTok{=}\StringTok{"constant"}\NormalTok{,}
      \NormalTok{    boundary}\OperatorTok{=}\VariableTok{None}\NormalTok{, padded}\OperatorTok{=}\VariableTok{False}\NormalTok{,}
      \NormalTok{)}
      \NormalTok{P }\OperatorTok{=}\NormalTok{ np.}\BuiltInTok{abs}\NormalTok{(Z)}\OperatorTok{**}\DecValTok{2}
    \end{Highlighting}
  \end{Shaded}

  \begin{tcolorbox}[enhanced jigsaw, arc=.35mm, bottomrule=.15mm, bottomtitle=1mm, breakable, colback=white, colbacktitle=quarto-callout-note-color!10!white, colframe=quarto-callout-note-color-frame, coltitle=black, left=2mm, leftrule=.75mm, opacityback=0, opacitybacktitle=0.6, rightrule=.15mm, title=\textcolor{quarto-callout-note-color}{\faInfo}\hspace{0.5em}{Nota}, titlerule=0mm, toprule=.15mm, toptitle=1mm]

    La conversión de \(|Z|^2\) a PSD requiere revisar la normalización de la
    función y el propósito del análisis.

  \end{tcolorbox}
\end{frame}

\begin{frame}{Validar un espectrograma}
  \protect\phantomsection\label{validar-un-espectrograma}
  \begin{enumerate}
    \tightlist
    \item
          contrastar eventos con el dominio temporal;
    \item
          verificar frecuencias con señales sintéticas;
    \item
          revisar efectos de ventana y borde;
    \item
          repetir con otra longitud razonable;
    \item
          comparar potencia integrada con estimaciones temporales;
    \item
          inspeccionar canales de artefacto.
  \end{enumerate}

  \begin{tcolorbox}[enhanced jigsaw, arc=.35mm, bottomrule=.15mm, bottomtitle=1mm, breakable, colback=white, colbacktitle=quarto-callout-tip-color!10!white, colframe=quarto-callout-tip-color-frame, coltitle=black, left=2mm, leftrule=.75mm, opacityback=0, opacitybacktitle=0.6, rightrule=.15mm, title=\textcolor{quarto-callout-tip-color}{\faLightbulb}\hspace{0.5em}{Tip}, titlerule=0mm, toprule=.15mm, toptitle=1mm]

    Una interpretación robusta debe persistir bajo cambios moderados de
    parámetros compatibles con la pregunta.

  \end{tcolorbox}
\end{frame}

\begin{frame}{Verificación de la sesión 18}
  \protect\phantomsection\label{verificaciuxf3n-de-la-sesiuxf3n-18}
  Una ráfaga de 20 Hz dura 300 ms en un EEG a 250 Hz.

  \begin{enumerate}
    \tightlist
    \item
          compare ventanas de 128 ms y 1 s;
    \item
          indique cuál localiza mejor el inicio;
    \item
          indique cuál separa mejor 20 Hz de 22 Hz;
    \item
          discuta el efecto del solapamiento;
    \item
          proponga canales o marcas para validar origen.
  \end{enumerate}

  \begin{tcolorbox}[enhanced jigsaw, arc=.35mm, bottomrule=.15mm, bottomtitle=1mm, breakable, colback=white, colbacktitle=quarto-callout-tip-color!10!white, colframe=quarto-callout-tip-color-frame, coltitle=black, left=2mm, leftrule=.75mm, opacityback=0, opacitybacktitle=0.6, rightrule=.15mm, title=\textcolor{quarto-callout-tip-color}{\faLightbulb}\hspace{0.5em}{Criterio de logro}, titlerule=0mm, toprule=.15mm, toptitle=1mm]

    La respuesta debe conectar duración del evento, lóbulo temporal,
    resolución y evidencia fisiológica.

  \end{tcolorbox}
\end{frame}

\begin{frame}{Cierre de la semana 9}
  \protect\phantomsection\label{cierre-de-la-semana-9}
  \begin{itemize}
    \tightlist
    \item
          Welch promedia periodogramas modificados para reducir variabilidad.
    \item
          La longitud de segmento controla resolución y aproximación local.
    \item
          Ancho de banda necesita una definición reproducible.
    \item
          La STFT localiza cambios espectrales con resolución limitada.
    \item
          Ventana, solapamiento, padding y escala forman parte del resultado.
  \end{itemize}

  \begin{tcolorbox}[enhanced jigsaw, arc=.35mm, bottomrule=.15mm, bottomtitle=1mm, breakable, colback=white, colbacktitle=quarto-callout-note-color!10!white, colframe=quarto-callout-note-color-frame, coltitle=black, left=2mm, leftrule=.75mm, opacityback=0, opacitybacktitle=0.6, rightrule=.15mm, title=\textcolor{quarto-callout-note-color}{\faInfo}\hspace{0.5em}{Pregunta de salida}, titlerule=0mm, toprule=.15mm, toptitle=1mm]

    ¿Por qué un espectrograma con más columnas puede no contener mejor
    resolución temporal?

  \end{tcolorbox}
\end{frame}

\section{Semana 10 · Ventanas tiempo-frecuencia y sistemas
  LTI}\label{semana-10-ventanas-tiempo-frecuencia-y-sistemas-lti}

\begin{frame}{Sesión 19 · La ventana debe responder a la pregunta}
  \protect\phantomsection\label{sesiuxf3n-19-la-ventana-debe-responder-a-la-pregunta}
  Tres preguntas dominan la selección:

  \begin{enumerate}
    \tightlist
    \item
          ¿qué duración tiene el evento?
    \item
          ¿qué separación frecuencial se necesita?
    \item
          ¿qué diferencia de amplitud existe entre componentes?
  \end{enumerate}

  \begin{tcolorbox}[enhanced jigsaw, arc=.35mm, bottomrule=.15mm, bottomtitle=1mm, breakable, colback=white, colbacktitle=quarto-callout-important-color!10!white, colframe=quarto-callout-important-color-frame, coltitle=black, left=2mm, leftrule=.75mm, opacityback=0, opacitybacktitle=0.6, rightrule=.15mm, title=\textcolor{quarto-callout-important-color}{\faExclamation}\hspace{0.5em}{Importante}, titlerule=0mm, toprule=.15mm, toptitle=1mm]

    La forma y longitud de ventana se eligen conjuntamente. Ninguna ventana
    domina en todos los criterios.

  \end{tcolorbox}
\end{frame}

\begin{frame}{La longitud fija una escala de observación}
  \protect\phantomsection\label{la-longitud-fija-una-escala-de-observaciuxf3n}
  Para \(L\) muestras:

  \[T_w=\frac{L}{f_s},\qquad \Delta f_{nom}=\frac{f_s}{L}.\]

  \begin{tcolorbox}[enhanced jigsaw, arc=.35mm, bottomrule=.15mm, bottomtitle=1mm, breakable, colback=white, colbacktitle=quarto-callout-warning-color!10!white, colframe=quarto-callout-warning-color-frame, coltitle=black, left=2mm, leftrule=.75mm, opacityback=0, opacitybacktitle=0.6, rightrule=.15mm, title=\textcolor{quarto-callout-warning-color}{\faExclamationTriangle}\hspace{0.5em}{Advertencia}, titlerule=0mm, toprule=.15mm, toptitle=1mm]

    La resolución efectiva depende además de la forma de la ventana y del
    criterio usado para distinguir componentes.

  \end{tcolorbox}
\end{frame}

\begin{frame}{Rectangular prioriza un lóbulo principal estrecho}
  \protect\phantomsection\label{rectangular-prioriza-un-luxf3bulo-principal-estrecho}
  Ventajas:

  \begin{itemize}
    \tightlist
    \item
          máxima conservación de un bloque alineado;
    \item
          buena separación nominal de tonos comparables.
  \end{itemize}

  Limitaciones:

  \begin{itemize}
    \tightlist
    \item
          lóbulos laterales altos;
    \item
          sensibilidad a discontinuidades de borde.
  \end{itemize}

  \begin{tcolorbox}[enhanced jigsaw, arc=.35mm, bottomrule=.15mm, bottomtitle=1mm, breakable, colback=white, colbacktitle=quarto-callout-note-color!10!white, colframe=quarto-callout-note-color-frame, coltitle=black, left=2mm, leftrule=.75mm, opacityback=0, opacitybacktitle=0.6, rightrule=.15mm, title=\textcolor{quarto-callout-note-color}{\faInfo}\hspace{0.5em}{Nota}, titlerule=0mm, toprule=.15mm, toptitle=1mm]

    Resulta útil como referencia para comprender por qué el recorte directo
    produce fuga.

  \end{tcolorbox}
\end{frame}

\begin{frame}{Hann ofrece un compromiso general}
  \protect\phantomsection\label{hann-ofrece-un-compromiso-general}
  Ventajas:

  \begin{itemize}
    \tightlist
    \item
          extremos nulos;
    \item
          menor fuga lateral que rectangular;
    \item
          uso frecuente en Welch y STFT.
  \end{itemize}

  Limitación principal: lóbulo central más ancho.

  \begin{tcolorbox}[enhanced jigsaw, arc=.35mm, bottomrule=.15mm, bottomtitle=1mm, breakable, colback=white, colbacktitle=quarto-callout-important-color!10!white, colframe=quarto-callout-important-color-frame, coltitle=black, left=2mm, leftrule=.75mm, opacityback=0, opacitybacktitle=0.6, rightrule=.15mm, title=\textcolor{quarto-callout-important-color}{\faExclamation}\hspace{0.5em}{Importante}, titlerule=0mm, toprule=.15mm, toptitle=1mm]

    Hann suele ser un punto de partida defendible, no una selección
    automática.

  \end{tcolorbox}
\end{frame}

\begin{frame}{Hamming mantiene extremos pequeños, no nulos}
  \protect\phantomsection\label{hamming-mantiene-extremos-pequeuxf1os-no-nulos}
  Hamming reduce especialmente el primer lóbulo lateral, pero su extensión
  periódica mantiene una discontinuidad pequeña.

  \begin{tcolorbox}[enhanced jigsaw, arc=.35mm, bottomrule=.15mm, bottomtitle=1mm, breakable, colback=white, colbacktitle=quarto-callout-note-color!10!white, colframe=quarto-callout-note-color-frame, coltitle=black, left=2mm, leftrule=.75mm, opacityback=0, opacitybacktitle=0.6, rightrule=.15mm, title=\textcolor{quarto-callout-note-color}{\faInfo}\hspace{0.5em}{Nota}, titlerule=0mm, toprule=.15mm, toptitle=1mm]

    La diferencia frente a Hann puede ser relevante para estimación tonal y
    respuesta de filtros, aunque no siempre cambia una conclusión biomédica.

  \end{tcolorbox}
\end{frame}

\begin{frame}{Blackman prioriza rango dinámico}
  \protect\phantomsection\label{blackman-prioriza-rango-dinuxe1mico}
  Sus lóbulos laterales bajos ayudan a observar componentes débiles cerca
  de componentes intensas.

  Costo: lóbulo principal más ancho y menor separación de frecuencias
  próximas.

  \begin{tcolorbox}[enhanced jigsaw, arc=.35mm, bottomrule=.15mm, bottomtitle=1mm, breakable, colback=white, colbacktitle=quarto-callout-warning-color!10!white, colframe=quarto-callout-warning-color-frame, coltitle=black, left=2mm, leftrule=.75mm, opacityback=0, opacitybacktitle=0.6, rightrule=.15mm, title=\textcolor{quarto-callout-warning-color}{\faExclamationTriangle}\hspace{0.5em}{Advertencia}, titlerule=0mm, toprule=.15mm, toptitle=1mm]

    Reducir fuga lateral no mejora simultáneamente todas las formas de
    resolución.

  \end{tcolorbox}
\end{frame}

\begin{frame}{Una ventana gaussiana localiza de forma equilibrada}
  \protect\phantomsection\label{una-ventana-gaussiana-localiza-de-forma-equilibrada}
  La familia gaussiana permite ajustar el compromiso mediante su
  dispersión y posee buenas propiedades de concentración
  tiempo-frecuencia.

  \begin{tcolorbox}[enhanced jigsaw, arc=.35mm, bottomrule=.15mm, bottomtitle=1mm, breakable, colback=white, colbacktitle=quarto-callout-note-color!10!white, colframe=quarto-callout-note-color-frame, coltitle=black, left=2mm, leftrule=.75mm, opacityback=0, opacitybacktitle=0.6, rightrule=.15mm, title=\textcolor{quarto-callout-note-color}{\faInfo}\hspace{0.5em}{Nota}, titlerule=0mm, toprule=.15mm, toptitle=1mm]

    La elección del parámetro debe quedar documentada. ``Gaussiana'' no
    define por sí sola una ventana única.

  \end{tcolorbox}
\end{frame}

\begin{frame}{Ventana y evento deben tener escalas compatibles}
  \protect\phantomsection\label{ventana-y-evento-deben-tener-escalas-compatibles}
  Ejemplos:

  \begin{itemize}
    \tightlist
    \item
          QRS: decenas a pocos cientos de milisegundos;
    \item
          activación EMG: depende de la tarea y músculo;
    \item
          oscilaciones respiratorias: varios segundos;
    \item
          ritmos EEG sostenidos: duración definida por hipótesis.
  \end{itemize}

  \begin{tcolorbox}[enhanced jigsaw, arc=.35mm, bottomrule=.15mm, bottomtitle=1mm, breakable, colback=white, colbacktitle=quarto-callout-important-color!10!white, colframe=quarto-callout-important-color-frame, coltitle=black, left=2mm, leftrule=.75mm, opacityback=0, opacitybacktitle=0.6, rightrule=.15mm, title=\textcolor{quarto-callout-important-color}{\faExclamation}\hspace{0.5em}{Importante}, titlerule=0mm, toprule=.15mm, toptitle=1mm]

    Una ventana mucho mayor que el evento diluye su localización. Una
    ventana demasiado corta pierde detalle frecuencial.

  \end{tcolorbox}
\end{frame}

\begin{frame}{La frecuencia de interés determina ciclos observados}
  \protect\phantomsection\label{la-frecuencia-de-interuxe9s-determina-ciclos-observados}
  Número aproximado de ciclos dentro de la ventana:

  \[N_c=f_0T_w.\]

  \begin{tcolorbox}[enhanced jigsaw, arc=.35mm, bottomrule=.15mm, bottomtitle=1mm, breakable, colback=white, colbacktitle=quarto-callout-warning-color!10!white, colframe=quarto-callout-warning-color-frame, coltitle=black, left=2mm, leftrule=.75mm, opacityback=0, opacitybacktitle=0.6, rightrule=.15mm, title=\textcolor{quarto-callout-warning-color}{\faExclamationTriangle}\hspace{0.5em}{Advertencia}, titlerule=0mm, toprule=.15mm, toptitle=1mm]

    Estimar una frecuencia a partir de menos de un ciclo completo es
    inestable y depende fuertemente del modelo.

  \end{tcolorbox}
\end{frame}

\begin{frame}{Solapamiento mejora seguimiento visual}
  \protect\phantomsection\label{solapamiento-mejora-seguimiento-visual}
  Reducir el avance \(H\) produce estimaciones más densas en el tiempo.

  \begin{tcolorbox}[enhanced jigsaw, arc=.35mm, bottomrule=.15mm, bottomtitle=1mm, breakable, colback=white, colbacktitle=quarto-callout-warning-color!10!white, colframe=quarto-callout-warning-color-frame, coltitle=black, left=2mm, leftrule=.75mm, opacityback=0, opacitybacktitle=0.6, rightrule=.15mm, title=\textcolor{quarto-callout-warning-color}{\faExclamationTriangle}\hspace{0.5em}{Advertencia}, titlerule=0mm, toprule=.15mm, toptitle=1mm]

    Las columnas vecinas comparten datos y son dependientes. No deben
    contarse como observaciones independientes en análisis estadístico.

  \end{tcolorbox}
\end{frame}

\begin{frame}{Zero padding facilita localizar máximos}
  \protect\phantomsection\label{zero-padding-facilita-localizar-muxe1ximos}
  Una malla frecuencial densa puede ayudar a interpolar el máximo de un
  lóbulo.

  \begin{tcolorbox}[enhanced jigsaw, arc=.35mm, bottomrule=.15mm, bottomtitle=1mm, breakable, colback=white, colbacktitle=quarto-callout-important-color!10!white, colframe=quarto-callout-important-color-frame, coltitle=black, left=2mm, leftrule=.75mm, opacityback=0, opacitybacktitle=0.6, rightrule=.15mm, title=\textcolor{quarto-callout-important-color}{\faExclamation}\hspace{0.5em}{Importante}, titlerule=0mm, toprule=.15mm, toptitle=1mm]

    La precisión de la estimación sigue limitada por SNR, duración, ventana
    y estabilidad de la frecuencia.

  \end{tcolorbox}
\end{frame}

\begin{frame}{Normalizar entre ventanas requiere coherencia}
  \protect\phantomsection\label{normalizar-entre-ventanas-requiere-coherencia}
  Opciones distintas responden preguntas distintas:

  \begin{itemize}
    \tightlist
    \item
          potencia absoluta;
    \item
          potencia relativa por ventana;
    \item
          cambio respecto de una línea base;
    \item
          puntuación estandarizada por frecuencia.
  \end{itemize}

  \begin{tcolorbox}[enhanced jigsaw, arc=.35mm, bottomrule=.15mm, bottomtitle=1mm, breakable, colback=white, colbacktitle=quarto-callout-warning-color!10!white, colframe=quarto-callout-warning-color-frame, coltitle=black, left=2mm, leftrule=.75mm, opacityback=0, opacitybacktitle=0.6, rightrule=.15mm, title=\textcolor{quarto-callout-warning-color}{\faExclamationTriangle}\hspace{0.5em}{Advertencia}, titlerule=0mm, toprule=.15mm, toptitle=1mm]

    Normalizar cada columna por su propia potencia elimina cambios globales
    de amplitud y puede ocultar activaciones reales.

  \end{tcolorbox}
\end{frame}

\begin{frame}{Aplicación: activación muscular durante marcha}
  \protect\phantomsection\label{aplicaciuxf3n-activaciuxf3n-muscular-durante-marcha}
  Una STFT puede describir cuándo aumenta energía EMG y si su distribución
  cambia dentro del ciclo.

  Diseño:

  \begin{itemize}
    \tightlist
    \item
          alinear ciclos con eventos de marcha;
    \item
          conservar escala o usar referencia fisiológica;
    \item
          separar análisis temporal y espectral;
    \item
          validar artefactos de movimiento.
  \end{itemize}

  \begin{tcolorbox}[enhanced jigsaw, arc=.35mm, bottomrule=.15mm, bottomtitle=1mm, breakable, colback=white, colbacktitle=quarto-callout-important-color!10!white, colframe=quarto-callout-important-color-frame, coltitle=black, left=2mm, leftrule=.75mm, opacityback=0, opacitybacktitle=0.6, rightrule=.15mm, title=\textcolor{quarto-callout-important-color}{\faExclamation}\hspace{0.5em}{Importante}, titlerule=0mm, toprule=.15mm, toptitle=1mm]

    El espectrograma debe interpretarse junto con el ciclo de marcha y no
    como una imagen aislada.

  \end{tcolorbox}
\end{frame}

\begin{frame}{Aplicación: cambio de ritmo respiratorio}
  \protect\phantomsection\label{aplicaciuxf3n-cambio-de-ritmo-respiratorio}
  Ventanas largas separan frecuencias respiratorias próximas. Ventanas
  cortas localizan transiciones de ritmo.

  \begin{tcolorbox}[enhanced jigsaw, arc=.35mm, bottomrule=.15mm, bottomtitle=1mm, breakable, colback=white, colbacktitle=quarto-callout-note-color!10!white, colframe=quarto-callout-note-color-frame, coltitle=black, left=2mm, leftrule=.75mm, opacityback=0, opacitybacktitle=0.6, rightrule=.15mm, title=\textcolor{quarto-callout-note-color}{\faInfo}\hspace{0.5em}{Nota}, titlerule=0mm, toprule=.15mm, toptitle=1mm]

    Una estrategia útil compara varias longitudes plausibles y reporta qué
    conclusión permanece estable.

  \end{tcolorbox}
\end{frame}

\begin{frame}{Flujo de selección documentada}
  \protect\phantomsection\label{flujo-de-selecciuxf3n-documentada}
  \begin{enumerate}
    \tightlist
    \item
          formular evento o banda de interés;
    \item
          estimar su escala temporal;
    \item
          fijar separación frecuencial necesaria;
    \item
          elegir longitud y forma;
    \item
          definir avance y padding;
    \item
          validar con señales conocidas;
    \item
          realizar análisis de sensibilidad.
  \end{enumerate}

  \begin{tcolorbox}[enhanced jigsaw, arc=.35mm, bottomrule=.15mm, bottomtitle=1mm, breakable, colback=white, colbacktitle=quarto-callout-tip-color!10!white, colframe=quarto-callout-tip-color-frame, coltitle=black, left=2mm, leftrule=.75mm, opacityback=0, opacitybacktitle=0.6, rightrule=.15mm, title=\textcolor{quarto-callout-tip-color}{\faLightbulb}\hspace{0.5em}{Tip}, titlerule=0mm, toprule=.15mm, toptitle=1mm]

    La selección final debe poder reproducirse a partir de parámetros
    guardados, no de ajustes visuales manuales.

  \end{tcolorbox}
\end{frame}

\begin{frame}{Verificación de la sesión 19}
  \protect\phantomsection\label{verificaciuxf3n-de-la-sesiuxf3n-19}
  Se desea localizar una activación de 250 ms y distinguir componentes
  separadas 2 Hz.

  \begin{enumerate}
    \tightlist
    \item
          analice la tensión entre ambos objetivos;
    \item
          calcule la duración mínima nominal para 2 Hz;
    \item
          proponga una longitud inicial;
    \item
          seleccione una ventana y justifique;
    \item
          describa un análisis de sensibilidad.
  \end{enumerate}

  \begin{tcolorbox}[enhanced jigsaw, arc=.35mm, bottomrule=.15mm, bottomtitle=1mm, breakable, colback=white, colbacktitle=quarto-callout-tip-color!10!white, colframe=quarto-callout-tip-color-frame, coltitle=black, left=2mm, leftrule=.75mm, opacityback=0, opacitybacktitle=0.6, rightrule=.15mm, title=\textcolor{quarto-callout-tip-color}{\faLightbulb}\hspace{0.5em}{Criterio de logro}, titlerule=0mm, toprule=.15mm, toptitle=1mm]

    La solución debe reconocer que una única STFT puede no satisfacer
    simultáneamente ambos objetivos.

  \end{tcolorbox}
\end{frame}

\begin{frame}{Sesión 20 · Un sistema transforma una señal}
  \protect\phantomsection\label{sesiuxf3n-20-un-sistema-transforma-una-seuxf1al}
  \[y=\mathcal{T}\{x\}.\]

  El análisis pregunta qué propiedades de \(\mathcal{T}\) permiten
  predecir la salida para entradas diferentes.

  \begin{tcolorbox}[enhanced jigsaw, arc=.35mm, bottomrule=.15mm, bottomtitle=1mm, breakable, colback=white, colbacktitle=quarto-callout-important-color!10!white, colframe=quarto-callout-important-color-frame, coltitle=black, left=2mm, leftrule=.75mm, opacityback=0, opacitybacktitle=0.6, rightrule=.15mm, title=\textcolor{quarto-callout-important-color}{\faExclamation}\hspace{0.5em}{Importante}, titlerule=0mm, toprule=.15mm, toptitle=1mm]

    Las propiedades pertenecen al modelo del sistema dentro de condiciones
    definidas. Un sistema fisiológico puede ser aproximadamente lineal solo
    en un rango.

  \end{tcolorbox}
\end{frame}

\begin{frame}{Linealidad combina aditividad y homogeneidad}
  \protect\phantomsection\label{linealidad-combina-aditividad-y-homogeneidad}
  \[\mathcal{T}\{a x_1+b x_2\}
    =a\mathcal{T}\{x_1\}+b\mathcal{T}\{x_2\}.\]

  \begin{tcolorbox}[enhanced jigsaw, arc=.35mm, bottomrule=.15mm, bottomtitle=1mm, breakable, colback=white, colbacktitle=quarto-callout-warning-color!10!white, colframe=quarto-callout-warning-color-frame, coltitle=black, left=2mm, leftrule=.75mm, opacityback=0, opacitybacktitle=0.6, rightrule=.15mm, title=\textcolor{quarto-callout-warning-color}{\faExclamationTriangle}\hspace{0.5em}{Advertencia}, titlerule=0mm, toprule=.15mm, toptitle=1mm]

    Verificar una sola entrada no demuestra linealidad. La propiedad debe
    sostenerse para las combinaciones consideradas.

  \end{tcolorbox}
\end{frame}

\begin{frame}{Saturación viola linealidad}
  \protect\phantomsection\label{saturaciuxf3n-viola-linealidad}
  Un amplificador ideal puede cumplir \(y=Gx\) dentro de su rango. Al
  alcanzar los rieles, la salida se recorta.

  \begin{tcolorbox}[enhanced jigsaw, arc=.35mm, bottomrule=.15mm, bottomtitle=1mm, breakable, colback=white, colbacktitle=quarto-callout-important-color!10!white, colframe=quarto-callout-important-color-frame, coltitle=black, left=2mm, leftrule=.75mm, opacityback=0, opacitybacktitle=0.6, rightrule=.15mm, title=\textcolor{quarto-callout-important-color}{\faExclamation}\hspace{0.5em}{Importante}, titlerule=0mm, toprule=.15mm, toptitle=1mm]

    La linealidad depende del rango operativo. Un modelo local puede ser
    útil aunque el dispositivo completo sea no lineal.

  \end{tcolorbox}
\end{frame}

\begin{frame}{Invariancia temporal conserva la regla}
  \protect\phantomsection\label{invariancia-temporal-conserva-la-regla}
  Si

  \[y(t)=\mathcal{T}\{x(t)\},\]

  un sistema invariante cumple

  \[\mathcal{T}\{x(t-t_0)\}=y(t-t_0).\]

  \begin{tcolorbox}[enhanced jigsaw, arc=.35mm, bottomrule=.15mm, bottomtitle=1mm, breakable, colback=white, colbacktitle=quarto-callout-note-color!10!white, colframe=quarto-callout-note-color-frame, coltitle=black, left=2mm, leftrule=.75mm, opacityback=0, opacitybacktitle=0.6, rightrule=.15mm, title=\textcolor{quarto-callout-note-color}{\faInfo}\hspace{0.5em}{Nota}, titlerule=0mm, toprule=.15mm, toptitle=1mm]

    El desplazamiento de la entrada debe producir el mismo desplazamiento de
    la salida, sin cambiar forma ni ganancia.

  \end{tcolorbox}
\end{frame}

\begin{frame}{Adaptación fisiológica rompe invariancia}
  \protect\phantomsection\label{adaptaciuxf3n-fisioluxf3gica-rompe-invariancia}
  La respuesta a un mismo estímulo puede cambiar por fatiga, aprendizaje,
  medicación o estado autonómico.

  \begin{tcolorbox}[enhanced jigsaw, arc=.35mm, bottomrule=.15mm, bottomtitle=1mm, breakable, colback=white, colbacktitle=quarto-callout-warning-color!10!white, colframe=quarto-callout-warning-color-frame, coltitle=black, left=2mm, leftrule=.75mm, opacityback=0, opacitybacktitle=0.6, rightrule=.15mm, title=\textcolor{quarto-callout-warning-color}{\faExclamationTriangle}\hspace{0.5em}{Advertencia}, titlerule=0mm, toprule=.15mm, toptitle=1mm]

    Aplicar herramientas LTI a un sistema biológico exige justificar
    intervalo, amplitud y estado donde la aproximación resulta válida.

  \end{tcolorbox}
\end{frame}

\begin{frame}{Memoria y causalidad son propiedades distintas}
  \protect\phantomsection\label{memoria-y-causalidad-son-propiedades-distintas}
  \begin{itemize}
    \tightlist
    \item
          sin memoria: \(y(t)\) depende solo de \(x(t)\);
    \item
          con memoria: depende de otros instantes;
    \item
          causal: no usa valores futuros de la entrada.
  \end{itemize}

  \begin{tcolorbox}[enhanced jigsaw, arc=.35mm, bottomrule=.15mm, bottomtitle=1mm, breakable, colback=white, colbacktitle=quarto-callout-important-color!10!white, colframe=quarto-callout-important-color-frame, coltitle=black, left=2mm, leftrule=.75mm, opacityback=0, opacitybacktitle=0.6, rightrule=.15mm, title=\textcolor{quarto-callout-important-color}{\faExclamation}\hspace{0.5em}{Importante}, titlerule=0mm, toprule=.15mm, toptitle=1mm]

    Un promedio móvil causal tiene memoria. Un filtrado de fase cero fuera
    de línea usa muestras futuras y no es causal.

  \end{tcolorbox}
\end{frame}

\begin{frame}{Estabilidad BIBO limita la salida}
  \protect\phantomsection\label{estabilidad-bibo-limita-la-salida}
  Si \(|x(t)|\leq M_x<\infty\), estabilidad BIBO exige
  \(|y(t)|\leq M_y<\infty\).

  Para un sistema LTI continuo, una condición suficiente y necesaria
  habitual es

  \[\int_{-\infty}^{\infty}|h(t)|dt<\infty.\]

  \begin{tcolorbox}[enhanced jigsaw, arc=.35mm, bottomrule=.15mm, bottomtitle=1mm, breakable, colback=white, colbacktitle=quarto-callout-warning-color!10!white, colframe=quarto-callout-warning-color-frame, coltitle=black, left=2mm, leftrule=.75mm, opacityback=0, opacitybacktitle=0.6, rightrule=.15mm, title=\textcolor{quarto-callout-warning-color}{\faExclamationTriangle}\hspace{0.5em}{Advertencia}, titlerule=0mm, toprule=.15mm, toptitle=1mm]

    Estabilidad matemática no implica seguridad clínica ni validez
    fisiológica.

  \end{tcolorbox}
\end{frame}

\begin{frame}{El impulso ideal concentra área unitaria}
  \protect\phantomsection\label{el-impulso-ideal-concentra-uxe1rea-unitaria}
  \[\delta(t)=0\quad\text{para }t\neq0,\]

  \[\int_{-\infty}^{\infty}\delta(t)dt=1.\]

  Propiedad de selección:

  \[\int x(\tau)\delta(t-\tau)d\tau=x(t).\]

  \begin{tcolorbox}[enhanced jigsaw, arc=.35mm, bottomrule=.15mm, bottomtitle=1mm, breakable, colback=white, colbacktitle=quarto-callout-note-color!10!white, colframe=quarto-callout-note-color-frame, coltitle=black, left=2mm, leftrule=.75mm, opacityback=0, opacitybacktitle=0.6, rightrule=.15mm, title=\textcolor{quarto-callout-note-color}{\faInfo}\hspace{0.5em}{Nota}, titlerule=0mm, toprule=.15mm, toptitle=1mm]

    El impulso de Dirac es una distribución matemática. Los experimentos
    físicos usan pulsos breves con amplitud y energía limitadas.

  \end{tcolorbox}
\end{frame}

\begin{frame}{En discreto, el impulso es una secuencia ordinaria}
  \protect\phantomsection\label{en-discreto-el-impulso-es-una-secuencia-ordinaria}
  \[\delta[n]=\begin{cases}1,&n=0,\\0,&n\neq0.\end{cases}\]

  Toda secuencia puede escribirse como

  \[x[n]=\sum_{k=-\infty}^{\infty}x[k]\delta[n-k].\]

  \begin{tcolorbox}[enhanced jigsaw, arc=.35mm, bottomrule=.15mm, bottomtitle=1mm, breakable, colback=white, colbacktitle=quarto-callout-important-color!10!white, colframe=quarto-callout-important-color-frame, coltitle=black, left=2mm, leftrule=.75mm, opacityback=0, opacitybacktitle=0.6, rightrule=.15mm, title=\textcolor{quarto-callout-important-color}{\faExclamation}\hspace{0.5em}{Importante}, titlerule=0mm, toprule=.15mm, toptitle=1mm]

    La descomposición expresa la entrada como suma de impulsos desplazados y
    ponderados.

  \end{tcolorbox}
\end{frame}

\begin{frame}{La respuesta al impulso caracteriza un sistema LTI}
  \protect\phantomsection\label{la-respuesta-al-impulso-caracteriza-un-sistema-lti}
  \[h(t)=\mathcal{T}\{\delta(t)\}.\]

  Por linealidad e invariancia:

  \[y(t)=\int_{-\infty}^{\infty}x(\tau)h(t-\tau)d\tau.\]

  \begin{tcolorbox}[enhanced jigsaw, arc=.35mm, bottomrule=.15mm, bottomtitle=1mm, breakable, colback=white, colbacktitle=quarto-callout-important-color!10!white, colframe=quarto-callout-important-color-frame, coltitle=black, left=2mm, leftrule=.75mm, opacityback=0, opacitybacktitle=0.6, rightrule=.15mm, title=\textcolor{quarto-callout-important-color}{\faExclamation}\hspace{0.5em}{Importante}, titlerule=0mm, toprule=.15mm, toptitle=1mm]

    Conocer \(h\) permite calcular la respuesta a cualquier entrada dentro
    del modelo LTI.

  \end{tcolorbox}
\end{frame}

\begin{frame}{La salida es una convolución}
  \protect\phantomsection\label{la-salida-es-una-convoluciuxf3n}
  \[y(t)=x(t)*h(t).\]

  En tiempo discreto:

  \[y[n]=\sum_{k=-\infty}^{\infty}x[k]h[n-k].\]

  \begin{tcolorbox}[enhanced jigsaw, arc=.35mm, bottomrule=.15mm, bottomtitle=1mm, breakable, colback=white, colbacktitle=quarto-callout-note-color!10!white, colframe=quarto-callout-note-color-frame, coltitle=black, left=2mm, leftrule=.75mm, opacityback=0, opacitybacktitle=0.6, rightrule=.15mm, title=\textcolor{quarto-callout-note-color}{\faInfo}\hspace{0.5em}{Nota}, titlerule=0mm, toprule=.15mm, toptitle=1mm]

    La semana siguiente desarrolla interpretación gráfica, soporte y cálculo
    de la convolución.

  \end{tcolorbox}
\end{frame}

\begin{frame}{Causalidad de un LTI se lee en \(h\)}
  \protect\phantomsection\label{causalidad-de-un-lti-se-lee-en-h}
  Un sistema LTI causal cumple

  \[h(t)=0\quad\text{para }t<0,\]

  o en discreto \(h[n]=0\) para \(n<0\).

  \begin{tcolorbox}[enhanced jigsaw, arc=.35mm, bottomrule=.15mm, bottomtitle=1mm, breakable, colback=white, colbacktitle=quarto-callout-warning-color!10!white, colframe=quarto-callout-warning-color-frame, coltitle=black, left=2mm, leftrule=.75mm, opacityback=0, opacitybacktitle=0.6, rightrule=.15mm, title=\textcolor{quarto-callout-warning-color}{\faExclamationTriangle}\hspace{0.5em}{Advertencia}, titlerule=0mm, toprule=.15mm, toptitle=1mm]

    Una respuesta al impulso simétrica alrededor de cero no puede
    implementarse causalmente sin añadir retardo.

  \end{tcolorbox}
\end{frame}

\begin{frame}{Las exponenciales complejas son autofunciones LTI}
  \protect\phantomsection\label{las-exponenciales-complejas-son-autofunciones-lti}
  Si \(x(t)=e^{j2\pi ft}\), entonces

  \[y(t)=H(f)e^{j2\pi ft}.\]

  \begin{tcolorbox}[enhanced jigsaw, arc=.35mm, bottomrule=.15mm, bottomtitle=1mm, breakable, colback=white, colbacktitle=quarto-callout-important-color!10!white, colframe=quarto-callout-important-color-frame, coltitle=black, left=2mm, leftrule=.75mm, opacityback=0, opacitybacktitle=0.6, rightrule=.15mm, title=\textcolor{quarto-callout-important-color}{\faExclamation}\hspace{0.5em}{Importante}, titlerule=0mm, toprule=.15mm, toptitle=1mm]

    El sistema conserva la frecuencia y modifica amplitud y fase mediante
    \(H(f)\).

  \end{tcolorbox}
\end{frame}

\begin{frame}{\(H(f)\) es la transformada de \(h(t)\)}
  \protect\phantomsection\label{hf-es-la-transformada-de-ht}
  \[H(f)=\int_{-\infty}^{\infty}h(t)e^{-j2\pi ft}dt.\]

  La relación

  \[Y(f)=X(f)H(f)\]

  coincide con el teorema de convolución estudiado en la semana 6.

  \begin{tcolorbox}[enhanced jigsaw, arc=.35mm, bottomrule=.15mm, bottomtitle=1mm, breakable, colback=white, colbacktitle=quarto-callout-tip-color!10!white, colframe=quarto-callout-tip-color-frame, coltitle=black, left=2mm, leftrule=.75mm, opacityback=0, opacitybacktitle=0.6, rightrule=.15mm, title=\textcolor{quarto-callout-tip-color}{\faLightbulb}\hspace{0.5em}{Tip}, titlerule=0mm, toprule=.15mm, toptitle=1mm]

    Tiempo y frecuencia describen el mismo sistema mediante \(h(t)\) y
    \(H(f)\).

  \end{tcolorbox}
\end{frame}

\begin{frame}{Ejemplos de sistemas de procesamiento}
  \protect\phantomsection\label{ejemplos-de-sistemas-de-procesamiento}
  {\def\LTcaptype{none} % do not increment counter
    \begin{longtable}[]{@{}lll@{}}
      \toprule\noalign{}
      Sistema        & Respuesta temporal esperada & Acción frecuencial          \\
      \midrule\noalign{}
      \endhead
      promedio móvil & suaviza cambios rápidos     & atenúa frecuencias altas    \\
      diferenciador  & resalta pendientes          & amplifica frecuencias altas \\
      retardo        & desplaza la señal           & añade fase lineal           \\
      ganancia       & escala amplitud             & multiplica por constante    \\
      \bottomrule\noalign{}
    \end{longtable}
  }

  \begin{tcolorbox}[enhanced jigsaw, arc=.35mm, bottomrule=.15mm, bottomtitle=1mm, breakable, colback=white, colbacktitle=quarto-callout-warning-color!10!white, colframe=quarto-callout-warning-color-frame, coltitle=black, left=2mm, leftrule=.75mm, opacityback=0, opacitybacktitle=0.6, rightrule=.15mm, title=\textcolor{quarto-callout-warning-color}{\faExclamationTriangle}\hspace{0.5em}{Advertencia}, titlerule=0mm, toprule=.15mm, toptitle=1mm]

    La acción útil depende de preservar los eventos y medidas relevantes
    para la tarea biomédica.

  \end{tcolorbox}
\end{frame}

\begin{frame}[fragile]{Código: probar linealidad numéricamente}
  \protect\phantomsection\label{cuxf3digo-probar-linealidad-numuxe9ricamente}
  \begin{Shaded}
    \begin{Highlighting}[]
      \ImportTok{import}\NormalTok{ numpy }\ImportTok{as}\NormalTok{ np}

      \NormalTok{y\_combo }\OperatorTok{=}\NormalTok{ T(a}\OperatorTok{*}\NormalTok{x1 }\OperatorTok{+}\NormalTok{ b}\OperatorTok{*}\NormalTok{x2)}
      \NormalTok{y\_super }\OperatorTok{=}\NormalTok{ a}\OperatorTok{*}\NormalTok{T(x1) }\OperatorTok{+}\NormalTok{ b}\OperatorTok{*}\NormalTok{T(x2)}
      \NormalTok{error }\OperatorTok{=}\NormalTok{ np.linalg.norm(y\_combo }\OperatorTok{{-}}\NormalTok{ y\_super)}
      \NormalTok{referencia }\OperatorTok{=}\NormalTok{ np.linalg.norm(y\_combo)}
      \NormalTok{error\_rel }\OperatorTok{=}\NormalTok{ error }\OperatorTok{/} \BuiltInTok{max}\NormalTok{(referencia, }\FloatTok{1e{-}12}\NormalTok{)}
    \end{Highlighting}
  \end{Shaded}

  \begin{tcolorbox}[enhanced jigsaw, arc=.35mm, bottomrule=.15mm, bottomtitle=1mm, breakable, colback=white, colbacktitle=quarto-callout-warning-color!10!white, colframe=quarto-callout-warning-color-frame, coltitle=black, left=2mm, leftrule=.75mm, opacityback=0, opacitybacktitle=0.6, rightrule=.15mm, title=\textcolor{quarto-callout-warning-color}{\faExclamationTriangle}\hspace{0.5em}{Advertencia}, titlerule=0mm, toprule=.15mm, toptitle=1mm]

    Un error pequeño para algunos vectores apoya la hipótesis, pero no
    demuestra linealidad para todas las entradas.

  \end{tcolorbox}
\end{frame}

\begin{frame}{Verificación de la sesión 20}
  \protect\phantomsection\label{verificaciuxf3n-de-la-sesiuxf3n-20}
  Para \(y[n]=0.5x[n]+0.5x[n-1]\):

  \begin{enumerate}
    \tightlist
    \item
          verifique linealidad;
    \item
          determine memoria y causalidad;
    \item
          obtenga \(h[n]\);
    \item
          evalúe estabilidad BIBO;
    \item
          interprete su efecto sobre un ECG;
    \item
          anticipe su respuesta frecuencial.
  \end{enumerate}

  \begin{tcolorbox}[enhanced jigsaw, arc=.35mm, bottomrule=.15mm, bottomtitle=1mm, breakable, colback=white, colbacktitle=quarto-callout-tip-color!10!white, colframe=quarto-callout-tip-color-frame, coltitle=black, left=2mm, leftrule=.75mm, opacityback=0, opacitybacktitle=0.6, rightrule=.15mm, title=\textcolor{quarto-callout-tip-color}{\faLightbulb}\hspace{0.5em}{Criterio de logro}, titlerule=0mm, toprule=.15mm, toptitle=1mm]

    La respuesta debe conectar la ecuación, las propiedades, \(h[n]\) y la
    interpretación biomédica.

  \end{tcolorbox}
\end{frame}

\begin{frame}{Cierre de la semana 10}
  \protect\phantomsection\label{cierre-de-la-semana-10}
  \begin{itemize}
    \tightlist
    \item
          La selección de ventana define una escala tiempo-frecuencia.
    \item
          Ninguna configuración maximiza ambas resoluciones.
    \item
          Linealidad e invariancia temporal permiten usar superposición.
    \item
          La respuesta al impulso caracteriza un sistema LTI.
    \item
          \(h(t)\) y \(H(f)\) son descripciones complementarias.
  \end{itemize}

  \begin{tcolorbox}[enhanced jigsaw, arc=.35mm, bottomrule=.15mm, bottomtitle=1mm, breakable, colback=white, colbacktitle=quarto-callout-note-color!10!white, colframe=quarto-callout-note-color-frame, coltitle=black, left=2mm, leftrule=.75mm, opacityback=0, opacitybacktitle=0.6, rightrule=.15mm, title=\textcolor{quarto-callout-note-color}{\faInfo}\hspace{0.5em}{Pregunta de salida}, titlerule=0mm, toprule=.15mm, toptitle=1mm]

    ¿Qué evidencia experimental necesitaría para tratar un sistema
    fisiológico como LTI durante un intervalo específico?

  \end{tcolorbox}
\end{frame}

\section{Síntesis y estudio
  autónomo}\label{suxedntesis-y-estudio-autuxf3nomo}

\begin{frame}{La cadena completa de decisiones}
  \protect\phantomsection\label{la-cadena-completa-de-decisiones}
  \[\begin{aligned}
      \text{pregunta} & \rightarrow\text{adquisición}\rightarrow\text{muestreo}\rightarrow\text{segmento}                            \\
                      & \rightarrow\text{ventana}\rightarrow\text{estimador}\rightarrow\text{validación}\rightarrow\text{inferencia}
    \end{aligned}\]

  \begin{tcolorbox}[enhanced jigsaw, arc=.35mm, bottomrule=.15mm, bottomtitle=1mm, breakable, colback=white, colbacktitle=quarto-callout-important-color!10!white, colframe=quarto-callout-important-color-frame, coltitle=black, left=2mm, leftrule=.75mm, opacityback=0, opacitybacktitle=0.6, rightrule=.15mm, title=\textcolor{quarto-callout-important-color}{\faExclamation}\hspace{0.5em}{Importante}, titlerule=0mm, toprule=.15mm, toptitle=1mm]

    Cada resultado espectral debe conservar trazabilidad hacia el registro,
    los parámetros y el estado fisiológico observado.

  \end{tcolorbox}
\end{frame}

\begin{frame}{Errores que deben poder diagnosticarse}
  \protect\phantomsection\label{errores-que-deben-poder-diagnosticarse}
  \begin{enumerate}
    \tightlist
    \item
          sumar magnitudes cuando corresponde sumar espectros complejos;
    \item
          interpretar zero padding como mayor resolución real;
    \item
          atribuir fuga a nuevas oscilaciones fisiológicas;
    \item
          usar Nyquist sin margen para el filtro analógico;
    \item
          confundir potencia por bin con PSD;
    \item
          promediar estados incompatibles con Welch;
    \item
          leer resolución desde el número de columnas del espectrograma;
    \item
          asumir que un sistema biológico es LTI sin verificar rango y estado.
  \end{enumerate}
\end{frame}

\begin{frame}{Autoevaluación acumulativa}
  \protect\phantomsection\label{autoevaluaciuxf3n-acumulativa}
  Puede explicar, sin consultar notas:

  \begin{itemize}
    \tightlist
    \item
          cómo un retraso modifica la fase;
    \item
          por qué \(\Delta f=1/T\);
    \item
          qué cambia al aplicar Hann;
    \item
          cómo se produce aliasing;
    \item
          qué unidades tiene una PSD;
    \item
          qué intercambia Welch entre resolución y varianza;
    \item
          cómo se selecciona una ventana STFT;
    \item
          por qué \(h(t)\) caracteriza un sistema LTI.
  \end{itemize}

  \begin{tcolorbox}[enhanced jigsaw, arc=.35mm, bottomrule=.15mm, bottomtitle=1mm, breakable, colback=white, colbacktitle=quarto-callout-note-color!10!white, colframe=quarto-callout-note-color-frame, coltitle=black, left=2mm, leftrule=.75mm, opacityback=0, opacitybacktitle=0.6, rightrule=.15mm, title=\textcolor{quarto-callout-note-color}{\faInfo}\hspace{0.5em}{Nota}, titlerule=0mm, toprule=.15mm, toptitle=1mm]

    Una explicación completa incluye ecuación, unidades, supuesto, efecto
    observable y límite de interpretación.

  \end{tcolorbox}
\end{frame}

\begin{frame}{Actividad integradora sugerida}
  \protect\phantomsection\label{actividad-integradora-sugerida}
  Con un registro biomédico real o simulado:

  \begin{enumerate}
    \tightlist
    \item
          audite muestreo, unidades y duración;
    \item
          compare periodograma y Welch;
    \item
          analice dos ventanas;
    \item
          verifique potencia mediante integración;
    \item
          construya una STFT con dos longitudes;
    \item
          interprete cambios y artefactos;
    \item
          documente parámetros en Python;
    \item
          formule una conclusión limitada por la evidencia.
  \end{enumerate}

  \begin{tcolorbox}[enhanced jigsaw, arc=.35mm, bottomrule=.15mm, bottomtitle=1mm, breakable, colback=white, colbacktitle=quarto-callout-tip-color!10!white, colframe=quarto-callout-tip-color-frame, coltitle=black, left=2mm, leftrule=.75mm, opacityback=0, opacitybacktitle=0.6, rightrule=.15mm, title=\textcolor{quarto-callout-tip-color}{\faLightbulb}\hspace{0.5em}{Producto esperado}, titlerule=0mm, toprule=.15mm, toptitle=1mm]

    Un cuaderno reproducible que permita reconstruir cada figura y
    justificar cada decisión.

  \end{tcolorbox}
\end{frame}

\begin{frame}{Referencias y alcance}
  \protect\phantomsection\label{referencias-y-alcance}
  Contenido alineado con el programa de \textbf{Sistemas y Señales
    Biomédicos}, semanas 6--10, y con el desarrollo de Fourier, adquisición,
  potencia y sistemas lineales del texto guía
  (\citeproc{ref-semmlow2018}{Semmlow 2018}). Para formulación clásica
  consultar (\citeproc{ref-oppenheim1997}{Oppenheim et~al. 1997}). Los
  ejemplos computacionales siguen prácticas de Python descritas en
  (\citeproc{ref-esakkirajan2024}{Esakkirajan et~al. 2024}). Para
  aplicaciones de bioseñales consultar
  (\citeproc{ref-rangayyan2015}{Rangayyan 2015}).

  \begin{tcolorbox}[enhanced jigsaw, arc=.35mm, bottomrule=.15mm, bottomtitle=1mm, breakable, colback=white, colbacktitle=quarto-callout-warning-color!10!white, colframe=quarto-callout-warning-color-frame, coltitle=black, left=2mm, leftrule=.75mm, opacityback=0, opacitybacktitle=0.6, rightrule=.15mm, title=\textcolor{quarto-callout-warning-color}{\faExclamationTriangle}\hspace{0.5em}{Advertencia}, titlerule=0mm, toprule=.15mm, toptitle=1mm]

    Estas diapositivas apoyan formación en ingeniería. Los ejemplos no
    establecen criterios diagnósticos ni recomendaciones clínicas.

  \end{tcolorbox}
\end{frame}

\begin{frame}{Cierre: criterio de dominio}
  \protect\phantomsection\label{cierre-criterio-de-dominio}
  El estudiante debe sostener esta afirmación con evidencia:

  \begin{quote}
    Un espectro no es una propiedad aislada del organismo. Es el resultado
    de observar, segmentar y estimar una señal bajo decisiones explícitas.
  \end{quote}

  \begin{tcolorbox}[enhanced jigsaw, arc=.35mm, bottomrule=.15mm, bottomtitle=1mm, breakable, colback=white, colbacktitle=quarto-callout-important-color!10!white, colframe=quarto-callout-important-color-frame, coltitle=black, left=2mm, leftrule=.75mm, opacityback=0, opacitybacktitle=0.6, rightrule=.15mm, title=\textcolor{quarto-callout-important-color}{\faExclamation}\hspace{0.5em}{Importante}, titlerule=0mm, toprule=.15mm, toptitle=1mm]

    El dominio del tema se demuestra cuando una conclusión puede rastrearse
    hasta la adquisición, las unidades, los parámetros y los supuestos.

  \end{tcolorbox}
\end{frame}

\begin{frame}[allowframebreaks]{Referencias}
  \protect\phantomsection\label{referencias}
  \protect\phantomsection\label{refs}
  \begin{CSLReferences}{1}{1}
    \bibitem[\citeproctext]{ref-esakkirajan2024}
    Esakkirajan, S., T. Veerakumar, y Badri N. Subudhi. 2024. \emph{Digital
      Signal Processing: Illustration Using Python}. Springer.
    \url{https://doi.org/10.1007/978-981-99-6752-0}.

    \bibitem[\citeproctext]{ref-oppenheim1997}
    Oppenheim, Alan V., Alan S. Willsky, y S. Hamid Nawab. 1997.
    \emph{Signals and Systems}. 2.ª ed. Prentice Hall.

    \bibitem[\citeproctext]{ref-rangayyan2015}
    Rangayyan, Rangaraj M. 2015. \emph{Biomedical Signal Analysis: A
      Case-Study Approach}. 2.ª ed. Wiley-IEEE Press.
    \url{https://doi.org/10.1002/9781119068129}.

    \bibitem[\citeproctext]{ref-semmlow2018}
    Semmlow, John L. 2018. \emph{Circuits, Signals, and Systems for
      Bioengineers: A MATLAB-Based Introduction}. 3.ª ed. Academic Press.

  \end{CSLReferences}
\end{frame}




\end{document}
